% Generated by docs/book/figscripts/make_app_E.py from the chapters named in docs/book/main.tex. Do not edit by hand.
% Every displayed equation, in book order; bodies copied from the chapters with \label, \nonumber, \notag and \tag removed.
\chapter{Formula sheet}\label{app:formulas}
\newcounter{iamfs}\renewcommand{\theiamfs}{\thechapter.\arabic{iamfs}}
\newcommand{\iamfsentry}[5]{\refstepcounter{iamfs}\label{app:fs:#1}%
  \begin{equation*}#2\tag{\theiamfs}\end{equation*}\nopagebreak%
  \noindent\parbox{\linewidth}{\small\setlength{\baselineskip}{1.15em}\setlength{\parskip}{2pt}#3\par\emph{Status:} #4.\quad\emph{In the text:} #5.}\par\medskip}
This appendix collects every displayed equation of the seven parts (559 displays), in the order in which they appear, grouped by part and
chapter. Each entry gives the equation as printed in the chapter, its status, and where it appears in the text; the entries that were in the
earlier selection of this sheet keep their line on what the equation is. The numbers in parentheses belong to this appendix; the chapter's
own equation number is given where the chapter numbers the display. The status is the label the chapter attaches to the display. Where the
chapter attaches none, the status written for the earlier selection is used, then the first label later in the same paragraph; an entry
with none of these says so. A relation that appears in more than one place is listed each time, with pointers to its other entries.
Inline formulas are not collected; symbols and units are in Appendix~\ref{app:notation} and terms in Appendix~\ref{app:glossary}.
Chapters with no displayed equation, and so no entry here: \ref{ch:bhinformation}, \ref{ch:saturation}, \ref{ch:higgsrecord}, \ref{ch:translation}, \ref{ch:instrument}, \ref{ch:identity}, \ref{ch:serial}, \ref{ch:discipline}, \ref{ch:chain}, \ref{ch:report}, \ref{ch:leukocyte}, \ref{part4:ch:reach}, \ref{ch:status}, \ref{ch:synthesis}, \ref{ch:reach}, \ref{ch:open}, \ref{ch:statusall}, \ref{ch:conclusion}; the front matter has none either.

\section{\texorpdfstring{Part~\ref{part:1}}{Part 1}: Introduction to IAM's Law and the Virial Theorem}
\subsection*{Chapter~\ref{ch:iams_law}: IAM's Law: the thermodynamic cost of a classical record}
\iamfsentry{0}{\Delta E_{\rm rec}=k_BT_H\ln2\quad\text{per bit}.}{IAM's Law: the minimum cost of one classical bit, paid at the temperature of the nearest encoding surface (horizon temperature $T_H$).}{\conjecture{} (Assumption~2: the horizon is the reservoir); in its laboratory limit it is Landauer's bound, \observed}{Eq.~\eqref{eq:law}; Chapter~\ref{ch:iams_law}, section ``Statement of the law''}
\iamfsentry{1}{|\psi_S\rangle\otimes|E_0\rangle\;\longrightarrow\;\sum_ic_i\,|s_i\rangle\otimes|E_i\rangle,}{Decoherence: a system--environment product state evolves into an entangled sum over pointer states.}{\derived{} (standard quantum mechanics)}{Eq.~\eqref{eq:law_entangle}; Chapter~\ref{ch:iams_law}, section ``The quantum boundary''; see also (\ref{app:fs:76})}
\iamfsentry{2}{\rho_S={\rm Tr}_E|\Psi\rangle\langle\Psi|=\sum_{ij}c_ic_j^*\langle E_j|E_i\rangle\,|s_i\rangle\langle s_j|\;\longrightarrow\;\sum_i|c_i|^2|s_i\rangle\langle s_i|.}{}{\derived{}}{Eq.~\eqref{eq:law_rhodiag}; Chapter~\ref{ch:iams_law}, section ``The quantum boundary''}
\iamfsentry{3}{I=\log_2N\ \text{bits}.}{}{\observed{}}{Eq.~\eqref{eq:law_bits}; Chapter~\ref{ch:iams_law}, section ``The quantum boundary''; see also (\ref{app:fs:79})}
\iamfsentry{4}{\Delta E_{\rm Landauer}=k_BT\ln2}{}{\observed{}}{Eq.~\eqref{eq:law_landauer}; Chapter~\ref{ch:iams_law}, section ``The quantum boundary''; see also (\ref{app:fs:80})}
\iamfsentry{5}{\Delta E_{\rm rec}(a)=k_BT_H(a)\ln2=\frac{\hbar H(a)}{2\pi}\ln2 .}{Cost per bit at the cosmic horizon at scale factor $a$, using the Gibbons--Hawking temperature.}{\derived{} (given Eq.~\eqref{eq:law}, a \conjecture, with $T=T_{GH}$)}{Eq.~\eqref{eq:epochcost}; Chapter~\ref{ch:iams_law}, section ``The quantum boundary''}
\iamfsentry{6}{E_{\rm Landauer}=Q=T\Delta S_{\min}=-E_f=\langle K\rangle=\tfrac12|\langle V\rangle| .}{}{\derived{} (the equalities)}{Eq.~\eqref{eq:law_virialchain}; Chapter~\ref{ch:iams_law}, section ``The local expression: the kinetic half''}
\iamfsentry{7}{\delta Q=\int_{\mathcal H}T_{ab}\chi^a\,d\Sigma^b=-\kappa\int_{\mathcal H}\lambda\,T_{ab}k^ak^b\,d\lambda\,dA.}{}{\derived{}}{Eq.~\eqref{eq:law_dQ}; Chapter~\ref{ch:iams_law}, section ``The global expression: completing the horizon entropy''; see also (\ref{app:fs:82})}
\iamfsentry{8}{\delta S=-\eta\int_{\mathcal H}\lambda\,R_{ab}k^ak^b\,d\lambda\,dA.}{}{\derived{}}{Eq.~\eqref{eq:law_dS}; Chapter~\ref{ch:iams_law}, section ``The global expression: completing the horizon entropy''}
\iamfsentry{9}{R_{ab}-\tfrac12Rg_{ab}+\Lambda g_{ab}=\frac{2\pi}{\hbar c\,\eta}T_{ab}\equiv\frac{8\pi G}{c^4}T_{ab},\qquad G=\frac{c^3}{4\hbar\eta}.}{}{\derived{}}{Eq.~\eqref{eq:law_jacobson}; Chapter~\ref{ch:iams_law}, section ``The global expression: completing the horizon entropy''}
\iamfsentry{10}{\frac{\hbar\eta}{2\pi}=\frac{c^3}{8\pi G}.}{}{\derived{}}{Eq.~\eqref{eq:law_structure}; Chapter~\ref{ch:iams_law}, section ``The global expression: completing the horizon entropy''}
\iamfsentry{11}{\eta=\frac{c^3}{4\hbar G}=\frac{1}{4\ell_P^2},\qquad \frac14=\frac{2\pi}{8\pi}.}{The Bekenstein--Hawking coefficient as the ratio of the Euclidean period $2\pi$ to the $8\pi$ of the field equations.}{\calc{} (as labeled in the chapter)}{Eq.~\eqref{eq:law_eta}; Chapter~\ref{ch:iams_law}, section ``The global expression: completing the horizon entropy''}
\iamfsentry{12}{\delta A_{\min}=\frac{8\pi G}{\kappa c^2}\cdot\frac{\hbar\kappa}{2\pi c}=\frac{4\hbar G}{c^3}=4\ell_P^2,}{}{\derived{}}{Eq.~\eqref{eq:law_dAmin}; Chapter~\ref{ch:iams_law}, section ``The global expression: completing the horizon entropy''}
\iamfsentry{13}{S_{\rm total}=S_{\rm geo}+S_{\rm info},}{}{\conjecture{}}{Eq.~\eqref{eq:law_Stotal}; Chapter~\ref{ch:iams_law}, section ``The global expression: completing the horizon entropy''; see also (\ref{app:fs:103}), (\ref{app:fs:404})}
\iamfsentry{14}{\dot H=-4\pi G(\rho+P).}{Cai--Kim: the first law at the apparent horizon gives the second Friedmann equation.}{\derived{} (Cai and Kim 2005)}{Eq.~\eqref{eq:law_F2}; Chapter~\ref{ch:iams_law}, section ``The global expression: completing the horizon entropy''; see also (\ref{app:fs:100}), (\ref{app:fs:169})}
\iamfsentry{15}{S_{\rm total}=S_{\rm geo}+S_{\rm info},\qquad -dE=T_H\,d\!\left(S_{\rm geo}+S_{\rm info}\right).}{Horizon entropy completed by the record term, and the first law applied to it at the apparent horizon. $S_{\rm total}$ is defined here.}{\conjecture{} (Assumption~2)}{Eq.~\eqref{eq:firstlaw}; Chapter~\ref{ch:iams_law}, section ``The global expression: completing the horizon entropy''}
\iamfsentry{16}{\frac{4\pi(\rho+P)}{H^2}=-\frac{\dot H}{GH^2}+\frac{H}{2\pi}\dot S_{\rm info}\quad\Longrightarrow\quad
\dot H=-4\pi G(\rho+P)+\frac{G}{2\pi}H^3\dot S_{\rm info}.}{}{\derived{}}{Eq.~\eqref{eq:law_F2info}; Chapter~\ref{ch:iams_law}, section ``The global expression: completing the horizon entropy''}
\iamfsentry{17}{\dot S_{\rm info}=\frac{8\pi^2\rho_{\rm info}}{3aH^3},\qquad \frac{dS_{\rm info}}{d\ln a}\Big|_{a=1}=\beta_m\frac{\pi}{GH_0^2}=\beta_m\,S_{\rm geo}(a{=}1),}{}{\calc{}}{Eq.~\eqref{eq:law_Sneed}; Chapter~\ref{ch:iams_law}, section ``The global expression: completing the horizon entropy''}
\iamfsentry{18}{H_m^2=H_{\Lambda{\rm CDM}}^2+\beta_mE(a)H_0^2=\frac{8\pi G}{3}\rho_m+\frac\Lambda3+\beta_mE(a)H_0^2 .}{Matter-sector expansion rate: the record term adds $\beta_mE(a)H_0^2$ to the rate that matter perturbations feel.}{\derived{} (given the completed first law, Eq.~\eqref{eq:firstlaw}, a \conjecture, and the sector split, a \conjecture)}{Eq.~\eqref{eq:Hm}; Chapter~\ref{ch:iams_law}, section ``The global expression: completing the horizon entropy''}
\iamfsentry{19}{\dot I\propto\rho_m(a)\,D(a)^n\,f(a)\,H(a),}{Assumed form of the rate of record production from structure formation.}{\conjecture{} (assumed form of the writing rate)}{Eq.~\eqref{eq:rate}; Chapter~\ref{ch:iams_law}, section ``The global expression: completing the horizon entropy''; see also (\ref{app:fs:104})}
\iamfsentry{20}{\frac{dS_{\rm info}}{dt}\propto\frac{\dot I}{T_HA_H},}{Rate of change of the record entropy: records priced at $k_BT_H\ln2$ and stored per unit horizon area.}{\conjecture{} (assumed accounting)}{Eq.~\eqref{eq:Sdot}; Chapter~\ref{ch:iams_law}, section ``The global expression: completing the horizon entropy''}
\iamfsentry{21}{\dot\rho_{\rm info}=\rho_{\rm info}\,\frac{H}{a}.}{}{\conjecture{} (the record constraint, a premise)}{Eq.~\eqref{eq:law_constraint}; Chapter~\ref{ch:iams_law}, section ``The global expression: completing the horizon entropy''}
\iamfsentry{22}{\ln\frac{\rho_{\rm info}(a)}{\rho_{\rm info}(1)}=\int_1^a\frac{da'}{a'^2}=1-\frac1a,\qquad
E(a)\equiv\frac{\rho_{\rm info}(a)}{\rho_{\rm info}(1)}=\exp\!\left(1-\frac1a\right)=e^{-z}.}{The activation function $E(a)$, normalized to $E(1)=1$.}{\derived{} (given the record-density constraint, a \conjecture)}{Eq.~\eqref{eq:Ea}; Chapter~\ref{ch:iams_law}, section ``The global expression: completing the horizon entropy''}
\iamfsentry{23}{\phi(t)\equiv\ln E(a)=1-\frac1a ,}{}{\derived{}}{Eq.~\eqref{eq:law_phi}; Chapter~\ref{ch:iams_law}, section ``The global expression: completing the horizon entropy''}
\iamfsentry{24}{S_{\rm info}=\int\left[-\frac{3H_0^2}{8\pi G}\beta_me^{\phi}+\lambda\left(\dot\phi-\frac{H}{a}\right)\right]\sqrt{-g}\,d^4x,}{Constrained action for the record term ($\phi=\ln E$, Lagrange multiplier $\lambda$); its metric variation returns Eq.~\eqref{eq:Hm}.}{definition (the action); that its variation returns Eq.~\eqref{eq:Hm} is \derived}{Eq.~\eqref{eq:law_action}; Chapter~\ref{ch:iams_law}, section ``The global expression: completing the horizon entropy''}
\iamfsentry{25}{L=-\frac{3}{8\pi G}\frac{aa'^2}{N}-Na^3\Big(\frac{\rho_{m0}}{a^3}+\rho_\Lambda+\rho_0\beta_me^{\phi}\Big)+\lambda a^3\Big(\phi'-\frac{a'}{a^2}\Big).}{}{\derived{}}{Eq.~\eqref{eq:law_L}; Chapter~\ref{ch:iams_law}, section ``The global expression: completing the horizon entropy''}
\iamfsentry{26}{w_{\rm info}(a)=-1-\frac{1}{3a}:\qquad -1.67\ (a=0.5),\quad -1.33\ (a=1),\quad -1.17\ (a=2).}{Effective equation of state of the record term, with values at $a=0.5,1,2$.}{\derived{} (given $\rho_{\rm info}\propto E(a)$ from the record constraint, a \conjecture); values \calc}{Eq.~\eqref{eq:law_winfo}; Chapter~\ref{ch:iams_law}, section ``The global expression: completing the horizon entropy''}
\iamfsentry{27}{\rho_{\rm info}(a{=}1)=\tfrac12\rho_m=\frac{\Omega_m}{2}\,\rho_{\rm crit}.}{}{\prediction{}}{Eq.~\eqref{eq:law_rhoinfo1}; Chapter~\ref{ch:iams_law}, section ``The global expression: completing the horizon entropy''}
\iamfsentry{28}{\beta_m=\frac{\Omega_m}{2}=0.15765\quad(\Omega_m=0.3153).}{}{\prediction{}}{Eq.~\eqref{eq:law_beta}; Chapter~\ref{ch:iams_law}, section ``The global expression: completing the horizon entropy''}
\iamfsentry{29}{\mu(a)=\frac{H^2_{\Lambda{\rm CDM}}(a)}{H^2_{\Lambda{\rm CDM}}(a)+\beta_mE(a)H_0^2}<1.}{Effective growth coupling of matter, $\mu(a)<1$; light keeps $\Sigma=1$.}{\derived{} (given Eq.~\eqref{eq:Hm} read as an effective Newton constant; that reading is a \conjecture)}{Eq.~\eqref{eq:mu}; Chapter~\ref{ch:iams_law}, section ``The global expression: completing the horizon entropy''}
\iamfsentry{30}{\Sigma(a)=1.}{}{\prediction{}}{Eq.~\eqref{eq:law_Sigma}; Chapter~\ref{ch:iams_law}, section ``The global expression: completing the horizon entropy''}
\iamfsentry{31}{\delta\phi=0}{The record field is unperturbed in linear theory.}{\derived{} (given the record constraint, a \conjecture, and an unperturbed background horizon at first order)}{Eq.~\eqref{eq:law_dphi}; Chapter~\ref{ch:iams_law}, section ``Perturbations: why growth is weaker and light is not''}
\iamfsentry{32}{\begin{aligned}&\text{Poisson equation:} && k^2\Psi=-4\pi Ga^2\rho_m\delta_m, \\
&\text{no anisotropic stress:} && \Psi=\Phi, \\
&\text{matter growth:} && \ddot\delta_m+2H_m\dot\delta_m-4\pi G\rho_m\delta_m=0,\end{aligned}}{}{\derived{}}{Eqs.~\eqref{eq:law_poisson}, \eqref{eq:law_noaniso}, \eqref{eq:law_growth}; Chapter~\ref{ch:iams_law}, section ``Perturbations: why growth is weaker and light is not''}
\iamfsentry{33}{\mu(a)=\frac{E^2_{\Lambda{\rm CDM}}(a)}{E^2_{{\rm eff},m}(a)}=\frac{H^2_{\Lambda{\rm CDM}}(a)}{H^2_{\Lambda{\rm CDM}}(a)+\beta_mE(a)H_0^2}<1,}{}{\interp{}}{Eq.~\eqref{eq:law_muratio}; Chapter~\ref{ch:iams_law}, section ``Perturbations: why growth is weaker and light is not''}
\iamfsentry{34}{\begin{aligned}\text{(i) }G_{\rm eff}=\mu G:&\quad D''+\Big(2+\frac{d\ln H}{d\ln a}\Big)D'-\tfrac32\Omega_m(a)\,\mu\,D=0,\\
\text{(ii) friction, }\Lambda\text{CDM clock:}&\quad D''+\Big(\frac{d\ln H}{d\ln a}+2\frac{H_m}{H}\Big)D'-\tfrac32\Omega_m(a)D=0,\\
\text{(iii) whole equation on }H_m:&\quad D''+\Big(2+\frac{d\ln H_m}{d\ln a}\Big)D'-\tfrac32\frac{\Omega_ma^{-3}H_0^2}{H_m^2}D=0 .\end{aligned}}{}{\calc{}}{Chapter~\ref{ch:iams_law}, section ``Perturbations: why growth is weaker and light is not''}
\iamfsentry{35}{\frac{dS_{\rm info}}{d\ln a}\propto\frac{\rho_m\,D^n f}{T_H\,A_H}\propto\frac{a^{-3}a^n}{a^{-3/2}a^3}=a^{\,n-9/2}.}{}{\derived{}}{Eq.~\eqref{eq:law_dSdlna}; Chapter~\ref{ch:iams_law}, section ``The exponent: from the horizon and from the halos''}
\iamfsentry{36}{S_{\rm info}(a)\propto\frac{a^{\,n-9/2}}{n-9/2}.}{}{\derived{}}{Eq.~\eqref{eq:law_Sa}; Chapter~\ref{ch:iams_law}, section ``The exponent: from the horizon and from the halos''}
\iamfsentry{37}{n-\frac92=-1\quad\Longrightarrow\quad n=\frac72 .}{}{\derived{} (given the form and the rate law)}{Eq.~\eqref{eq:law_n}; Chapter~\ref{ch:iams_law}, section ``The exponent: from the horizon and from the halos''; see also (\ref{app:fs:114})}
\iamfsentry{38}{\frac{S_{\rm info}(R)}{A(R)}\ge\frac{k_B}{4\ell_P^2}.}{}{\derived{}}{Eq.~\eqref{eq:law_saturation}; Chapter~\ref{ch:iams_law}, section ``Black holes: the local saturation limit''}
\iamfsentry{39}{\frac{4\pi Gk_BM^2/(\hbar c)}{4\pi\cdot4G^2M^2/c^4}=\frac{k_Bc^3}{4\hbar G}=\frac{k_B}{4\ell_P^2}.}{}{\derived{}}{Eq.~\eqref{eq:law_hoop}; Chapter~\ref{ch:iams_law}, section ``Black holes: the local saturation limit''}
\iamfsentry{40}{M_{\rm eq}=\frac{c^3}{4GH},}{}{\derived{}}{Eq.~\eqref{eq:law_Meq}; Chapter~\ref{ch:iams_law}, section ``Black holes: the local saturation limit''; see also (\ref{app:fs:164})}
\iamfsentry{41}{\Gamma=\frac{P}{k_BT_{\rm BH}\ln2}=\frac{c^3}{1920\,GM\ln2}\ \text{bits\,s}^{-1},\qquad P=\frac{\hbar c^6}{15360\pi G^2M^2},}{}{\calc{}}{Eq.~\eqref{eq:law_Gamma}; Chapter~\ref{ch:iams_law}, section ``Black holes: the local saturation limit''}
\iamfsentry{42}{H_0^{\rm sirens}=H_0^{\rm CMB}\sqrt{1+\beta_m}=67.16\times1.0759=72.26\ \text{km\,s}^{-1}\text{Mpc}^{-1}.}{}{\prediction{}}{Eq.~\eqref{eq:law_sirens}; Chapter~\ref{ch:iams_law}, section ``Observational consistency''}
\iamfsentry{43}{\frac{\Lambda_{\rm obs}}{\rho_{\rm vac}}=\frac2\pi\Big(\frac{\ell_P}{\ell_H}\Big)^2\frac{\Omega_b}{\Omega_m},}{}{\calc{}}{Eq.~\eqref{eq:law_ccbase}; Chapter~\ref{ch:iams_law}, section ``Observational consistency''}
\iamfsentry{44}{\frac{T_{dS}}{T_H}=\sqrt{\Omega_\Lambda}=\sqrt{0.6847}=0.8275 .}{}{\calc{}}{Eq.~\eqref{eq:law_TdS}; Chapter~\ref{ch:iams_law}, section ``Observational consistency''}
\iamfsentry{45}{\Omega_bh^2=0.02232\pm0.00014,\qquad \eta=(6.113\pm0.037)\times10^{-10}.}{}{\measured{}}{Eq.~\eqref{eq:law_eta_chain}; Chapter~\ref{ch:iams_law}, section ``Observational consistency''}
\iamfsentry{46}{\Delta E_{\rm rec}=k_BT_H\ln2\quad\text{per bit},}{}{\interp{}}{Eq.~\eqref{eq:law}, restated; Chapter~\ref{ch:iams_law}, section ``Discussion''}
\iamfsentry{47}{\mathcal M=\frac{E_{\rm drive}}{k_BT_{\rm local}},}{The Mahaffey number: drive energy over the thermal quantum.}{definition}{Eq.~\eqref{eq:law_mahaffey}; Chapter~\ref{ch:iams_law}, section ``The Mahaffey number''}
\subsection*{Chapter~\ref{ch:virial_law}: One half at every scale: the virial partition}
\iamfsentry{48}{2\langle K\rangle=k\,\langle V\rangle .}{Virial theorem for a potential homogeneous of degree $k$.}{\derived{} (Clausius; Euler's theorem)}{Eq.~\eqref{eq:vl_euler}; Chapter~\ref{ch:virial_law}, section ``The theorem''; see also (\ref{app:fs:429})}
\iamfsentry{49}{2\langle K\rangle+\langle V\rangle=0,\qquad \langle K\rangle=\tfrac12|\langle V\rangle|,\qquad E=\langle K\rangle+\langle V\rangle=-\langle K\rangle .}{}{\derived{}}{Eq.~\eqref{eq:vl_virial}; Chapter~\ref{ch:virial_law}, section ``The theorem''}
\iamfsentry{50}{\langle T\rangle=\frac n2\,\langle|V|\rangle .}{}{\derived{}}{Eq.~\eqref{eq:vl_n}; Chapter~\ref{ch:virial_law}, section ``The theorem''}
\iamfsentry{51}{\eta=-T/E=1.0000000000\qquad(\text{SD}=0\ \text{across the 20 elements and 10 molecules}).}{}{\calc{}}{Chapter~\ref{ch:virial_law}, section ``Atoms and molecules''}
\iamfsentry{52}{M_{\rm Ch}=\frac{\omega_3^0\sqrt{3\pi}}{2}\left(\frac{\hbar c}{G}\right)^{3/2}\frac{1}{(\mu_em_u)^2}=1.456\,M_\odot\qquad(\mu_e=2),}{}{\derived{}}{Eq.~\eqref{eq:vl_chandra}; Chapter~\ref{ch:virial_law}, section ``Stars''}
\iamfsentry{53}{S=\frac{k_Bc^3A}{4G\hbar}=\frac{4\pi Gk_BM^2}{\hbar c},\qquad T_H=\frac{\hbar c^3}{8\pi Gk_BM}.}{}{\derived{}}{Eq.~\eqref{eq:vl_bh}; Chapter~\ref{ch:virial_law}, section ``The black-hole horizon''}
\iamfsentry{54}{N\,k_BT_H\ln2=T_HS=\tfrac12Mc^2 .}{}{\derived{}}{Eq.~\eqref{eq:vl_smarr}; Chapter~\ref{ch:virial_law}, section ``The black-hole horizon''}
\iamfsentry{55}{\beta_m=\frac{\Omega_m}{2}=0.15765\qquad(\Omega_m=0.3153,\ \text{Planck 2018}).}{The cosmological coupling from the virial half, with Planck 2018 $\Omega_m=0.3153$.}{\prediction{}}{Eq.~\eqref{eq:vl_beta}; Chapter~\ref{ch:virial_law}, section ``The cosmic horizon''; see also (\ref{app:fs:62})}
\subsection*{Chapter~\ref{ch:virial_identity}: What the kinetic half is: the thermodynamic identity}
\iamfsentry{56}{2\langle K\rangle+\langle V\rangle=0,\qquad\langle K\rangle=\tfrac12|\langle V\rangle|,}{Virial theorem for $1/r$ binding ($k=-1$): the kinetic half.}{\derived{} (Clausius; Euler's theorem)}{Eq.~\eqref{eq:vi_virial}; Chapter~\ref{ch:virial_identity}, section ``What is $K$?''}
\iamfsentry{57}{Q=E_i-E_f=-(\langle K\rangle+\langle V\rangle).}{}{\derived{}}{Eq.~\eqref{eq:vi_Q}; Chapter~\ref{ch:virial_identity}, section ``Derivation''}
\iamfsentry{58}{\Delta S\ge Q/T .}{}{\derived{}}{Eq.~\eqref{eq:vi_dS}; Chapter~\ref{ch:virial_identity}, section ``Derivation''}
\iamfsentry{59}{E_{\rm Landauer}=T\times\frac QT=Q .}{}{\derived{}}{Eq.~\eqref{eq:vi_EL}; Chapter~\ref{ch:virial_identity}, section ``Derivation''}
\iamfsentry{60}{\boxed{\;\langle K\rangle=Q=T\Delta S_{\min}=E_{\rm Landauer}=\tfrac12|\langle V\rangle|\;}}{}{\derived{} (sympy check: \texttt{verify\_virial\_papers.py}, section F)}{Eq.~\eqref{eq:vi_identity}; Chapter~\ref{ch:virial_identity}, section ``Derivation''}
\section{\texorpdfstring{Part~\ref{part:2}}{Part 2}: Informational Actualization of General Relativity: The Cosmological Dynamic}
\subsection*{Chapter~\ref{ch:virial}: The cosmological coupling $\beta_m=\Omega_m/2$}
\iamfsentry{61}{-dE=T_H\,d\!\left(S_{\rm geo}+S_{\rm info}\right).}{Horizon first law with the record entropy added (as Eq.~\eqref{eq:firstlaw}).}{\conjecture{} (Assumption~2)}{Eq.~\eqref{eq:vc_firstlaw}; Chapter~\ref{ch:virial}, section ``From the theorem to the cosmological coupling''; see also (\ref{app:fs:125}), (\ref{app:fs:202})}
\iamfsentry{62}{\beta_m=\frac{\Omega_m}{2}=0.15765\qquad(\Omega_m=0.3153,\ \text{Planck 2018}).}{The cosmological coupling from the virial half, with Planck 2018 $\Omega_m=0.3153$.}{\prediction{}}{Eq.~\eqref{eq:vc_beta}; Chapter~\ref{ch:virial}, section ``From the theorem to the cosmological coupling''; see also (\ref{app:fs:55})}
\iamfsentry{63}{\beta_m=\Omega_m\times f_{\rm coll}\times\eta_{\rm vir},}{}{\derived{}}{Eq.~\eqref{eq:vc_decompose}; Chapter~\ref{ch:virial}, section ``The half in a real ensemble of halos''}
\iamfsentry{64}{\eta_{\rm vir}=\frac{1}{2f_{\rm coll}}=\frac{1}{2\times0.62}=0.81 .}{}{\calc{}}{Eq.~\eqref{eq:vc_eta}; Chapter~\ref{ch:virial}, section ``The half in a real ensemble of halos''}
\iamfsentry{65}{n-\tfrac92=-1\quad\Longrightarrow\quad n=\tfrac72\qquad\text{(analytical, pure matter domination).}}{}{\derived{}}{Eq.~\eqref{eq:vc_n}; Chapter~\ref{ch:virial}, section ``The exponent of the record rate''}
\iamfsentry{66}{E(a)=\exp\!\left(1-\frac1a\right)=e^{-z}.}{The activation function in scale factor and redshift (as Eq.~\eqref{eq:Ea}).}{\derived{} (given the record-density constraint, a \conjecture)}{Eq.~\eqref{eq:vc_E}; Chapter~\ref{ch:virial}, section ``The activation function and the growth coupling''}
\iamfsentry{67}{\ddot\delta_m+2H\dot\delta_m-\tfrac32\Omega_mH_0^2a^{-3}\mu(a)\,\delta_m=0,\qquad
\mu(a)=\frac{H^2_{\Lambda{\rm CDM}}(a)}{H^2_{\Lambda{\rm CDM}}(a)+\beta_mE(a)H_0^2},}{Linear growth equation in the $\mu$--$\Sigma$ form used by the Level~1 chains, with the growth coupling $\mu(a)$.}{\derived{} (given Eq.~\eqref{eq:Hm} read as an effective Newton constant; that reading is a \conjecture)}{Eq.~\eqref{part2:eq:mu_virial}; Chapter~\ref{ch:virial}, section ``The activation function and the growth coupling''}
\iamfsentry{68}{f\sigma_8^{\rm IAM}(z)=f_{\rm IAM}(z)\,\sigma_8\,\frac{D_{\rm IAM}(z)}{D_{\rm IAM}(0)},}{}{\prediction{}}{Eq.~\eqref{eq:vc_fs8}; Chapter~\ref{ch:virial}, section ``Growth and the amplitude of structure''}
\iamfsentry{69}{H_0^{(m)}=H_0^{(\gamma)}\sqrt{1+\beta_m}=67.16\times1.0759=72.26\ \text{km\,s}^{-1}\text{Mpc}^{-1}}{Matter-sector Hubble constant from the photon-sector value of the Level~2 chain.}{\calc{} from the Level~2 chain value of $H_0$ (\measured); that matter-sector probes read this rate is a \conjecture{} (Ch.~\ref{ch:dual})}{Eq.~\eqref{eq:vc_h0}; Chapter~\ref{ch:virial}, section ``Two expansion rates''}
\iamfsentry{70}{\Omega_m^{\rm growth}(z)=\Omega_m\,\mu(z),}{}{\calc{}}{Eq.~\eqref{eq:vc_omgrowth}; Chapter~\ref{ch:virial}, section ``Growth and geometry give different $\Omega_m$''}
\iamfsentry{71}{R(a)=\frac{\Omega_ma^{-3}}{\beta_mE(a)},}{}{\calc{}}{Eq.~\eqref{eq:vc_R}; Chapter~\ref{ch:virial}, section ``The virial ratio through cosmic history''}
\iamfsentry{72}{R(1)=\frac{\Omega_m}{\beta_m\,E(1)}=\frac{\Omega_m}{(\Omega_m/2)\cdot1}=2 .}{}{\calc{}}{Eq.~\eqref{eq:vc_R1}; Chapter~\ref{ch:virial}, section ``The virial ratio through cosmic history''}
\subsection*{Chapter~\ref{ch:virial_tests}: Tests of the virial coupling}
\iamfsentry{73}{E_G(z)\equiv\frac{\Upsilon_{gm}(k,z)}{\beta(z)\,\Upsilon_{gg}(k,z)},}{}{\derived{}}{Eq.~\eqref{eq:vt_eg}; Chapter~\ref{ch:virial_tests}, section ``The $E_G$ statistic''}
\iamfsentry{74}{E_G^{\rm IAM}(z)=\frac{\Omega_m}{f_{\rm IAM}(z)},}{}{\derived{}}{Eq.~\eqref{eq:vt_egiam}; Chapter~\ref{ch:virial_tests}, section ``The $E_G$ statistic''}
\iamfsentry{75}{R_1=Y_{\rm SZ}/M_{\rm lens},\qquad R_2=L_X/Y_{\rm SZ},\qquad R_3=L_X/M_{\rm lens}.}{}{\prediction{}}{Eq.~\eqref{eq:vt_R}; Chapter~\ref{ch:virial_tests}, section ``The three-way cluster test''}
\subsection*{Chapter~\ref{ch:theory}: Horizon thermodynamics and gravitational decoherence}
\iamfsentry{76}{|\psi_S\rangle\otimes|E_0\rangle\;\longrightarrow\;\sum_i c_i\,|s_i\rangle\otimes|E_i\rangle ,}{Decoherence: a system--environment product state evolves into an entangled sum over pointer states.}{\derived{} (standard quantum mechanics)}{Eq.~\eqref{eq:th:entangle}; Chapter~\ref{ch:theory}, section ``Decoherence, information, and irreversibility''; see also (\ref{app:fs:1})}
\iamfsentry{77}{\rho_S=\mathrm{Tr}_E|\Psi\rangle\langle\Psi|=\sum_{i,j}c_ic_j^*\,\langle E_j|E_i\rangle\,|s_i\rangle\langle s_j| .}{}{\derived{}}{Chapter~\ref{ch:theory}, section ``Decoherence, information, and irreversibility''}
\iamfsentry{78}{\rho_S\;\longrightarrow\;\sum_i|c_i|^2\,|s_i\rangle\langle s_i| .}{}{\derived{}}{Eq.~\eqref{eq:th:diag}; Chapter~\ref{ch:theory}, section ``Decoherence, information, and irreversibility''}
\iamfsentry{79}{I=\log_2N\ \text{bits}.}{}{\observed{}}{Chapter~\ref{ch:theory}, section ``Decoherence, information, and irreversibility''; see also (\ref{app:fs:3})}
\iamfsentry{80}{\Delta E_{\rm Landauer}=k_BT\ln2}{}{\observed{}}{Eq.~\eqref{eq:th:landauer}; Chapter~\ref{ch:theory}, section ``Decoherence, information, and irreversibility''; see also (\ref{app:fs:4})}
\iamfsentry{81}{T=\frac{\hbar\kappa}{2\pi k_B}.}{}{\derived{}}{Eq.~\eqref{eq:th:unruh}; Chapter~\ref{ch:theory}, section ``Gravity as horizon thermodynamics''}
\iamfsentry{82}{\delta Q=\int_{\mathcal H}T_{ab}\,\chi^a\,d\Sigma^b=-\kappa\int_{\mathcal H}\lambda\,T_{ab}k^ak^b\,d\lambda\,dA ,}{Heat flux across a local Rindler horizon, written with the boost Killing vector.}{\derived{} (Jacobson 1995)}{Eq.~\eqref{eq:th:dQ}; Chapter~\ref{ch:theory}, section ``Gravity as horizon thermodynamics''; see also (\ref{app:fs:7})}
\iamfsentry{83}{\delta S=\eta\,\delta A=\eta\int_{\mathcal H}\theta\,d\lambda\,dA ,}{}{\derived{}}{Eq.~\eqref{eq:th:dS}; Chapter~\ref{ch:theory}, section ``Gravity as horizon thermodynamics''}
\iamfsentry{84}{\frac{d\theta}{d\lambda}=-\frac12\theta^2-\sigma_{ab}\sigma^{ab}-R_{ab}k^ak^b .}{}{\derived{}}{Eq.~\eqref{eq:th:raych}; Chapter~\ref{ch:theory}, section ``Gravity as horizon thermodynamics''}
\iamfsentry{85}{\delta A=-\int_{\mathcal H}\lambda\,R_{ab}k^ak^b\,d\lambda\,dA .}{}{\derived{}}{Eq.~\eqref{eq:th:dA}; Chapter~\ref{ch:theory}, section ``Gravity as horizon thermodynamics''}
\iamfsentry{86}{-\kappa\int\lambda\,T_{ab}k^ak^b\,d\lambda\,dA=\frac{\hbar\kappa}{2\pi}\,\eta\Big(-\int\lambda\,R_{ab}k^ak^b\,d\lambda\,dA\Big).}{}{\derived{}}{Eq.~\eqref{eq:th:clausius}; Chapter~\ref{ch:theory}, section ``Gravity as horizon thermodynamics''}
\iamfsentry{87}{T_{ab}=\frac{\hbar\eta}{2\pi}\big(R_{ab}+f\,g_{ab}\big),}{Clausius relation for every null vector forces $T_{ab}$ proportional to $R_{ab}$ up to a term in $g_{ab}$.}{\derived{} (Jacobson 1995)}{Eq.~\eqref{eq:th:Tab}; Chapter~\ref{ch:theory}, section ``Gravity as horizon thermodynamics''}
\iamfsentry{88}{R_{ab}-\frac12Rg_{ab}+\Lambda g_{ab}=\frac{8\pi G}{c^4}T_{ab},\qquad G=\frac{c^3}{4\hbar\eta}.}{The Einstein equations as an equation of state, with $G$ fixed by $\eta$.}{\derived{} (Jacobson 1995)}{Eq.~\eqref{eq:th:einstein}; Chapter~\ref{ch:theory}, section ``Gravity as horizon thermodynamics''}
\iamfsentry{89}{\kappa\tau\sim\kappa\tau+2\pi\quad\Longrightarrow\quad\tau\sim\tau+\frac{2\pi}{\kappa}.}{}{\derived{}}{Eq.~\eqref{eq:th:period}; Chapter~\ref{ch:theory}, section ``Gravity as horizon thermodynamics''}
\iamfsentry{90}{\underbrace{\frac{\hbar\eta}{2\pi}}_{\text{Rindler geometry}}=\underbrace{\frac{c^3}{8\pi G}}_{\text{Einstein geometry}} .}{}{\derived{}}{Eq.~\eqref{eq:th:structural}; Chapter~\ref{ch:theory}, section ``Gravity as horizon thermodynamics''; see also (\ref{app:fs:394})}
\iamfsentry{91}{\eta=\frac{c^3}{4\hbar G}=\frac{1}{4\ell_P^2},\qquad\frac14=\frac{2\pi}{8\pi},\qquad \ell_P^2=\frac{\hbar G}{c^3}.}{}{\derived{}}{Eq.~\eqref{eq:th:eta}; Chapter~\ref{ch:theory}, section ``Gravity as horizon thermodynamics''}
\iamfsentry{92}{\delta A_{\min}=\frac{8\pi G}{\kappa c^3}\cdot\frac{\hbar\kappa}{2\pi}=\frac{4\hbar G}{c^3}=4\ell_P^2}{Smallest horizon area that holds one unit of entropy $k_B$: $4\ell_P^2$, independent of surface gravity.}{\derived}{Eq.~\eqref{eq:th:dAmin}; Chapter~\ref{ch:theory}, section ``Gravity as horizon thermodynamics''}
\iamfsentry{93}{\tilde r_A=\frac1H ,}{}{\derived{}}{Eq.~\eqref{eq:th:rA}; Chapter~\ref{ch:theory}, section ``Gravity as horizon thermodynamics''}
\iamfsentry{94}{S_{\rm geo}=\frac{A_H}{4G}=\frac{\pi}{GH^2},}{}{\derived{}}{Eq.~\eqref{eq:th:Sgeo}; Chapter~\ref{ch:theory}, section ``Gravity as horizon thermodynamics''}
\iamfsentry{95}{T_H=\frac{H}{2\pi}.}{}{\derived{}}{Eq.~\eqref{eq:th:TH}; Chapter~\ref{ch:theory}, section ``Gravity as horizon thermodynamics''}
\iamfsentry{96}{E=\frac{\tilde r_A}{2G}=\frac{4\pi}{3}\rho\,\tilde r_A^3=\frac{4\pi}{3}\frac{\rho}{H^3},}{}{\derived{}}{Eq.~\eqref{eq:th:EMS}; Chapter~\ref{ch:theory}, section ``Gravity as horizon thermodynamics''}
\iamfsentry{97}{dS_{\rm geo}=-\frac{2\pi}{GH^3}\,dH .}{}{\derived{}}{Eq.~\eqref{eq:th:dSgeo}; Chapter~\ref{ch:theory}, section ``Gravity as horizon thermodynamics''}
\iamfsentry{98}{-dE=A_H\,(\rho+P)\,H\tilde r_A\,dt=4\pi\tilde r_A^3(\rho+P)H\,dt=\frac{4\pi(\rho+P)}{H^2}\,dt .}{}{\derived{}}{Eq.~\eqref{eq:th:dE}; Chapter~\ref{ch:theory}, section ``Gravity as horizon thermodynamics''}
\iamfsentry{99}{\frac{4\pi(\rho+P)}{H^2}=\frac{H}{2\pi}\cdot\frac{2\pi}{GH^3}\cdot(-\dot H)=-\frac{\dot H}{GH^2},}{}{\derived{}}{Eq.~\eqref{eq:th:firstlaw}; Chapter~\ref{ch:theory}, section ``Gravity as horizon thermodynamics''}
\iamfsentry{100}{\dot H=-4\pi G(\rho+P).}{Cai--Kim: the first law at the apparent horizon gives the second Friedmann equation.}{\derived{} (Cai and Kim 2005)}{Eq.~\eqref{eq:th:Hdot}; Chapter~\ref{ch:theory}, section ``Gravity as horizon thermodynamics''; see also (\ref{app:fs:14}), (\ref{app:fs:169})}
\iamfsentry{101}{H^2=\frac{8\pi G}{3}\rho+\frac\Lambda3 ,}{Matter-sector expansion rate $H_{\rm IAM}$: the rate in the friction term of matter perturbations; the background stays $\Lambda$CDM.}{\derived{} (given the completed first law, Eq.~\eqref{eq:firstlaw}, a \conjecture, and the sector split, a \conjecture)}{Eq.~\eqref{eq:th:friedmann}; Chapter~\ref{ch:theory}, section ``Gravity as horizon thermodynamics''; see also (\ref{app:fs:170})}
\iamfsentry{102}{\Delta E\ge k_BT_H\ln2=\frac{\hbar H}{2\pi}\ln2 .}{}{\calc{}}{Eq.~\eqref{eq:th:bitcost}; Chapter~\ref{ch:theory}, section ``Gravity as horizon thermodynamics''; see also (\ref{app:fs:174})}
\iamfsentry{103}{S_{\rm total}=S_{\rm geo}+S_{\rm info}.}{}{\conjecture{}}{Eq.~\eqref{eq:th:Stotal}; Chapter~\ref{ch:theory}, section ``Informational entropy from gravitational decoherence''; see also (\ref{app:fs:13}), (\ref{app:fs:404})}
\iamfsentry{104}{\dot I\propto\rho_m(a)\,D(a)^n\,f(a)\,H(a),}{Assumed rate of information production by structure formation (as Eq.~\eqref{eq:rate}).}{\conjecture{} (assumed form of the writing rate)}{Eq.~\eqref{eq:Idot}; Chapter~\ref{ch:theory}, section ``Informational entropy from gravitational decoherence''; see also (\ref{app:fs:19})}
\iamfsentry{105}{\frac{dS_{\rm info}}{dt}\propto\frac{\dot I}{T_H}\cdot\frac{1}{A_H}.}{Accumulation of record entropy on the horizon: Landauer factor $1/T_H$ and holographic factor $1/A_H$.}{\conjecture{} (assumed accounting: Landauer factor $1/T$ and storage per unit horizon area, Assumption~2)}{Eq.~\eqref{eq:dSdt}; Chapter~\ref{ch:theory}, section ``Informational entropy from gravitational decoherence''}
\iamfsentry{106}{2K+U=0,}{}{\derived{}}{Eq.~\eqref{eq:th:virial}; Chapter~\ref{ch:theory}, section ``Informational entropy from gravitational decoherence''}
\iamfsentry{107}{\beta_m=\frac{\Omega_m}{2}.}{The coupling from the virial partition $\rho_{\rm info}(1)=\tfrac12\rho_m$.}{\derived{} (given the virial partition, Section~\ref{sec:th:coupling})}{Eq.~\eqref{eq:th:beta}; Chapter~\ref{ch:theory}, section ``Informational entropy from gravitational decoherence''; see also (\ref{app:fs:147})}
\iamfsentry{108}{\begin{gathered}H(a)\propto a^{-3/2},\\
A_H(a)=\frac{4\pi}{H^2}\propto a^{3},\\
T_H(a)=\frac{H}{2\pi}\propto a^{-3/2},\\
D(a)\propto a.\end{gathered}}{}{\derived{}}{Eqs.~\eqref{eq:th:Hmd}, \eqref{eq:th:Dmd}; Chapter~\ref{ch:theory}, section ``Derivation of the activation function''}
\iamfsentry{109}{\frac{dS_{\rm info}}{d\ln a}\propto\frac{\rho_m\,D^n\,f}{T_H\,A_H}.}{}{\derived{}}{Eq.~\eqref{eq:th:perlna}; Chapter~\ref{ch:theory}, section ``Derivation of the activation function''}
\iamfsentry{110}{\frac{dS_{\rm info}}{d\ln a}\propto\frac{a^{-3}\cdot a^n\cdot1}{a^{-3/2}\cdot a^{3}}=a^{\,n-9/2}.}{}{\derived{}}{Eq.~\eqref{eq:th:power}; Chapter~\ref{ch:theory}, section ``Derivation of the activation function''}
\iamfsentry{111}{\frac{dS_{\rm info}}{da}\propto a^{\,n-11/2}.}{}{\derived{}}{Eq.~\eqref{eq:th:perda}; Chapter~\ref{ch:theory}, section ``Derivation of the activation function''}
\iamfsentry{112}{S_{\rm info}(a)\propto\int^{a}a'^{\,n-11/2}\,da'=\frac{a^{\,n-9/2}}{n-9/2}.}{Matter-era scaling of the record entropy per $\ln a$ and per $a$.}{\derived{}}{Eq.~\eqref{eq:Sn}; Chapter~\ref{ch:theory}, section ``Derivation of the activation function''}
\iamfsentry{113}{S_{\rm info}(a)\propto-\frac1a+{\rm const}.}{}{\derived{}}{Eq.~\eqref{eq:th:target}; Chapter~\ref{ch:theory}, section ``Derivation of the activation function''}
\iamfsentry{114}{n-\frac92=-1\quad\Longrightarrow\quad n=\frac72 .}{}{\derived{} (given the form)}{Eq.~\eqref{eq:th:n}; Chapter~\ref{ch:theory}, section ``Derivation of the activation function''; see also (\ref{app:fs:37})}
\iamfsentry{115}{S_{\rm info}(a)\propto-\frac1a+C .}{}{\conjecture{}}{Eq.~\eqref{eq:th:S72}; Chapter~\ref{ch:theory}, section ``Derivation of the activation function''}
\iamfsentry{116}{-dE_{\rm info}=T_H\,dS_{\rm info}.}{}{\conjecture{}}{Eq.~\eqref{eq:th:dEinfo}; Chapter~\ref{ch:theory}, section ``Derivation of the activation function''}
\iamfsentry{117}{\dot\rho_{\rm info}=\rho_{\rm info}\cdot\frac Ha .}{}{\conjecture{} (the identification of the energy density's growth rate with the record rate)}{Eq.~\eqref{eq:th:rhodot}; Chapter~\ref{ch:theory}, section ``Derivation of the activation function''; see also (\ref{app:fs:140})}
\iamfsentry{118}{\rho_{\rm info}(a)\propto\exp\!\Big(\int\frac Ha\,\frac{da}{aH}\Big)=\exp\!\Big(\int\frac{da}{a^2}\Big)=\exp\!\Big(-\frac1a\Big).}{}{\conjecture{}}{Eq.~\eqref{eq:th:rhoint}; Chapter~\ref{ch:theory}, section ``Derivation of the activation function''}
\iamfsentry{119}{\Delta H^2\propto\exp\!\Big(C-\frac1a\Big).}{}{\conjecture{}}{Eq.~\eqref{eq:th:dH2}; Chapter~\ref{ch:theory}, section ``Derivation of the activation function''}
\iamfsentry{120}{E(1)=\exp(C-1)=1\quad\Longrightarrow\quad C=1 .}{}{\derived{}}{Eq.~\eqref{eq:th:C}; Chapter~\ref{ch:theory}, section ``Derivation of the activation function''}
\iamfsentry{121}{\boxed{E(a)=\exp\!\Big(1-\frac1a\Big)}}{The activation function, boxed result of the decoherence constraint with $E(1)=1$.}{\derived{} (given the record-density constraint, a \conjecture)}{Eq.~\eqref{part2:eq:Ea}; Chapter~\ref{ch:theory}, section ``Derivation of the activation function''}
\iamfsentry{122}{\frac{\text{structure formation rate}}{A_H}\propto\frac{D(a)}{A_H(a)}\propto\frac{a}{a^3}=\frac{1}{a^2}.}{}{\derived{}}{Eq.~\eqref{eq:th:surfdens}; Chapter~\ref{ch:theory}, section ``Derivation of the activation function''}
\iamfsentry{123}{E(z)=\exp\big(1-(1+z)\big)=\exp(-z)=e\cdot e^{-(1+z)}.}{}{\derived{}}{Eq.~\eqref{eq:th:Ez}; Chapter~\ref{ch:theory}, section ``Derivation of the activation function''}
\iamfsentry{124}{-dE=T\,dS_{\rm geo}+T\,dS_{\rm info}.}{}{\calc{}}{Eq.~\eqref{eq:th:firstlaw2}; Chapter~\ref{ch:theory}, section ``Derivation of the activation function''}
\iamfsentry{125}{-dE=T_H\,d(S_{\rm geo}+S_{\rm info}).}{Horizon first law with the record entropy (the sector split enters through (\ref{app:fs:60})).}{\conjecture{} (Assumption~2)}{Eq.~\eqref{eq:th:firstlaw3}; Chapter~\ref{ch:theory}, section ``The dual-sector split and the $\mu$--$\Sigma$ mapping''; see also (\ref{app:fs:61}), (\ref{app:fs:202})}
\iamfsentry{126}{H^2=\frac{8\pi G}{3}\rho+\frac\Lambda3+\beta_mE(a)H_0^2 .}{Matter-sector expansion rate $H_{\rm IAM}$: the rate in the friction term of matter perturbations; the background stays $\Lambda$CDM.}{\derived{} (given the completed first law, Eq.~\eqref{eq:firstlaw}, a \conjecture, and the sector split, a \conjecture)}{Eq.~\eqref{eq:HIAM}; Chapter~\ref{ch:theory}, section ``The dual-sector split and the $\mu$--$\Sigma$ mapping''}
\iamfsentry{127}{\rho_{\rm info}(a)=\frac{3H_0^2}{8\pi G}\,\beta_mE(a),}{}{\derived{}}{Eq.~\eqref{eq:th:rhoinfo}; Chapter~\ref{ch:theory}, section ``The dual-sector split and the $\mu$--$\Sigma$ mapping''}
\iamfsentry{128}{\begin{aligned}H^2&=\frac{8\pi G}{3}(\rho_m+\rho_{\rm info})+\frac\Lambda3\\&=\frac{8\pi G}{3}\rho_m+\frac\Lambda3+\frac{8\pi G}{3}\cdot\frac{3H_0^2}{8\pi G}\beta_mE(a)
\\&=\frac{8\pi G}{3}\rho_m+\frac\Lambda3+\beta_mE(a)H_0^2 .\end{aligned}}{}{\derived{}}{Eq.~\eqref{eq:th:cancel}; Chapter~\ref{ch:theory}, section ``The dual-sector split and the $\mu$--$\Sigma$ mapping''}
\iamfsentry{129}{\begin{gathered}k^2\Psi=-4\pi G\,\mu(a)\,a^2\bar\rho_m\delta_m,\\
k^2(\Phi+\Psi)=-8\pi G\,\Sigma(a)\,a^2\bar\rho_m\delta_m.\end{gathered}}{Modified Poisson equations of the $\mu$--$\Sigma$ framework in Fourier space; GR is $\mu=\Sigma=1$.}{definition ($\mu$--$\Sigma$ parametrization)}{Eqs.~\eqref{eq:th:mudef}, \eqref{eq:th:sigmadef}; Chapter~\ref{ch:theory}, section ``The dual-sector split and the $\mu$--$\Sigma$ mapping''; see also (\ref{app:fs:188})}
\iamfsentry{130}{\mu(a)=\frac{H^2_{\Lambda\rm CDM}(a)}{H^2_{\Lambda\rm CDM}(a)+\beta_mE(a)H_0^2}.}{}{\derived{}}{Eq.~\eqref{eq:th:mu}; Chapter~\ref{ch:theory}, section ``The dual-sector split and the $\mu$--$\Sigma$ mapping''; see also (\ref{app:fs:333})}
\iamfsentry{131}{\mu_0\equiv\mu(z{=}0)-1=\frac{-\beta_m}{1+\beta_m}=-0.136 .}{}{\derived{} ($\mu(0)=1/(1+\beta_m)=0.8638$, $\mu_0=-0.1362$ \calc)}{Eq.~\eqref{eq:th:mu0}; Chapter~\ref{ch:theory}, section ``The dual-sector split and the $\mu$--$\Sigma$ mapping''}
\iamfsentry{132}{\Sigma(a)=1\quad\text{(exact)}.}{}{\derived{}}{Eq.~\eqref{eq:th:Sigma}; Chapter~\ref{ch:theory}, section ``The dual-sector split and the $\mu$--$\Sigma$ mapping''}
\iamfsentry{133}{\varphi(t)\equiv\ln E(a)=1-\frac1a .}{}{\derived{}}{Eq.~\eqref{eq:th:phi}; Chapter~\ref{ch:theory}, section ``Variational formulation''}
\iamfsentry{134}{\dot\varphi=\frac{\dot a}{a^2}=\frac Ha .}{}{\derived{}}{Eq.~\eqref{eq:th:phidot}; Chapter~\ref{ch:theory}, section ``Variational formulation''}
\iamfsentry{135}{S=S_{\rm EH}+S_\Lambda+S_{\rm matter}+S_{\rm info},}{}{\conjecture{}}{Eq.~\eqref{eq:th:Stot}; Chapter~\ref{ch:theory}, section ``Variational formulation''}
\iamfsentry{136}{S_{\rm info}=\int\Big[-\frac{3H_0^2}{8\pi G}\beta_m e^{\varphi}+\lambda\Big(\dot\varphi-\frac Ha\Big)\Big]\sqrt{-g}\,d^4x .}{Constrained action for the record term in the variational form.}{\conjecture{} (the action)}{Eq.~\eqref{eq:th:action}; Chapter~\ref{ch:theory}, section ``Variational formulation''}
\iamfsentry{137}{L_{\rm info}=-Na^3\rho_{\rm info}(\varphi)+\lambda\,a^3\Big(\varphi'-\frac{a'}{a^2}\Big):}{}{\derived{}}{Eq.~\eqref{eq:th:Lmini}; Chapter~\ref{ch:theory}, section ``Variational formulation''}
\iamfsentry{138}{\frac{d}{dt}\big(a^3\lambda\big)=-a^3\rho_{\rm info}\qquad\Longleftrightarrow\qquad\dot\lambda+3H\lambda=-\frac{3H_0^2}{8\pi G}\beta_me^{\varphi},}{}{\derived{}}{Eq.~\eqref{eq:th:lambdadot}; Chapter~\ref{ch:theory}, section ``Variational formulation''}
\iamfsentry{139}{H^2=\frac{8\pi G}{3}(\rho_m+\rho_r)+\frac\Lambda3+\beta_me^{\varphi}H_0^2 .}{}{\derived{}}{Eq.~\eqref{eq:th:Hvar}; Chapter~\ref{ch:theory}, section ``Variational formulation''}
\iamfsentry{140}{\dot\rho_{\rm info}=\rho_{\rm info}\cdot\frac Ha .}{}{\derived{}}{Eq.~\eqref{eq:th:rhodot2}; Chapter~\ref{ch:theory}, section ``Variational formulation''; see also (\ref{app:fs:117})}
\iamfsentry{141}{w_{\rm info}(a)=-1-\frac{1}{3a}.}{}{\derived{}}{Eq.~\eqref{eq:th:winfo}; Chapter~\ref{ch:theory}, section ``Variational formulation''}
\iamfsentry{142}{w^{\rm eff}_{\rm DE}(a)=\frac{P_\Lambda+P_{\rm info}}{\rho_\Lambda+\rho_{\rm info}}=\frac{-\rho_\Lambda+w_{\rm info}\,\rho_{\rm info}}{\rho_\Lambda+\rho_{\rm info}}.}{}{\derived{}}{Eq.~\eqref{eq:th:weff}; Chapter~\ref{ch:theory}, section ``Variational formulation''}
\iamfsentry{143}{w^{\rm eff}_{\rm DE}(1)=\frac{-\Omega_\Lambda-\tfrac43\beta_m}{\Omega_\Lambda+\beta_m}=-\frac{1-\Omega_m/3}{1-\Omega_m/2}=-1.062 .}{}{\derived{}}{Chapter~\ref{ch:theory}, section ``Variational formulation''}
\iamfsentry{144}{w_0=-1.062,\qquad w_a=-\frac{dw^{\rm eff}}{da}\Big|_{a=1}=-\frac{\Omega_m^2}{3(2-\Omega_m)^2}=-0.012 .}{CPL parameters matched to $w_{\rm eff}$ and its slope at $a=1$.}{\derived{} (form)}{Chapter~\ref{ch:theory}, section ``Variational formulation''}
\iamfsentry{145}{\langle T\rangle=-\frac12\langle V\rangle .}{}{\conjecture{}}{Eq.~\eqref{eq:th:virialavg}; Chapter~\ref{ch:theory}, section ``The coupling constant: $\beta_m=\Omega_m/2$''}
\iamfsentry{146}{\rho_{\rm info}(a{=}1)=\frac12\rho_m=\frac{\Omega_m}{2}\rho_{\rm crit}.}{}{\derived{}}{Eq.~\eqref{eq:th:rhoinfo1}; Chapter~\ref{ch:theory}, section ``The coupling constant: $\beta_m=\Omega_m/2$''}
\iamfsentry{147}{\beta_m=\frac{\Omega_m}{2}.}{}{\derived{}}{Eq.~\eqref{eq:th:betam}; Chapter~\ref{ch:theory}, section ``The coupling constant: $\beta_m=\Omega_m/2$''; see also (\ref{app:fs:107})}
\iamfsentry{148}{f^{\rm ST}_{\rm coll}=0.64,\qquad f^{\rm Tinker}_{\rm coll}=0.71 .}{}{\calc{}}{Eq.~\eqref{eq:th:fcoll}; Chapter~\ref{ch:theory}, section ``The coupling constant: $\beta_m=\Omega_m/2$''}
\iamfsentry{149}{\beta_m=\Omega_m\cdot f_{\rm coll}\cdot\eta_{\rm vir}}{}{\derived{}}{Eq.~\eqref{eq:th:betadecomp}; Chapter~\ref{ch:theory}, section ``The coupling constant: $\beta_m=\Omega_m/2$''}
\iamfsentry{150}{I(a)=\int_0^{t(a)}\frac{R(a')}{T_H(a')\,A_H(a')}\,dt'=\int_0^a\frac{R(a')}{T_H(a')\,A_H(a')}\,\frac{da'}{a'H(a')}\;\propto\;\int_0^a\frac{R(a')}{a'}\,da',}{}{\derived{}}{Eq.~\eqref{eq:th:Iint}; Chapter~\ref{ch:theory}, section ``Numerical verification''}
\iamfsentry{151}{f_{\rm ST}(\nu)=A\sqrt{\frac{2q}{\pi}}\Big[1+(q\nu^2)^{-p}\Big]\,\nu\,e^{-q\nu^2/2},}{}{\calc{}}{Eq.~\eqref{eq:th:fST}; Chapter~\ref{ch:theory}, section ``Numerical verification''}
\iamfsentry{152}{\delta\varphi=0}{}{\derived{}}{Eq.~\eqref{eq:th:dphi}; Chapter~\ref{ch:theory}, section ``Perturbation theory: $\mu(a)<1$, $\Sigma(a)=1$''}
\iamfsentry{153}{k^2\Psi=-4\pi Ga^2\rho_m\delta_m ;}{}{\derived{}}{Eq.~\eqref{eq:th:poisson}; Chapter~\ref{ch:theory}, section ``Perturbation theory: $\mu(a)<1$, $\Sigma(a)=1$''}
\iamfsentry{154}{\Psi=\Phi}{}{\derived{}}{Eq.~\eqref{eq:th:noaniso}; Chapter~\ref{ch:theory}, section ``Perturbation theory: $\mu(a)<1$, $\Sigma(a)=1$''}
\iamfsentry{155}{\ddot\delta_m+2H_{\rm IAM}\dot\delta_m-4\pi G\rho_m\delta_m=0 .}{Growth equation of matter perturbations with friction from the matter-sector rate.}{\derived{} (standard perturbation theory with $H_m$ as the friction term; placing $H_m$ there is the sector split, a \conjecture)}{Eq.~\eqref{eq:th:growth}; Chapter~\ref{ch:theory}, section ``Perturbation theory: $\mu(a)<1$, $\Sigma(a)=1$''}
\iamfsentry{156}{\mu(a)=\frac{E^2_{\Lambda\rm CDM}(a)}{E^2_{\rm IAM}(a)}<1,\qquad\Sigma(a)=1,}{}{\derived{}}{Eq.~\eqref{eq:th:mu2}; Chapter~\ref{ch:theory}, section ``Perturbation theory: $\mu(a)<1$, $\Sigma(a)=1$''}
\iamfsentry{157}{\ddot D_2+2H_{\rm IAM}\dot D_2-4\pi G\bar\rho_m\,D_2=-4\pi G\bar\rho_m\,D_1^2,\qquad4\pi G\bar\rho_m=\tfrac32H_0^2\Omega_ma^{-3},}{}{\derived{}}{Eq.~\eqref{eq:th:D2}; Chapter~\ref{ch:theory}, section ``Perturbation theory: $\mu(a)<1$, $\Sigma(a)=1$''}
\iamfsentry{158}{F_2(\mathbf k_1,\mathbf k_2)=\frac57+\frac{\hat{\mathbf k}_1\cdot\hat{\mathbf k}_2}{2}\Big(\frac{k_1}{k_2}+\frac{k_2}{k_1}\Big)+\frac27(\hat{\mathbf k}_1\cdot\hat{\mathbf k}_2)^2}{}{\derived{}}{Eq.~\eqref{eq:th:F2}; Chapter~\ref{ch:theory}, section ``Perturbation theory: $\mu(a)<1$, $\Sigma(a)=1$''}
\iamfsentry{159}{\nabla_\mu T^{\mu\nu}_{\rm total}=\nabla_\mu\big(T^{\mu\nu}_{\rm matter}+T^{\mu\nu}_\Lambda+T^{\mu\nu}_{\rm info}\big)=0 .}{}{\interp{}}{Eq.~\eqref{eq:th:conservation}; Chapter~\ref{ch:theory}, section ``Observational validation''}
\iamfsentry{160}{w_{\rm eff}(a)=\frac{\Omega_\Lambda\cdot(-1)+\beta_mE(a)\cdot w_{\rm info}(a)}{\Omega_\Lambda+\beta_mE(a)},}{}{\calc{}}{Eq.~\eqref{eq:th:weff2}; Chapter~\ref{ch:theory}, section ``Observational validation''}
\iamfsentry{161}{H_0^{\rm sirens}=H_0^{\rm photon}\sqrt{1+\beta_m}=67.16\sqrt{1.15765}=72.26\ \text{km\,s}^{-1}\text{Mpc}^{-1},}{}{\calc{}}{Eq.~\eqref{eq:th:sirens}; Chapter~\ref{ch:theory}, section ``Observational validation''}
\iamfsentry{162}{\frac{S_{\rm info}(R)}{A(R)}\ge\frac{1}{4\ell_P^2}.}{Collapse criterion as holographic saturation: record entropy per area reaching $1/4\ell_P^2$ (met exactly at the Schwarzschild radius).}{\conjecture{}}{Eq.~\eqref{eq:th:hoop}; Chapter~\ref{ch:theory}, section ``Physical interpretation''; see also (\ref{app:fs:517})}
\iamfsentry{163}{\frac{4\pi GM^2/(\hbar c)}{4\pi\cdot4G^2M^2/c^4}=\frac{c^3}{4\hbar G}=\frac{1}{4\ell_P^2}.}{}{\derived{}}{Eq.~\eqref{eq:th:hoopcheck}; Chapter~\ref{ch:theory}, section ``Physical interpretation''}
\iamfsentry{164}{M_{\rm eq}=\frac{c^3}{4GH}.}{}{\derived{}}{Eq.~\eqref{eq:th:Meq}; Chapter~\ref{ch:theory}, section ``Physical interpretation''; see also (\ref{app:fs:40})}
\iamfsentry{165}{\Gamma=\frac{P}{k_BT_{BH}\ln2}=\frac{c^3}{1920\,GM\ln2}\ \text{bits\,s}^{-1}.}{}{\derived{}}{Eq.~\eqref{eq:th:Gamma}; Chapter~\ref{ch:theory}, section ``Physical interpretation''}
\subsection*{Chapter~\ref{ch:entropicgravity}: Entropy law or entropy source: the thermodynamic-gravity conjecture}
\iamfsentry{166}{-dE=A_H\,(\rho+P)\,H\tilde r_A\,dt=\frac{4\pi(\rho+P)}{H^2}\,dt .}{}{\derived{}}{Eq.~\eqref{eq:eg_flux}; Chapter~\ref{ch:entropicgravity}, section ``The shared starting point''}
\iamfsentry{167}{dS_{\rm BH}=\frac{d}{dt}\Big(\frac{\pi}{GH^2}\Big)dt=-\frac{2\pi}{GH^3}\,\dot H\,dt .}{}{\derived{}}{Eq.~\eqref{eq:eg_dS}; Chapter~\ref{ch:entropicgravity}, section ``The shared starting point''}
\iamfsentry{168}{-dE=T_H\,dS_{\rm BH},}{}{\derived{}}{Eq.~\eqref{eq:eg_firstlaw}; Chapter~\ref{ch:entropicgravity}, section ``The shared starting point''}
\iamfsentry{169}{\dot H=-4\pi G(\rho+P).}{Cai--Kim: the first law at the apparent horizon gives the second Friedmann equation.}{\derived{} (Cai and Kim 2005)}{Eq.~\eqref{eq:eg_Hdot}; Chapter~\ref{ch:entropicgravity}, section ``The shared starting point''; see also (\ref{app:fs:14}), (\ref{app:fs:100})}
\iamfsentry{170}{H^2=\frac{8\pi G}{3}\rho+\frac\Lambda3 ,}{}{\derived{}}{Eq.~\eqref{eq:eg_friedmann}; Chapter~\ref{ch:entropicgravity}, section ``The shared starting point''; see also (\ref{app:fs:101})}
\iamfsentry{171}{S_{\rm Barrow}=\Big(\frac{A}{A_0}\Big)^{1+\Delta/2},}{}{\conjecture{}}{Eq.~\eqref{eq:eg_barrow}; Chapter~\ref{ch:entropicgravity}, section ``Two roads from the conjecture''}
\iamfsentry{172}{S_{\rm Tsallis}=\gamma A^\delta,}{}{\conjecture{}}{Eq.~\eqref{eq:eg_tsallis}; Chapter~\ref{ch:entropicgravity}, section ``Two roads from the conjecture''}
\iamfsentry{173}{S_{\rm total}=S_{\rm BH}+S_{\rm info},}{}{\conjecture{}}{Eq.~\eqref{eq:eg_stotal}; Chapter~\ref{ch:entropicgravity}, section ``Two roads from the conjecture''}
\iamfsentry{174}{\Delta E\ge k_BT_H\ln2=\frac{\hbar H}{2\pi}\ln2}{}{\conjecture{}}{Eq.~\eqref{eq:eg_bitcost}; Chapter~\ref{ch:entropicgravity}, section ``Two roads from the conjecture''; see also (\ref{app:fs:102})}
\iamfsentry{175}{\ddot\delta+\big[2+\beta_mE(a)\big]H\dot\delta-4\pi G\bar\rho_m\delta=0,\qquad 4\pi G\bar\rho_m=\tfrac32H_0^2\Omega_ma^{-3}=\tfrac32\Omega_m(a)H^2 .}{}{\derived{}}{Eq.~\eqref{eq:eg_growth}; Chapter~\ref{ch:entropicgravity}, section ``Where the two roads lead''}
\iamfsentry{176}{\delta''+\Big[2+\beta_mE(a)+\frac{d\ln H}{d\ln a}\Big]\delta'-\frac32\,\frac{\Omega_ma^{-3}}{E^2_{\Lambda{\rm CDM}}(a)}\,\delta=0,
\qquad E^2_{\Lambda{\rm CDM}}=\Omega_ma^{-3}+\Omega_\Lambda ,}{}{\derived{}}{Eq.~\eqref{eq:eg_growthN}; Chapter~\ref{ch:entropicgravity}, section ``Where the two roads lead''}
\iamfsentry{177}{2\langle T\rangle+\langle V\rangle=0\quad\Longrightarrow\quad\langle T\rangle=\tfrac12|\langle V\rangle| .}{}{\derived{}}{Eq.~\eqref{eq:eg_virial}; Chapter~\ref{ch:entropicgravity}, section ``The coupling constant: \texorpdfstring{$\beta_m=\Omega_m/2$}{beta\_m = Omega\_m/2}''}
\iamfsentry{178}{\beta_m=\frac{\Omega_m}{2}=0.15765}{}{\derived{}}{Eq.~\eqref{eq:eg_beta}; Chapter~\ref{ch:entropicgravity}, section ``The coupling constant: \texorpdfstring{$\beta_m=\Omega_m/2$}{beta\_m = Omega\_m/2}''; see also (\ref{app:fs:319}), (\ref{app:fs:331}), (\ref{app:fs:399}), (\ref{app:fs:522})}
\iamfsentry{179}{E(a)=\exp\Big(1-\frac1a\Big).}{}{\derived{}}{Eq.~\eqref{eq:eg_Ea}; Chapter~\ref{ch:entropicgravity}, section ``The activation function \texorpdfstring{$E(a)$}{E(a)}''}
\subsection*{Chapter~\ref{ch:dual}: The matter ledger of the horizon: the dual-sector model and its chains}
\iamfsentry{180}{-dE=T_H\,d\!\left(S_{\mathrm{geo}}+S_{\mathrm{info}}\right).}{Horizon first law with the record entropy added; the record density is taken to grow as $\dot\rho_{\rm info}=\rho_{\rm info}H/a$.}{\conjecture{} (as labeled)}{Chapter~\ref{ch:dual}, section ``One term added to the horizon entropy''}
\iamfsentry{181}{\boxed{\;E(a)=\exp\!\left(1-\frac1a\right)=e^{-z}\;}}{The activation function as the solution of the record-density law, normalized to one today.}{\derived{} (given $\dot\rho_{\rm info}=\rho_{\rm info}H/a$, a \conjecture)}{Eq.~\eqref{part2:eq:Ea_chains}; Chapter~\ref{ch:dual}, section ``One term added to the horizon entropy''}
\iamfsentry{182}{\mu(a)=\frac{H^2_{\Lambda\mathrm{CDM}}(a)}{H^2_{\Lambda\mathrm{CDM}}(a)+\beta_m E(a)H_0^2},\qquad \Sigma(a)=1,}{The standard modified-growth form: $\mu(a)$ multiplies $G$ in the growth equation, $\Sigma=1$ in the lensing potential.}{\derived{} (as labeled; given Eq.~\eqref{eq:Hm}, a \conjecture)}{Eq.~\eqref{part2:eq:mu}; Chapter~\ref{ch:dual}, section ``Two sectors''}
\iamfsentry{183}{H_0^{\rm (m)}=H_0\sqrt{1+\beta},}{}{\derived{}}{Eq.~\eqref{eq:dual_free_H0m}; Chapter~\ref{ch:dual}, section ``The free-coupling fit in full''}
\subsection*{Chapter~\ref{ch:latetime}: Late-time growth suppression: confrontation with Planck and large-scale structure}
\iamfsentry{184}{\mu(a)=\frac{H^2_{\Lambda\rm CDM}(a)}{H^2_{\Lambda\rm CDM}(a)+\beta\,E(a)\,H_0^2},}{The form of $\mu(a)$ tested in the Level~1 MGCAMB chains, with $\beta=\Omega_m/2$.}{\derived{} (given Eq.~\eqref{eq:Hm} read as an effective Newton constant; that reading is a \conjecture)}{Eq.~\eqref{eq:lt_mu}; Chapter~\ref{ch:latetime}, section ``Introduction''}
\iamfsentry{185}{E(a)=\exp\!\left(1-\frac1a\right)}{}{\calc{}}{Eq.~\eqref{eq:lt_activation}; Chapter~\ref{ch:latetime}, section ``Introduction''; see also (\ref{app:fs:206}), (\ref{app:fs:318}), (\ref{app:fs:332}), (\ref{app:fs:408}), (\ref{app:fs:520})}
\iamfsentry{186}{\mu(z{=}0)=\frac{1}{1+\beta}=\frac{1}{1.15765}=0.864\qquad(\beta=\Omega_m/2=0.15765,\ \Omega_m=0.3153),}{}{\calc{}}{Eq.~\eqref{eq:lt_mu0}; Chapter~\ref{ch:latetime}, section ``Introduction''}
\iamfsentry{187}{\ln E(a)=-\int_a^1\frac{da'}{a'^2}=1-\frac1a,\qquad E(1)=1,\qquad \lim_{a\to0}E(a)=0,\qquad E=e^{-z}.}{}{\derived{} (given the $a^{-2}$ scaling of Chapter~\ref{ch:theory})}{Eq.~\eqref{eq:lt_lnE}; Chapter~\ref{ch:latetime}, section ``Introduction''}
\iamfsentry{188}{\begin{aligned}k^2\Psi&=-4\pi G\,\mu(a)\,a^2\bar\rho_m\delta_m,\\
k^2(\Phi+\Psi)&=-8\pi G\,\Sigma(a)\,a^2\bar\rho_m\delta_m,\end{aligned}}{Modified Poisson equations of the $\mu$--$\Sigma$ framework in Fourier space; GR is $\mu=\Sigma=1$.}{definition ($\mu$--$\Sigma$ parametrization)}{Eqs.~\eqref{eq:lt_poisson}, \eqref{eq:lt_lensing}; Chapter~\ref{ch:latetime}, section ``Model definition and implementation''; see also (\ref{app:fs:129})}
\iamfsentry{189}{\mu(a)=1+\mu_0\,\frac{\Omega_{\rm DE}(a)}{\Omega_\Lambda},\qquad\Sigma(a)=1+\Sigma_0\,\frac{\Omega_{\rm DE}(a)}{\Omega_\Lambda},}{The DES parametrization implemented in MGCAMB; the runs set $\mu_0=-0.13495$, $\Sigma_0=0$.}{definition (MGCAMB form); its match to Eq.~\eqref{eq:lt_mu} is \calc}{Eq.~\eqref{eq:lt_mgcamb}; Chapter~\ref{ch:latetime}, section ``Model definition and implementation''}
\subsection*{Chapter~\ref{ch:level2}: The mechanism in the Boltzmann code: a matter-sector expansion rate}
\iamfsentry{190}{H^2(a)=H_0^2\left[\Omega_ma^{-3}+\Omega_ra^{-4}+\Omega_\Lambda\right].}{The $\Lambda$CDM background expansion rate used for photons and the background in the Level~2 code.}{definition (standard $\Lambda$CDM)}{Eq.~\eqref{eq:l2_friedmann}; Chapter~\ref{ch:level2}, section ``Model and implementation''}
\iamfsentry{191}{H_m^2(a)=H^2(a)+\beta_mE(a)H_0^2,\qquad E(a)=\exp\!\left(1-\frac1a\right),}{Matter-sector rate used by the CDM and baryon perturbations in the modified CAMB, with $\beta_m=0.15765$.}{\derived{} (given the completed first law, Eq.~\eqref{eq:firstlaw}, a \conjecture, and the sector split, a \conjecture)}{Eq.~\eqref{eq:l2_Hm}; Chapter~\ref{ch:level2}, section ``Model and implementation''}
\iamfsentry{192}{\beta_m=\frac{\Omega_m}{2}=\frac{0.3153}{2}=0.15765,}{}{\measured{}}{Eq.~\eqref{eq:l2_beta}; Chapter~\ref{ch:level2}, section ``Model and implementation''}
\iamfsentry{193}{\frac{4\pi G\bar\rho_m}{H_m^2}\,\delta_m=\frac32\,\Omega_m(a)\,\frac{H^2}{H_m^2}\,\delta_m\equiv\frac32\,\Omega_m(a)\,\mu(a)\,\delta_m ,}{}{\derived{}}{Chapter~\ref{ch:level2}, section ``Model and implementation''}
\iamfsentry{194}{\begin{aligned}\mu(a)&=\frac{H^2_{\Lambda\rm CDM}(a)}{H^2_{\Lambda\rm CDM}(a)+\beta_mE(a)H_0^2},\\
\Sigma(a)&=1\quad\text{(exact)}.\end{aligned}}{Growth coupling if the change acted as $G_{\rm eff}=\mu G$; present value $1/(1+\beta_m)$.}{\derived{} (for a Newton constant scaled by $H^2/H_m^2$; how the code realizes it is given below)}{Eqs.~\eqref{eq:l2_mu}, \eqref{eq:l2_sigma}; Chapter~\ref{ch:level2}, section ``Model and implementation''}
\iamfsentry{195}{\mu(z{=}0)=\frac{1}{1+\beta_m}=\frac{1}{1.15765}=0.864,}{}{\calc{}}{Eq.~\eqref{eq:l2_mu0}; Chapter~\ref{ch:level2}, section ``Model and implementation''}
\iamfsentry{196}{\frac{H_m(a)}{H_\gamma(a)}=\sqrt{1+\frac{\beta_mE(a)H_0^2}{H^2(a)}},}{}{\derived{}}{Eq.~\eqref{eq:l2_ratio}; Chapter~\ref{ch:level2}, section ``The sector-dependent Hubble parameter''}
\iamfsentry{197}{\frac{H_m(z{=}0)}{H_\gamma(z{=}0)}=\sqrt{1+\beta_m}=\sqrt{1.15765}=1.0759 .}{}{\derived{}}{Eq.~\eqref{eq:l2_ratio0}; Chapter~\ref{ch:level2}, section ``The sector-dependent Hubble parameter''}
\iamfsentry{198}{H_0^{(m)}=H_0\sqrt{1+\beta_m}=67.161\times1.0759=72.26\pm0.50\ {\rm km\,s^{-1}\,Mpc^{-1}},}{}{\calc{}}{Eq.~\eqref{eq:l2_H0m}; Chapter~\ref{ch:level2}, section ``The sector-dependent Hubble parameter''}
\iamfsentry{199}{\begin{aligned}\frac{H_0^{(\gamma)}-H_0^{\rm Planck}}{\sigma_{\rm Planck}}&=\frac{67.16-67.36}{0.54}=-0.37\sigma,\\
\frac{H_0^{(m)}-H_0^{\rm SH0ES}}{\sigma_{\rm SH0ES}}&=\frac{72.26-73.04}{1.04}=-0.75\sigma\end{aligned}}{}{\calc{}}{Eqs.~\eqref{eq:l2_tension_planck}, \eqref{eq:l2_tension_shoes}; Chapter~\ref{ch:level2}, section ``The sector-dependent Hubble parameter''}
\iamfsentry{200}{H^2(a)=H_0^2\left[\Omega_ma^{-3}+\Omega_ra^{-4}+\Omega_\Lambda+\beta_mE(a)\right].}{}{\measured{}}{Eq.~\eqref{eq:l2_bg}; Chapter~\ref{ch:level2}, section ``The background test''}
\subsection*{Chapter~\ref{ch:dsnote}: Why matter and light are in different sectors}
\iamfsentry{201}{\begin{aligned}\text{Massive particles (matter):}\quad & g_{\mu\nu}\frac{dx^\mu}{d\tau}\frac{dx^\nu}{d\tau}=-1,\qquad d\tau>0, \\
\text{Massless particles (photons):}\quad & g_{\mu\nu}\frac{dx^\mu}{d\lambda}\frac{dx^\nu}{d\lambda}=0,\qquad d\tau=0\end{aligned}}{}{\observed{}}{Eqs.~\eqref{eq:dsn_timelike}, \eqref{eq:dsn_null}; Chapter~\ref{ch:dsnote}, section ``The boundary Einstein built into the theory''}
\iamfsentry{202}{-dE=T_H\,d\!\left(S_{\rm geo}+S_{\rm info}\right),}{Horizon first law with the record entropy added (as Eq.~\eqref{eq:firstlaw}).}{\conjecture{} (Assumption~2)}{Eq.~\eqref{eq:dsn_firstlaw}; Chapter~\ref{ch:dsnote}, section ``The sector split in IAM: physics, not assumption''; see also (\ref{app:fs:61}), (\ref{app:fs:125})}
\iamfsentry{203}{S_{\rm info}>0\quad\text{if and only if}\quad d\tau>0 .}{}{\derived{}}{Eq.~\eqref{eq:dsn_iff}; Chapter~\ref{ch:dsnote}, section ``The sector split in IAM: physics, not assumption''}
\iamfsentry{204}{\begin{aligned}\mu(a)&=\frac{H^2_{\Lambda\rm CDM}(a)}{H^2_{\Lambda\rm CDM}(a)+\beta_mE(a)H_0^2}<1\qquad\text{(matter sector: suppressed growth)}, \\
\Sigma(a)&=1\qquad\text{(photon sector: standard GR, exact)},\end{aligned}}{}{\derived{}}{Eqs.~\eqref{eq:dsn_mu}, \eqref{eq:dsn_sigma}; Chapter~\ref{ch:dsnote}, section ``The sector split in IAM: physics, not assumption''}
\subsection*{Chapter~\ref{ch:s8trend}: Growth across redshift: the $S_8$ trend}
\iamfsentry{205}{\mu(a)=\frac{H^2_{\Lambda\rm CDM}(a)}{H^2_{\Lambda\rm CDM}(a)+\beta_mE(a)H_0^2},\qquad H^2_{\Lambda\rm CDM}(a)=H_0^2\left(\Omega_ma^{-3}+\Omega_\Lambda\right),}{}{\derived{}}{Eq.~\eqref{eq:s8_mu}; Chapter~\ref{ch:s8trend}, section ``The coupling''}
\iamfsentry{206}{E(a)=\exp\!\left(1-\frac1a\right)}{}{\derived{}}{Eq.~\eqref{eq:s8_Ea}; Chapter~\ref{ch:s8trend}, section ``The coupling''; see also (\ref{app:fs:185}), (\ref{app:fs:318}), (\ref{app:fs:332}), (\ref{app:fs:408}), (\ref{app:fs:520})}
\iamfsentry{207}{\beta_m=\frac{\Omega_m}{2}=0.15765\qquad(\Omega_m=0.3153),}{}{\prediction{}}{Eq.~\eqref{eq:s8_beta}; Chapter~\ref{ch:s8trend}, section ``The coupling''}
\iamfsentry{208}{\mu_0\equiv\mu(a{=}1)-1=-\frac{\beta_m}{\Omega_m+\Omega_\Lambda+\beta_m}=-\frac{\beta_m}{1+\beta_m}=-0.1362 .}{}{\derived{}}{Eq.~\eqref{eq:s8_mu0}; Chapter~\ref{ch:s8trend}, section ``The coupling''}
\subsection*{Chapter~\ref{ch:sectortension}: Dark energy or two rulers?}
\iamfsentry{209}{S_{\rm total}=S_{\rm geometric}+S_{\rm informational},\qquad S_{\rm geometric}=\frac{k_BA}{4\ell_P^2},}{}{\interp{}}{Eq.~\eqref{eq:st_entropy}; Chapter~\ref{ch:sectortension}, section ``The informational term in brief''}
\iamfsentry{210}{\tilde H_m^2(a)=H^2_{\Lambda\rm CDM}(a)+\beta_mE(a)H_0^2,\qquad H^2_{\Lambda\rm CDM}(a)=H_0^2\left(\Omega_ma^{-3}+\Omega_ra^{-4}+\Omega_\Lambda\right),}{}{\interp{}}{Eq.~\eqref{eq:st_hm}; Chapter~\ref{ch:sectortension}, section ``The informational term in brief''}
\iamfsentry{211}{\mu(a)=\frac{H^2_{\Lambda\rm CDM}(a)}{H^2_{\Lambda\rm CDM}(a)+\beta_mE(a)H_0^2}<1,\qquad \Sigma(a)=1\ \text{(exact)},}{}{\interp{}}{Eq.~\eqref{eq:st_mu}; Chapter~\ref{ch:sectortension}, section ``The informational term in brief''}
\iamfsentry{212}{\frac{d^2D}{dx^2}+\left(2+\frac{d\ln H}{dx}\right)\frac{dD}{dx}-\frac32\,\mu(a)\,\Omega_m(a)\,D=0,}{}{\derived{}}{Eq.~\eqref{eq:st_ode}; Chapter~\ref{ch:sectortension}, section ``Growth at the DESI redshifts''}
\iamfsentry{213}{\sigma_8^{\rm IAM}=0.7998\pm0.0058\quad\text{against}\quad\sigma_8^{\Lambda\rm CDM}=0.8087\pm0.0059,}{}{\calc{}}{Eq.~\eqref{eq:st_sigma8}; Chapter~\ref{ch:sectortension}, section ``The clustering amplitude''}
\iamfsentry{214}{S_8^{\rm IAM}=0.822\pm0.011\quad\text{against}\quad S_8^{\Lambda\rm CDM}=0.830\pm0.011,}{}{\calc{}}{Eq.~\eqref{eq:st_S8}; Chapter~\ref{ch:sectortension}, section ``The clustering amplitude''}
\subsection*{Chapter~\ref{ch:phantom}: The phantom crossing and the two rulers}
\iamfsentry{215}{a_{\rm cross}=1+\frac{1+w_0}{w_a},\qquad z_{\rm cross}=\frac1{a_{\rm cross}}-1 .}{}{\derived{} (checked with sympy)}{Eq.~\eqref{eq:ph_zcross}; Chapter~\ref{ch:phantom}, section ``The crossing redshift''}
\subsection*{Chapter~\ref{ch:dsvalidation}: Type Ia supernovae in the dual-sector picture: $\Lambda$CDM distances with a locally calibrated $H_0$}
\iamfsentry{216}{H_m^2(a)=H_0^2\left[\Omega_ma^{-3}+\Omega_\Lambda+\beta E(a)\right],}{Matter-sector expansion rate in density-parameter form, used against the supernova Hubble diagram.}{\derived{} (given the completed first law, Eq.~\eqref{eq:firstlaw}, a \conjecture, and the sector split, a \conjecture)}{Eq.~\eqref{eq:dsv_Hm}; Chapter~\ref{ch:dsvalidation}, section ``Theoretical framework''}
\iamfsentry{217}{\begin{aligned}&\text{photon sector (CMB, photon propagation):}\\&\qquad\beta_\gamma<0.0052\quad(95\,\%,\ \text{acoustic scale}), \\
&\text{matter sector (local $H_0$, structure growth):}\\&\qquad\beta_m=\Omega_m/2=0.15765\quad(\text{zero free parameters}).\end{aligned}}{}{\calc{}}{Eqs.~\eqref{eq:dsv_bg}, \eqref{eq:dsv_bm}; Chapter~\ref{ch:dsvalidation}, section ``Theoretical framework''}
\iamfsentry{218}{\begin{aligned}H_0(\text{photon})&=67.16\pm0.47\ \text{km\,s}^{-1}\,\text{Mpc}^{-1}\quad\text{(Level 2 posterior)}, \\
H_0(\text{matter})&=H_0\sqrt{1+\beta_m}=67.161\times1.0759=72.26\pm0.50\ \text{km\,s}^{-1}\,\text{Mpc}^{-1}.\end{aligned}}{}{\measured{}}{Eqs.~\eqref{eq:dsv_H0g}, \eqref{eq:dsv_H0m}; Chapter~\ref{ch:dsvalidation}, section ``Theoretical framework''}
\iamfsentry{219}{\mu_{\rm obs}=m_b^{\rm corr}-M,}{}{\observed{}}{Eq.~\eqref{eq:dsv_mbcorr}; Chapter~\ref{ch:dsvalidation}, section ``Data and methodology''}
\iamfsentry{220}{m_b^{\rm model}(z)=M+5\log_{10}\!\left[\frac{d_L(z)}{10\ {\rm pc}}\right]=M+5\log_{10}\!\left[\frac{d_L(z)}{1\ {\rm Mpc}}\right]+25,}{}{\derived{}}{Eq.~\eqref{eq:dsv_mu}; Chapter~\ref{ch:dsvalidation}, section ``Data and methodology''}
\iamfsentry{221}{d_L(z)=(1+z)\int_0^z\frac{c\,dz'}{H(z';\beta)}}{}{\derived{}}{Eq.~\eqref{eq:dsv_dL}; Chapter~\ref{ch:dsvalidation}, section ``Data and methodology''}
\iamfsentry{222}{H(z;\beta)=H_0\sqrt{\Omega_m(1+z)^3+\Omega_\Lambda+\beta E\!\left(\frac{1}{1+z}\right)}.}{}{\derived{}}{Eq.~\eqref{eq:dsv_Hz}; Chapter~\ref{ch:dsvalidation}, section ``Data and methodology''}
\iamfsentry{223}{\chi^2=\sum_{i=1}^{1588}\left(\frac{m_{b,i}^{\rm corr}-m_b^{\rm model}(z_i;\Omega_m,H_0,\beta,M)}{\sigma_{m_b,i}}\right)^2+\chi^2_{\rm prior},}{}{\derived{}}{Eq.~\eqref{eq:dsv_chi2}; Chapter~\ref{ch:dsvalidation}, section ``Data and methodology''}
\iamfsentry{224}{\begin{gathered}
\text{Test A: }\chi^2_{\rm prior}=\left(\frac{H_0-67.4}{0.5}\right)^2\ \text{(Planck)},\qquad \text{Test B: }\chi^2_{\rm prior}=\left(\frac{H_0-73.04}{1.04}\right)^2\ \text{(SH0ES)},\\
\text{Test C: }\chi^2_{\rm prior}=0\ \text{(unconstrained)}.
\end{gathered}}{}{\derived{}}{Eq.~\eqref{eq:dsv_priors}; Chapter~\ref{ch:dsvalidation}, section ``Data and methodology''}
\iamfsentry{225}{m_b^{\rm model}=\left(M-5\log_{10}H_0\right)+5\log_{10}\!\left[c\,(1+z)\,I(z;\Omega_m,\beta)\right]+25 .}{}{\derived{}}{Chapter~\ref{ch:dsvalidation}, section ``Data and methodology''}
\iamfsentry{226}{\hat{\mathcal M}=\frac{\sum_iw_ir_i}{\sum_iw_i},\qquad \chi^2_{\min}(\Omega_m,\beta)=\sum_iw_i\left(r_i-\hat{\mathcal M}\right)^2 ;}{}{\derived{}}{Chapter~\ref{ch:dsvalidation}, section ``Data and methodology''}
\iamfsentry{227}{H_0(\text{local})=H_0(\text{base})\times\sqrt{1+\beta_m}=72.26\ \text{km\,s}^{-1}\,\text{Mpc}^{-1}.}{}{\prediction{}}{Eq.~\eqref{eq:dsv_H0local}; Chapter~\ref{ch:dsvalidation}, section ``Physical interpretation''}
\iamfsentry{228}{d_L(z)\propto(1+z)\int_0^z\frac{dz'}{H(z')} .}{}{\calc{}}{Eq.~\eqref{eq:dsv_shape}; Chapter~\ref{ch:dsvalidation}, section ``Physical interpretation''}
\iamfsentry{229}{\Omega_m(a;\beta)=\frac{\Omega_ma^{-3}}{\Omega_ma^{-3}+\Omega_\Lambda+\beta E(a)}<\Omega_m(a;0).}{Matter density fraction with the record term in the denominator; lower than in $\Lambda$CDM, which lowers growth.}{\derived}{Eq.~\eqref{eq:dsv_Omdil}; Chapter~\ref{ch:dsvalidation}, section ``Physical interpretation''}
\iamfsentry{230}{\frac{\Omega_m(a;\beta)}{\Omega_m(a;0)}=\frac{H^2_{\Lambda\rm CDM}(a)}{H^2_{\Lambda\rm CDM}(a)+\beta E(a)H_0^2}=\mu(a),}{}{\derived{}}{Eq.~\eqref{eq:dsv_Omdil_mu}; Chapter~\ref{ch:dsvalidation}, section ``Physical interpretation''}
\iamfsentry{231}{k^2\Psi=-4\pi G\,\mu(a)\,\rho_m\delta_m a^2,\qquad k^2(\Phi+\Psi)=-8\pi G\,\Sigma(a)\,\rho_m\delta_m a^2.}{}{\derived{}}{Eq.~\eqref{eq:dsv_poisson}; Chapter~\ref{ch:dsvalidation}, section ``Discussion''}
\iamfsentry{232}{\mu(a)=\frac{H^2_{\Lambda\rm CDM}(a)}{H^2_{\Lambda\rm CDM}(a)+\beta_mE(a)H_0^2},\qquad\Sigma(a)=1,}{}{\derived{}}{Eq.~\eqref{eq:dsv_mu_sigma}; Chapter~\ref{ch:dsvalidation}, section ``Discussion''}
\subsection*{Chapter~\ref{ch:darkenergy}: The equation of state of the record}
\iamfsentry{233}{S_{\rm info}(a)=S_0\,E(a)=S_0\exp\!\left(1-\frac1a\right),}{}{\prediction{}}{Eq.~\eqref{eq:de_sinfo}; Chapter~\ref{ch:darkenergy}, section ``The IAM equation of state''}
\iamfsentry{234}{H_m^2(a)=H^2(a)+\beta_mE(a)H_0^2,\qquad E(a)=e^{1-1/a},\qquad\beta_m=\Omega_m/2,}{Matter-sector expansion rate; the photon sector keeps the $\Lambda$CDM rate.}{\derived{} (given the completed first law, Eq.~\eqref{eq:firstlaw}, a \conjecture, and the sector split, a \conjecture)}{Eq.~\eqref{part2:eq:Hm}; Chapter~\ref{ch:darkenergy}, section ``The IAM equation of state''}
\iamfsentry{235}{\rho_{\rm info}(a)=\beta_m\,E(a)\,\rho_{\rm crit,0},\qquad\rho_{\rm crit,0}=\frac{3H_0^2}{8\pi G}.}{}{\derived{}}{Eq.~\eqref{eq:de_rhoinfo}; Chapter~\ref{ch:darkenergy}, section ``The IAM equation of state''}
\iamfsentry{236}{\dot\rho_{\rm info}+3H\,(1+w_{\rm info})\,\rho_{\rm info}=0
\quad\Longrightarrow\quad w_{\rm info}=-1-\frac13\,\frac{d\ln\rho_{\rm info}}{d\ln a}.}{}{\derived{}}{Eq.~\eqref{eq:de_cont}; Chapter~\ref{ch:darkenergy}, section ``The IAM equation of state''}
\iamfsentry{237}{\frac{d\ln E}{d\ln a}=a\,\frac{d}{da}\left(1-\frac1a\right)=a\cdot\frac{1}{a^2}=\frac1a .}{}{\derived{}}{Eq.~\eqref{eq:de_dlnE}; Chapter~\ref{ch:darkenergy}, section ``The IAM equation of state''}
\iamfsentry{238}{\boxed{w_{\rm info}(a)=-1-\frac{1}{3a}.}}{}{\derived{}}{Eq.~\eqref{eq:de_winfo}; Chapter~\ref{ch:darkenergy}, section ``The IAM equation of state''}
\iamfsentry{239}{\lim_{a\to\infty}w_{\rm info}(a)=-1 .}{}{\derived{}}{Eq.~\eqref{eq:de_wlim}; Chapter~\ref{ch:darkenergy}, section ``The IAM equation of state''}
\iamfsentry{240}{w_{\rm info}(1)=-1-\tfrac13=-\tfrac43 .}{}{\derived{}}{Eq.~\eqref{eq:de_w1}; Chapter~\ref{ch:darkenergy}, section ``The IAM equation of state''}
\iamfsentry{241}{w_0=w_{\rm info}(a)\big|_{a=1}=-\tfrac43\approx-1.333,\qquad w_a=-\left.\frac{dw_{\rm info}(a)}{da}\right|_{a=1}=-\tfrac13\approx-0.333 .}{CPL image of $w_{\rm info}$ at $a=1$: $w_0=-4/3$, $w_a=-1/3$.}{\derived{}}{Eq.~\eqref{eq:wz_cpl}; Chapter~\ref{ch:darkenergy}, section ``The IAM equation of state''}
\iamfsentry{242}{E(0^+)=0,\qquad E(1)=1\ \ (\text{today}),\qquad\lim_{a\to\infty}E(a)=e\approx2.71828\ldots}{}{\derived{}}{Eq.~\eqref{eq:de_Elimits}; Chapter~\ref{ch:darkenergy}, section ``$E(a)$ saturation and the asymptotic universe''}
\iamfsentry{243}{H^2(a)=H_0^2\left[\Omega_ma^{-3}+\Omega_\Lambda\right].}{}{\calc{}}{Eq.~\eqref{eq:de_friedmann}; Chapter~\ref{ch:darkenergy}, section ``$E(a)$ saturation and the asymptotic universe''}
\iamfsentry{244}{H_\infty=H_0\sqrt{\Omega_\Lambda}\approx H_0\times0.8275\approx55.6~\mathrm{km\,s^{-1}\,Mpc^{-1}},}{}{\calc{}}{Eq.~\eqref{eq:de_Hinf}; Chapter~\ref{ch:darkenergy}, section ``$E(a)$ saturation and the asymptotic universe''}
\iamfsentry{245}{\lim_{a\to\infty}w_{\rm eff}(a)=-1 .}{}{\derived{}}{Eq.~\eqref{eq:de_weff}; Chapter~\ref{ch:darkenergy}, section ``$E(a)$ saturation and the asymptotic universe''}
\iamfsentry{246}{\frac{E(1)}{e}=\frac{e^0}{e}=\frac1e\approx0.36788=36.8\,\%.}{}{\derived{}}{Eq.~\eqref{eq:de_maturity_today}; Chapter~\ref{ch:darkenergy}, section ``Informational maturity of the universe''}
\iamfsentry{247}{\frac{d(E/e)}{da}=\frac{E(a)}{e\,a^2}=\frac{e^{-1/a}}{a^2},}{}{\derived{}}{Eq.~\eqref{eq:de_rate_a}; Chapter~\ref{ch:darkenergy}, section ``Informational maturity of the universe''}
\iamfsentry{248}{\frac{d(E/e)}{d\ln a}=\frac{e^{-1/a}}{a},}{}{\derived{}}{Eq.~\eqref{eq:de_rate_lna}; Chapter~\ref{ch:darkenergy}, section ``Informational maturity of the universe''}
\iamfsentry{249}{a(f)=\frac{-1}{\ln f}.}{}{\derived{}}{Eq.~\eqref{eq:de_af}; Chapter~\ref{ch:darkenergy}, section ``Informational maturity of the universe''}
\iamfsentry{250}{\left.\frac{d(E/e)}{dt}\right|_{a=1}=\left.\frac{d(E/e)}{da}\cdot\frac{da}{dt}\right|_{a=1}=\frac1e\cdot H_0\approx2.54\,\%\ \mathrm{Gyr^{-1}}}{}{\calc{}}{Eq.~\eqref{eq:de_rate_t}; Chapter~\ref{ch:darkenergy}, section ``Informational maturity of the universe''}
\subsection*{Chapter~\ref{ch:wzfuture}: The equation of state through time, and the far future}
\iamfsentry{251}{\frac{dw_{\rm info}}{da}=\frac{1}{3a^2}>0\quad\text{for all }a>0,\qquad\rho_{\rm info}(a)\le e\,\rho_{\rm info}(1).}{No-Big-Rip condition: $w_{\rm info}$ rises with $a$ and $\rho_{\rm info}$ is bounded by $e\,\rho_{\rm info}(1)$.}{\derived{} (given $E(a)\le e$, from the record constraint, a \conjecture)}{Eq.~\eqref{eq:wz_norip}; Chapter~\ref{ch:wzfuture}, section ``The no-Big-Rip proof''}
\iamfsentry{252}{\frac{\rho_{\rm info}(a)}{\rho_{\rm info}(1)}=\exp\!\left[-3\int_1^a\bigl(1+w_{\rm info}\bigr)\,\frac{da'}{a'}\right]
=\exp\!\left[\int_1^a\frac{da'}{a'^2}\right]=e^{\,1-1/a}=E(a),}{}{\derived{}}{Eq.~\eqref{eq:wz_rho_from_w}; Chapter~\ref{ch:wzfuture}, section ``The no-Big-Rip proof''}
\iamfsentry{253}{\begin{aligned}w_0&=-0.838\pm0.055,& w_a&=-0.62^{+0.22}_{-0.19} &&\text{(DR2 + CMB + Pantheon+)},\\
w_0&=-0.667\pm0.088,& w_a&=-1.09^{+0.31}_{-0.27} &&\text{(DR2 + CMB + Union3)},\\
w_0&=-0.752\pm0.057,& w_a&=-0.86^{+0.23}_{-0.20} &&\text{(DR2 + CMB + DES Y5)}.\end{aligned}}{}{\observed{}}{Eqs.~\eqref{eq:wz_desi1}, \eqref{eq:wz_desi2}, \eqref{eq:wz_desi3}; Chapter~\ref{ch:wzfuture}, section ``Comparison to DESI DR2''}
\iamfsentry{254}{(w_0,w_a)_{\rm IAM}=\left(-\tfrac43,\,-\tfrac13\right)\approx(-1.333,\,-0.333).}{}{\calc{}}{Eq.~\eqref{eq:wz_cplpoint}; Chapter~\ref{ch:wzfuture}, section ``Comparison to DESI DR2''}
\subsection*{Chapter~\ref{ch:lambda}: The cosmological constant}
\iamfsentry{255}{\rho_{\rm vac}\sim\frac{E_P^4}{(\hbar c)^3}=\frac{c^7}{\hbar G^2}=4.633\times10^{113}\ {\rm J\,m^{-3}},}{}{\calc{}}{Eq.~\eqref{eq:lam_rhovac}; Chapter~\ref{ch:lambda}, section ``The problem''}
\iamfsentry{256}{\rho_\Lambda=\Omega_\Lambda\rho_cc^2=5.250\times10^{-10}\ {\rm J\,m^{-3}},\qquad \rho_c=\frac{3H_0^2}{8\pi G},}{}{\observed{}}{Eq.~\eqref{eq:lam_rhoL}; Chapter~\ref{ch:lambda}, section ``The problem''}
\iamfsentry{257}{\frac{\rho_\Lambda}{\rho_{\rm vac}}=1.133\times10^{-123},\qquad \log_{10}\frac{\rho_\Lambda}{\rho_{\rm vac}}=-122.95 .}{}{\calc{}}{Eq.~\eqref{eq:lam_ratio}; Chapter~\ref{ch:lambda}, section ``The problem''}
\iamfsentry{258}{E_{\rm bit}=\frac{\hbar H_0\ln2}{2\pi}=2.541\times10^{-53}\ {\rm J}.}{}{\derived{} (the principle)}{Eq.~\eqref{eq:lam_ebit}; Chapter~\ref{ch:lambda}, section ``Two physically distinct quantities''}
\iamfsentry{259}{N_{\max}=\frac{S}{k_B}=\frac{A_H}{4l_P^2}=\frac{4\pi(c/H_0)^2}{4l_P^2}=2.265\times10^{122}\ \text{nats}=3.268\times10^{122}\ \text{bits},}{}{\derived{}}{Eq.~\eqref{eq:lam_nmax}; Chapter~\ref{ch:lambda}, section ``Two physically distinct quantities''}
\iamfsentry{260}{f_{\rm geo}=\frac{l_P^2}{A_H}=\frac{l_P^2}{4\pi(c/H_0)^2}=\frac{l_P^2H_0^2}{4\pi c^2}=\frac{1}{4\pi}\left(\frac{l_P}{l_H}\right)^2,}{}{\derived{}}{Eq.~\eqref{eq:lam_fgeo}; Chapter~\ref{ch:lambda}, section ``The calculation''}
\iamfsentry{261}{f_b=\frac{\Omega_b}{\Omega_m}.}{}{\conjecture{} (that this is the writing fraction)}{Eq.~\eqref{eq:lam_fb}; Chapter~\ref{ch:lambda}, section ``The calculation''}
\iamfsentry{262}{\frac{\rho_\Lambda}{\rho_{\rm vac}}\bigg|_{\rm IAM}=\frac{2}{\pi}\left(\frac{l_P}{l_H}\right)^2\frac{\Omega_b}{\Omega_m}.}{The cosmological-constant expression with writing fraction $\Omega_b/\Omega_m$ and geometric factor $2/\pi$; $1.218$ times the measured ratio.}{\fitted{} (the writing fraction and the $2/\pi$ were introduced as the residual was reduced; the $2/\pi$ is not derived); value \calc}{Eq.~\eqref{eq:base}; Chapter~\ref{ch:lambda}, section ``The calculation''}
\iamfsentry{263}{\begin{aligned}l_P&=\sqrt{\hbar G/c^3}=1.616\times10^{-35}\ {\rm m}, \\
H_0&=67.4\ {\rm km\,s^{-1}\,Mpc^{-1}}=2.184\times10^{-18}\ {\rm s^{-1}}, \\
l_H&=c/H_0=1.372\times10^{26}\ {\rm m}, \\
\Omega_b&=0.0493,\qquad \Omega_m=0.3153 .\end{aligned}}{}{\measured{} ($l_P$, $H_0$, $\Omega_b$, $\Omega_m$)}{Eqs.~\eqref{eq:lam_in1}, \eqref{eq:lam_in2}, \eqref{eq:lam_in3}, \eqref{eq:lam_in4}; Chapter~\ref{ch:lambda}, section ``The calculation''}
\iamfsentry{264}{\left(\frac{l_P}{l_H}\right)^2=\left(\frac{1.616\times10^{-35}}{1.372\times10^{26}}\right)^2=1.387\times10^{-122},\qquad
\frac{\Omega_b}{\Omega_m}=\frac{0.0493}{0.3153}=0.1564 .}{}{\calc{}}{Eq.~\eqref{eq:lam_geo}; Chapter~\ref{ch:lambda}, section ``The calculation''}
\iamfsentry{265}{\frac{\rho_\Lambda}{\rho_{\rm vac}}\bigg|_{\rm IAM}=\frac{2}{\pi}\times1.387\times10^{-122}\times0.1564=1.380\times10^{-123},}{}{\calc{}}{Eq.~\eqref{eq:lam_result}; Chapter~\ref{ch:lambda}, section ``The calculation''}
\iamfsentry{266}{\frac{\rho_\Lambda}{\rho_{\rm vac}}\bigg|_{\rm obs}=\frac{5.250\times10^{-10}}{4.633\times10^{113}}=1.133\times10^{-123},}{}{\calc{}}{Eq.~\eqref{eq:lam_obs}; Chapter~\ref{ch:lambda}, section ``The calculation''}
\iamfsentry{267}{\frac{(\rho_\Lambda/\rho_{\rm vac})_{\rm IAM}}{(\rho_\Lambda/\rho_{\rm vac})_{\rm obs}}=1.218 .}{}{\calc{}}{Eq.~\eqref{eq:lam_factor}; Chapter~\ref{ch:lambda}, section ``The calculation''}
\iamfsentry{268}{T_H=\frac{\hbar H_0}{2\pi k_B},}{}{\derived{}}{Eq.~\eqref{eq:lam_TH}; Chapter~\ref{ch:lambda}, section ``The calculation''}
\iamfsentry{269}{T_{dS}=\frac{\hbar H_{dS}}{2\pi k_B}=T_H\sqrt{\Omega_\Lambda},}{}{\derived{} (the temperature of the horizon set by $\Lambda$ alone)}{Eq.~\eqref{eq:lam_TdS}; Chapter~\ref{ch:lambda}, section ``The calculation''}
\iamfsentry{270}{E_{\rm bit,eff}=k_BT_{dS}\ln2=k_BT_H\ln2\times\sqrt{\Omega_\Lambda}.}{}{\calc{}}{Eq.~\eqref{eq:lam_Eeff}; Chapter~\ref{ch:lambda}, section ``The calculation''}
\iamfsentry{271}{\frac{\rho_\Lambda}{\rho_{\rm vac}}\bigg|_{T_{dS}}=\frac{2}{\pi}\left(\frac{l_P}{l_H}\right)^2\sqrt{\Omega_\Lambda}\,\frac{\Omega_b}{\Omega_m},}{The same expression with the de Sitter temperature factor $\sqrt{\Omega_\Lambda}$; $0.79\,\%$ above the measured value.}{\fitted{} (the $\sqrt{\Omega_\Lambda}$ was applied to close the residual of Eq.~\eqref{eq:base}); value \calc}{Eq.~\eqref{eq:corr}; Chapter~\ref{ch:lambda}, section ``The calculation''}
\iamfsentry{272}{\sqrt{\Omega_\Lambda}=0.8274,\qquad \frac{\rho_\Lambda}{\rho_{\rm vac}}\bigg|_{T_{dS}}=1.380\times10^{-123}\times0.8274=1.142\times10^{-123},}{}{\calc{}}{Eq.~\eqref{eq:lam_corr_num}; Chapter~\ref{ch:lambda}, section ``The calculation''}
\iamfsentry{273}{\frac{\rho_\Lambda}{\rho_{\rm vac}}=\frac{\Omega_\Lambda\,3H_0^2c^2/8\pi G}{c^7/\hbar G^2}
=\frac{3\Omega_\Lambda}{8\pi}\,\frac{\hbar GH_0^2}{c^5}=\frac{3\Omega_\Lambda}{8\pi}\left(\frac{l_P}{l_H}\right)^2,\qquad l_H=c/H_0,}{Horizon identity: the ratio rewritten exactly in Planck units; the 123 orders are $(l_P/l_H)^2$.}{\derived{} (an identity)}{Eq.~\eqref{eq:ident}; Chapter~\ref{ch:lambda}, section ``The horizon identity''}
\iamfsentry{274}{T_{GH}\,S=\frac{\hbar H}{2\pi}\,\frac{\pi c^5}{\hbar GH^2}=\frac{c^5}{2GH}=M_Hc^2\qquad\text{(any }H),}{}{\derived{}}{Eq.~\eqref{eq:lam_smarr}; Chapter~\ref{ch:lambda}, section ``The horizon identity''}
\iamfsentry{275}{\frac{2}{\pi}\sqrt{\Omega_\Lambda}\,\frac{\Omega_b}{\Omega_m}=\frac{3\Omega_\Lambda}{8\pi}\quad\Longrightarrow\quad
\boxed{\;\frac{\Omega_b}{\Omega_m}=\frac{3}{16}\sqrt{\Omega_\Lambda}\;}}{What Eq.~\eqref{eq:corr} asserts once the identity is divided out: a present-epoch relation between baryon fraction and dark-energy fraction.}{\derived{} (algebra)}{Eq.~\eqref{eq:rel}; Chapter~\ref{ch:lambda}, section ``What the expression asserts''}
\iamfsentry{276}{A_{\rm eff}=\tfrac12\times4\pi l_H^2=2\pi l_H^2 .}{}{\derived{}}{Eq.~\eqref{eq:lam_aeff}; Chapter~\ref{ch:lambda}, section ``The coefficient $2/\pi$ and the de Sitter causal structure''}
\iamfsentry{277}{N_{\rm eff}=\frac{A_{\rm eff}}{4l_P^2}=\frac{2\pi l_H^2}{4l_P^2}=\frac{\pi l_H^2}{2l_P^2}\ \text{nats},}{}{\derived{} (arithmetic)}{Eq.~\eqref{eq:lam_nbits}; Chapter~\ref{ch:lambda}, section ``The coefficient $2/\pi$ and the de Sitter causal structure''}
\iamfsentry{278}{f_{\rm bit}=\frac{l_P^2}{A_{\rm eff}/(4\pi)}=\frac{4\pi l_P^2}{2\pi l_H^2}=2\left(\frac{l_P}{l_H}\right)^2 .}{}{\derived{} (arithmetic)}{Eq.~\eqref{eq:lam_fbit}; Chapter~\ref{ch:lambda}, section ``The coefficient $2/\pi$ and the de Sitter causal structure''}
\subsection*{Chapter~\ref{ch:lambda_history}: The cosmological constant over cosmic history}
\iamfsentry{279}{\frac{dE}{da}=\frac{E(a)}{a^2},\qquad \frac{dE}{da}\Big|_{a=1}=1 .}{}{\derived{}}{Eq.~\eqref{eq:lh_dE}; Chapter~\ref{ch:lambda_history}, section ``The cosmological constant is not constant''}
\iamfsentry{280}{w=-1-\frac13\frac{d\ln\rho}{d\ln a}.}{}{\derived{}}{Eq.~\eqref{eq:lh_w}; Chapter~\ref{ch:lambda_history}, section ``The cosmological constant is not constant''}
\iamfsentry{281}{w=-1-\epsilon(a),\qquad \epsilon(a)=\frac13\frac{d\ln\rho}{d\ln a}>0,}{}{\derived{}}{Eq.~\eqref{eq:lh_eps}; Chapter~\ref{ch:lambda_history}, section ``The cosmological constant is not constant''}
\iamfsentry{282}{w_{\rm info}(a)=-1-\frac{1}{3a},\qquad w_{\rm info}(1)=-\frac43,}{}{\derived{}}{Eq.~\eqref{eq:lh_winfo}; Chapter~\ref{ch:lambda_history}, section ``The cosmological constant is not constant''}
\iamfsentry{283}{\frac{\Lambda(a)}{\rho_{\rm vac}}\propto\int_{a_{\rm EW}}^{a}\frac{\Omega_b(a')/\Omega_{\rm tot}(a')}{l_H(a')^2/l_P^2}\,
\frac{da'}{a'^2H(a')/H_0},}{The accumulation integral for $\Lambda(a)$ over cosmic history; as written it overshoots by $\sim10^{30}$.}{\conjecture; fails as written (\openprob)}{Eq.~\eqref{eq:lh_integral}; Chapter~\ref{ch:lambda_history}, section ``The full decoherence history integral''}
\iamfsentry{284}{K_{\rm integral}=3.1\times10^{30}\qquad\text{against the required}\qquad K=\frac{3\Omega_\Lambda/8\pi}{\Omega_b/\Omega_m}=0.523 .}{}{\calc{}}{Eq.~\eqref{eq:lh_K}; Chapter~\ref{ch:lambda_history}, section ``The full decoherence history integral''}
\iamfsentry{285}{\begin{gathered}\int_0^1\!\Big(\frac{H_0}{H}\Big)^2dE=0.470,\quad \int_0^1\!\frac{H_0}{H}\,dE=0.641,\quad \int_0^1\!dE=1,\\
\int_0^1\!\frac{H}{H_0}\,dE=1.983,\quad \int_0^1\!\Big(\frac{H}{H_0}\Big)^2dE=5.797 .\end{gathered}}{}{\calc{}}{Eq.~\eqref{eq:lh_weights}; Chapter~\ref{ch:lambda_history}, section ``The full decoherence history integral''}
\iamfsentry{286}{\frac{\rho_\Lambda}{\rho_{\rm vac}}=\frac{2}{\pi}\left(\frac{l_P}{l_H}\right)^2\frac{\Omega_b}{\Omega_m}=1.380\times10^{-123}}{}{\calc{}}{Eq.~\eqref{eq:lh_base}; Chapter~\ref{ch:lambda_history}, section ``Conclusions''}
\iamfsentry{287}{\frac{\rho_\Lambda}{\rho_{\rm vac}}=\frac{2}{\pi}\left(\frac{l_P}{l_H}\right)^2\sqrt{\Omega_\Lambda}\,\frac{\Omega_b}{\Omega_m}=1.142\times10^{-123},}{}{\calc{}}{Eq.~\eqref{eq:lh_corr}; Chapter~\ref{ch:lambda_history}, section ``Conclusions''}
\subsection*{Chapter~\ref{ch:baryon}: The baryon density}
\iamfsentry{288}{\eta\equiv\frac{n_b-n_{\bar b}}{n_\gamma}\approx6.1\times10^{-10},}{}{\observed{}}{Eq.~\eqref{eq:bar_eta}; Chapter~\ref{ch:baryon}, section ``The size of the asymmetry''; see also (\ref{app:fs:291})}
\iamfsentry{289}{n_b=\eta\,n_\gamma .}{}{\derived{} (definition)}{Eq.~\eqref{eq:bar_nb}; Chapter~\ref{ch:baryon}, section ``The self-consistency loop''}
\iamfsentry{290}{\eta\;\longrightarrow\;n_b\;\longrightarrow\;\dot W_{\rm info}\;\longrightarrow\;\Lambda_{\rm acc}\;\longrightarrow\;
\{\rho_{\rm dm},\rho_{\rm de}\}\;\longrightarrow\;H(a)\;\longrightarrow\;\eta .}{}{\conjecture{}}{Eq.~\eqref{eq:bar_loop}; Chapter~\ref{ch:baryon}, section ``The self-consistency loop''}
\subsection*{Chapter~\ref{ch:baryon_chain}: The baryon density from the microwave background alone}
\iamfsentry{291}{\eta\equiv\frac{n_b-n_{\bar b}}{n_\gamma}\approx6.1\times10^{-10},}{}{\observed{}}{Eq.~\eqref{eq:bc_eta}; Chapter~\ref{ch:baryon_chain}, section ``Two measured numbers''; see also (\ref{app:fs:288})}
\iamfsentry{292}{\Delta E_{\rm act}=k_BT_H\ln2\quad\text{per bit}}{}{\derived{}}{Eq.~\eqref{eq:bc_law}; Chapter~\ref{ch:baryon_chain}, section ``Two measured numbers''}
\iamfsentry{293}{\frac{\rho_\Lambda}{\rho_{\rm vac}}=\frac{2}{\pi}\left(\frac{l_P}{l_H}\right)^2\frac{\Omega_b}{\Omega_m},}{}{\openprob{}}{Eq.~\eqref{eq:bc_cc}; Chapter~\ref{ch:baryon_chain}, section ``The constraint in terms of $\Omega_b$''}
\iamfsentry{294}{\frac{T_{dS}}{T_H}=\frac{H_{dS}}{H_0}=\sqrt{\Omega_\Lambda},}{}{\derived{} (the temperature of the horizon set by $\Lambda$ alone)}{Eq.~\eqref{eq:bc_tratio}; Chapter~\ref{ch:baryon_chain}, section ``The constraint in terms of $\Omega_b$''}
\iamfsentry{295}{\frac{\rho_\Lambda}{\rho_{\rm vac}}\bigg|_{T_{dS}}=\frac{2}{\pi}\left(\frac{l_P}{l_H}\right)^2\sqrt{\Omega_\Lambda}\,
\frac{\Omega_b}{\Omega_m}=1.142\times10^{-123}}{}{\fitted{}}{Eq.~\eqref{eq:bc_cccorr}; Chapter~\ref{ch:baryon_chain}, section ``The constraint in terms of $\Omega_b$''}
\iamfsentry{296}{\eta=2.739\times10^{-8}\,\Omega_bh^2 .}{}{\derived{}}{Eq.~\eqref{eq:bc_etaob}; Chapter~\ref{ch:baryon_chain}, section ``The constraint in terms of $\Omega_b$''}
\iamfsentry{297}{\begin{aligned}\text{other runs:}\quad &\Omega_bh^2\ \text{flat on }[0.020,\,0.025],\quad \text{start }\mathcal N(0.02242,\,0.00014^2),\\
\text{this chain:}\quad &\Omega_bh^2\ \text{flat on }[0.010,\,0.040],\quad \text{start }\mathcal N(0.02242,\,0.001^2).\end{aligned}}{}{\calc{}}{Eqs.~\eqref{eq:bc_std}, \eqref{eq:bc_test}; Chapter~\ref{ch:baryon_chain}, section ``The chain''}
\iamfsentry{298}{\Omega_bh^2=0.022320\pm0.000136,}{}{\measured{}}{Eq.~\eqref{eq:bc_ob}; Chapter~\ref{ch:baryon_chain}, section ``The chain''}
\iamfsentry{299}{\eta=(6.113\pm0.037)\times10^{-10}.}{}{\measured{}}{Eq.~\eqref{eq:bc_etares}; Chapter~\ref{ch:baryon_chain}, section ``The chain''}
\iamfsentry{300}{\frac{\Omega_b/\Omega_m}{(3/16)\sqrt{\Omega_\Lambda}}=1.0046\pm0.0068\quad(0.7\sigma).}{}{\measured{}}{Eq.~\eqref{eq:bc_ratio}; Chapter~\ref{ch:baryon_chain}, section ``The chain''}
\subsection*{Chapter~\ref{ch:surveys}: What the surveys will measure}
\iamfsentry{301}{k^2\Psi=-4\pi G\,a^2\,\mu(a,k)\,\bar\rho\,\Delta,\qquad k^2(\Phi+\Psi)=-8\pi G\,a^2\,\Sigma(a,k)\,\bar\rho\,\Delta .}{}{\prediction{}}{Eq.~\eqref{eq:sp_poisson}; Chapter~\ref{ch:surveys}, section ``The dual-sector signature in the $\mu$--$\Sigma$ language''}
\iamfsentry{302}{\mu(a)=\frac{H^2_{\Lambda\rm CDM}(a)}{H^2_{\Lambda\rm CDM}(a)+\beta_mE(a)H_0^2},\qquad E(a)=e^{1-1/a},\qquad\beta_m=\frac{\Omega_m}{2},\qquad\Sigma(a)=1,}{Matter coupling $\mu(a)$ with $E(a)$, $\beta_m=\Omega_m/2$ and $\Sigma=1$, as used for survey forecasts.}{\prediction{} (as labeled in the chapter's status table; $\beta_m=\Omega_m/2$ fixed)}{Eq.~\eqref{eq:sp_mu}; Chapter~\ref{ch:surveys}, section ``The dual-sector signature in the $\mu$--$\Sigma$ language''}
\iamfsentry{303}{\mu_0\equiv\mu(z{=}0)-1=-\frac{\beta_m}{1+\beta_m}=-0.136 .}{}{\prediction{}}{Eq.~\eqref{eq:sp_mu0}; Chapter~\ref{ch:surveys}, section ``The dual-sector signature in the $\mu$--$\Sigma$ language''}
\iamfsentry{304}{0.06\lesssim z\lesssim1.12 .}{}{\calc{}}{Eq.~\eqref{eq:sp_zone}; Chapter~\ref{ch:surveys}, section ``When the weakening switched on''}
\iamfsentry{305}{\left.\frac{\Delta T}{T}\right|_{\rm ISW}=\int\left(\frac{\partial\Phi}{\partial\tau}+\frac{\partial\Psi}{\partial\tau}\right)d\tau ,}{}{\derived{}}{Eq.~\eqref{eq:sp_iswT}; Chapter~\ref{ch:surveys}, section ``What light sees''}
\iamfsentry{306}{\Phi+\Psi\propto\frac{D}{a},\qquad\frac{\partial(\Phi+\Psi)}{\partial\tau}\propto aH\,\frac{d(D/a)}{d\ln a}=H\,D\,(f-1),\qquad f\equiv\frac{d\ln D}{d\ln a}.}{}{\derived{}}{Eq.~\eqref{eq:sp_iswsrc}; Chapter~\ref{ch:surveys}, section ``What light sees''}
\iamfsentry{307}{A_{\rm ISW}^{\rm IAM}/A_{\rm ISW}^{\Lambda\rm CDM}\approx1.03 .}{ISW--galaxy cross-correlation amplitude relative to $\Lambda$CDM.}{\calc{}}{Eq.~\eqref{eq:sp_isw}; Chapter~\ref{ch:surveys}, section ``What light sees''}
\iamfsentry{308}{H_0^{\rm m}=H_0^{\gamma}\sqrt{1+\beta_m}=67.16\times\sqrt{1.15765}=72.26~\mathrm{km\,s^{-1}\,Mpc^{-1}},}{Matter-sector Hubble rate read by standard sirens, $H_0^{\rm m}=H_0^\gamma\sqrt{1+\beta_m}=72.26$.}{\calc{} from the Level~2 chain value of $H_0^\gamma$; the siren value is a \prediction}{Eq.~\eqref{eq:sp_siren}; Chapter~\ref{ch:surveys}, section ``Standard sirens''}
\subsection*{Chapter~\ref{ch:lensdyn}: Lensing mass and dynamical mass}
\iamfsentry{309}{M_{\rm true}=\frac{M_{X\text{-ray}}}{1-b},}{}{\observed{}}{Eq.~\eqref{eq:ld_b}; Chapter~\ref{ch:lensdyn}, section ``The discrepancy''}
\iamfsentry{310}{\begin{aligned}\nabla^2\Psi&=4\pi Ga^2\,\mu(a)\,\bar\rho\,\delta\qquad\text{(Poisson equation; dynamics)},\\
\nabla^2(\Phi+\Psi)&=8\pi Ga^2\,\Sigma(a)\,\bar\rho\,\delta\qquad\text{(lensing potential)},\end{aligned}}{}{\derived{}}{Eqs.~\eqref{eq:ld_poisson}, \eqref{eq:ld_lens}; Chapter~\ref{ch:lensdyn}, section ``The prediction''}
\iamfsentry{311}{\mu(a)=\frac{H^2_\Lambda(a)}{H^2_\Lambda(a)+\beta_mE(a)H_0^2},\qquad\Sigma(a)=1\quad\text{(all redshifts)},}{}{\prediction{}}{Eq.~\eqref{eq:ld_mu}; Chapter~\ref{ch:lensdyn}, section ``The prediction''}
\iamfsentry{312}{M_{\rm dyn}\propto\Psi\propto\mu\,M_{\rm true}.}{}{\derived{}}{Eq.~\eqref{eq:ld_mdyn}; Chapter~\ref{ch:lensdyn}, section ``The prediction''}
\iamfsentry{313}{M_{\rm lens}\propto\frac{\Psi+\Phi}{2}\propto\Sigma\,M_{\rm true}=M_{\rm true}.}{}{\derived{}}{Eq.~\eqref{eq:ld_mlens}; Chapter~\ref{ch:lensdyn}, section ``The prediction''}
\iamfsentry{314}{\boxed{\frac{M_{\rm lens}}{M_{\rm dyn}}=\frac{\Sigma}{\mu}=\frac{1}{\mu(z)}}\,.}{}{\derived{} (in the Level~1 form, where $\mu$ stands in Eq.~\eqref{eq:ld_poisson})}{Eq.~\eqref{eq:ld_ratio_mu}; Chapter~\ref{ch:lensdyn}, section ``The prediction''}
\iamfsentry{315}{\frac{M_{\rm lens}}{M_{\rm dyn}}(a)=1+\frac{\beta_mE(a)H_0^2}{H^2_\Lambda(a)}=1+\frac{\beta_m\exp(1-1/a)}{\Omega_ma^{-3}+\Omega_\Lambda}.}{Lensing-to-dynamical mass ratio of a cluster, $1/\mu(z)$, in the Level~1 form.}{\derived{} (sympy: $1/\mu-(1+\beta_mE/H^2)=0$)}{Eq.~\eqref{eq:ld_ratio}; Chapter~\ref{ch:lensdyn}, section ``The prediction''}
\iamfsentry{316}{\frac{\Phi}{\Psi}=\frac{2-\mu}{\mu},}{}{\derived{} (sympy)}{Eq.~\eqref{eq:ld_slip}; Chapter~\ref{ch:lensdyn}, section ``The prediction''}
\subsection*{Chapter~\ref{ch:threeway}: Three cluster masses}
\iamfsentry{317}{\frac{\dot a}{a}\bigg|_{\rm matter}=a\sqrt{\frac{\rho_{\rm tot}+\beta_mE(a)\,\rho_0}{3}}\,,}{}{\derived{}}{Eq.~\eqref{eq:tw_adot}; Chapter~\ref{ch:threeway}, section ``The perturbation mechanism''}
\iamfsentry{318}{E(a)=\exp\!\left(1-\frac1a\right),}{}{\prediction{}}{Eq.~\eqref{eq:tw_E}; Chapter~\ref{ch:threeway}, section ``The perturbation mechanism''; see also (\ref{app:fs:185}), (\ref{app:fs:206}), (\ref{app:fs:332}), (\ref{app:fs:408}), (\ref{app:fs:520})}
\iamfsentry{319}{\beta_m=\frac{\Omega_m}{2}=0.15765,}{}{\prediction{}}{Eq.~\eqref{eq:tw_beta}; Chapter~\ref{ch:threeway}, section ``The perturbation mechanism''; see also (\ref{app:fs:178}), (\ref{app:fs:331}), (\ref{app:fs:399}), (\ref{app:fs:522})}
\iamfsentry{320}{\mu(a)=\frac{H^2_\Lambda(a)}{H^2_\Lambda(a)+\beta_mE(a)H_0^2},\qquad H^2_\Lambda(a)=H_0^2(\Omega_ma^{-3}+\Omega_\Lambda),}{}{\interp{}}{Eq.~\eqref{eq:tw_mu}; Chapter~\ref{ch:threeway}, section ``The perturbation mechanism''}
\iamfsentry{321}{M_{\rm hydro}=\mu\,M_{\rm true}.}{}{\derived{}}{Eq.~\eqref{eq:tw_mhydro}; Chapter~\ref{ch:threeway}, section ``The three-way ratio in the Level 1 form''}
\iamfsentry{322}{M_{\rm SZ}\approx\mu\,M_{\rm true}.}{}{\derived{} (Level~1 form)}{Eq.~\eqref{eq:tw_msz}; Chapter~\ref{ch:threeway}, section ``The three-way ratio in the Level 1 form''}
\iamfsentry{323}{M_{\rm lens}=\Sigma\,M_{\rm true}=M_{\rm true}.}{}{\derived{} (Level~1 form)}{Eq.~\eqref{eq:tw_mlens}; Chapter~\ref{ch:threeway}, section ``The three-way ratio in the Level 1 form''}
\iamfsentry{324}{R(z)\equiv\frac{M_{\rm lens}}{M_{\rm hydro}}=\frac{1}{\mu(z)},}{Lensing-to-hydrostatic mass ratio $R(z)=1/\mu(z)$ for the three-estimator cluster test.}{\calc{} (Level~1 form); which implementation form holds inside clusters is \openprob}{Eq.~\eqref{eq:tw_R}; Chapter~\ref{ch:threeway}, section ``The three-way ratio in the Level 1 form''}
\iamfsentry{325}{\frac{M_{\rm SZ}}{M_{\rm hydro}}\approx\frac{\mu}{\mu}=1,}{}{\calc{}}{Eq.~\eqref{eq:tw_szx}; Chapter~\ref{ch:threeway}, section ``The three-way ratio in the Level 1 form''}
\iamfsentry{326}{M_{\rm lens}>M_{\rm SZ}\approx M_{\rm hydro}.}{}{\calc{}}{Eq.~\eqref{eq:tw_order}; Chapter~\ref{ch:threeway}, section ``The three-way ratio in the Level 1 form''}
\iamfsentry{327}{1-b_{\rm hydro}(z)=\bigl[1-b_{\rm NT}(z)\bigr]\bigl[1-b_{\rm IAM}(z)\bigr],\qquad b_{\rm IAM}(z)=1-\mu(z),}{}{\derived{}}{Eq.~\eqref{eq:tw_bias}; Chapter~\ref{ch:threeway}, section ``The redshift slope''}
\iamfsentry{328}{\begin{aligned}M_{\rm SZ}/M_{\rm hydro}&\approx1&&\text{both matter sector; holds by calibration,}\\
M_{\rm lens}/M_{\rm hydro}&>1&&\text{lensing in the photon sector; observed,}\\
d(M_{\rm lens}/M_{\rm hydro})/dz&<0&&\text{redshift dependence: the test.}\end{aligned}}{}{\calibrated{}}{Eqs.~\eqref{eq:tw_test1}, \eqref{eq:tw_test2}, \eqref{eq:tw_test3}; Chapter~\ref{ch:threeway}, section ``The SZ-to-X-ray ratio''}
\iamfsentry{329}{R(z)=M_{\rm lens}/M_{\rm hydro}=1/\mu(z),}{}{\prediction{}}{Chapter~\ref{ch:threeway}, section ``What would test it''}
\subsection*{Chapter~\ref{ch:satellites}: Missing satellites and growth suppression from $\mu<1$}
\iamfsentry{330}{2K+V=0,\qquad K=\tfrac12|V|.}{}{\derived{}}{Eq.~\eqref{eq:ms_virial}; Chapter~\ref{ch:satellites}, section ``The virial partition and the coupling''}
\iamfsentry{331}{\beta_m=\frac{\Omega_m}{2}=0.15765,}{}{\prediction{}}{Eq.~\eqref{eq:ms_beta}; Chapter~\ref{ch:satellites}, section ``The virial partition and the coupling''; see also (\ref{app:fs:178}), (\ref{app:fs:319}), (\ref{app:fs:399}), (\ref{app:fs:522})}
\iamfsentry{332}{E(a)=\exp\!\left(1-\frac1a\right),}{}{\interp{}}{Eq.~\eqref{eq:ms_E}; Chapter~\ref{ch:satellites}, section ``Growth suppression from $\mu<1$''; see also (\ref{app:fs:185}), (\ref{app:fs:206}), (\ref{app:fs:318}), (\ref{app:fs:408}), (\ref{app:fs:520})}
\iamfsentry{333}{\mu(a)=\frac{H^2_{\Lambda\rm CDM}(a)}{H^2_{\Lambda\rm CDM}(a)+\beta_mE(a)H_0^2},}{}{\interp{}}{Eq.~\eqref{eq:ms_mu}; Chapter~\ref{ch:satellites}, section ``Growth suppression from $\mu<1$''; see also (\ref{app:fs:130})}
\iamfsentry{334}{\mu(z=0)=\frac1{1+\beta_m}=\frac1{1+\Omega_m/2}=0.864,}{}{\calc{}}{Eq.~\eqref{eq:ms_mu0}; Chapter~\ref{ch:satellites}, section ``Growth suppression from $\mu<1$''}
\iamfsentry{335}{\ddot\delta_m+2H\dot\delta_m-\tfrac32\,\Omega_m(a)H^2\,\mu(a)\,\delta_m=0,}{}{\derived{}}{Eq.~\eqref{eq:ms_growth}; Chapter~\ref{ch:satellites}, section ``Growth suppression from $\mu<1$''}
\iamfsentry{336}{\left.\frac{\Delta D}{D}\right|_{z=0}=\frac{D_{\rm IAM}(0)-D_{\Lambda\rm CDM}(0)}{D_{\Lambda\rm CDM}(0)}=-0.78\,\%}{}{\calc{}}{Eq.~\eqref{eq:ms_dD}; Chapter~\ref{ch:satellites}, section ``Growth suppression from $\mu<1$''}
\iamfsentry{337}{\frac{dn}{dM}=\sqrt{\frac2\pi}\,\frac{\bar\rho}{M^2}\,\nu\,\left|\frac{d\ln\sigma_M}{d\ln M}\right|\,e^{-\nu^2/2},\qquad\nu=\frac{\delta_c}{\sigma_M},}{}{\derived{}}{Eq.~\eqref{eq:ms_psmf}; Chapter~\ref{ch:satellites}, section ``Growth suppression from $\mu<1$''}
\iamfsentry{338}{\Delta\ln n=(\nu^2-1)\,\epsilon .}{Press--Schechter change in halo abundance for a fractional change $\epsilon$ in $\sigma_M$.}{\derived{} (sympy)}{Eq.~\eqref{eq:ms_ps}; Chapter~\ref{ch:satellites}, section ``Growth suppression from $\mu<1$''}
\section{\texorpdfstring{Part~\ref{part:bh}}{Part 3}: Black Holes and Horizons}
\subsection*{Chapter~\ref{ch:blackholes}: Black-hole horizons: Bekenstein, Hawking and the price of a bit}
\iamfsentry{339}{S_{\rm BH}=\frac{k_Bc^3A}{4G\hbar}=\frac{4\pi Gk_BM^2}{\hbar c},\qquad T_{\rm BH}=\frac{\hbar c^3}{8\pi Gk_BM},\qquad A=\frac{16\pi G^2M^2}{c^4}.}{Bekenstein--Hawking entropy, Hawking temperature and horizon area of a Schwarzschild black hole of mass $M$.}{\derived{} (Bekenstein 1973; Hawking 1975)}{Chapter~\ref{ch:blackholes}, section ``A horizon has a temperature and an entropy''}
\iamfsentry{340}{T=\frac{\hbar\kappa}{2\pi k_Bc},}{Temperature of any causal horizon from its surface gravity $\kappa$ (black hole, de Sitter, Unruh).}{\derived{} (standard physics)}{Eq.~\eqref{eq:bh_Tuniv}; Chapter~\ref{ch:blackholes}, section ``A horizon has a temperature and an entropy''}
\iamfsentry{341}{\Gamma_{\rm thermal}=\frac{P}{\kB T_{\rm BH}\ln2}.}{}{\derived{}}{Eq.~\eqref{eq:bh_gamma_def}; Chapter~\ref{ch:blackholes}, section ``Encoding throughput of a black-hole horizon''}
\iamfsentry{342}{\Gamma_{\rm thermal}=\frac{\hbar c^6/(15360\pi G^2M^2)}{\kB\cdot\hbar c^3/(8\pi GM\kB)\cdot\ln2}=\frac{8\pi}{15360\pi}\,\frac{c^3}{GM\ln2}=\frac{c^3}{1920\,GM\ln2}\ \text{bits\,s}^{-1}.}{}{\derived{}}{Eq.~\eqref{eq:bh_gamma}; Chapter~\ref{ch:blackholes}, section ``Encoding throughput of a black-hole horizon''}
\iamfsentry{343}{P_{\rm SB}=\sigma_{\rm SB}\,A_{\rm BH}\,T_{\rm BH}^4,}{}{\derived{}}{Chapter~\ref{ch:blackholes}, section ``The Stefan--Boltzmann luminosity of the horizon''}
\iamfsentry{344}{\begin{aligned}P_{\rm SB}&=\frac{\pi^2\kB^4}{60\hbar^3c^2}\cdot\frac{16\pi G^2M^2}{c^4}\cdot\left(\frac{\hbar c^3}{8\pi GM\kB}\right)^4\\
&=\frac{\pi^2\kB^4}{60\hbar^3c^2}\cdot\frac{16\pi G^2M^2}{c^4}\cdot\frac{\hbar^4c^{12}}{4096\pi^4G^4M^4\kB^4}\\
&=\frac{16}{60\times4096}\,\frac{\hbar c^6}{\pi G^2M^2}=\frac{\hbar c^6}{15360\pi G^2M^2}.\end{aligned}}{}{\derived{}}{Eq.~\eqref{eq:bh_PSB}; Chapter~\ref{ch:blackholes}, section ``The Stefan--Boltzmann luminosity of the horizon''}
\iamfsentry{345}{\frac{P_{\rm SB}}{P_{\rm Hawking}}=1}{}{\derived{}}{Chapter~\ref{ch:blackholes}, section ``The Stefan--Boltzmann luminosity of the horizon''}
\iamfsentry{346}{\frac{dS_{\rm transfer}}{dt}=\frac{P_{\rm Hawking}}{\kB T_{\rm BH}\ln2}=\frac{c^3}{1920\,GM\ln2}=\Gamma_{\rm thermal}.}{}{\derived{}}{Eq.~\eqref{eq:bh_transfer_rate}; Chapter~\ref{ch:blackholes}, section ``The Stefan--Boltzmann luminosity of the horizon''}
\iamfsentry{347}{-\frac{d}{dt}\Big(\frac{S_{\rm BH}}{\kB\ln2}\Big)=\frac{8\pi GM}{\hbar c\ln2}\cdot\frac{\hbar c^4}{15360\pi G^2M^2}=\frac{c^3}{1920\,GM\ln2},}{}{\derived{}}{Chapter~\ref{ch:blackholes}, section ``The Stefan--Boltzmann luminosity of the horizon''}
\iamfsentry{348}{N\,k_BT_{\rm BH}\ln2=T_{\rm BH}S_{\rm BH}=\tfrac12Mc^2 ,}{Horizon bits priced at Landauer's minimum: the total is half the rest energy (Smarr relation).}{\derived{} (Smarr 1973); the Landauer reading is \interp}{Eq.~\eqref{eq:bh_smarr}; Chapter~\ref{ch:blackholes}, section ``The price of the horizon's bits''}
\iamfsentry{349}{\begin{aligned}\frac{T_{\rm BH}S_{\rm BH}}{Mc^2}&=\frac{\sqrt{1-\chi^2}}{2}:\\&\qquad 0.433\ (\chi=0.5),\quad 0.218\ (\chi=0.9),\quad 0.032\ (\chi=0.998).\end{aligned}}{Kerr black hole: fraction of rest energy held as $T_{\rm BH}S_{\rm BH}$ versus spin $\chi$, with three values.}{\derived; values \calc}{Chapter~\ref{ch:blackholes}, section ``The price of the horizon's bits''}
\iamfsentry{350}{\frac{dM}{dt}=-\frac{P_{\rm Hawking}(M)}{c^2}=-\frac{\hbar c^4}{15360\pi G^2M^2},}{}{\derived{}}{Eq.~\eqref{eq:bh_dMdt}; Chapter~\ref{ch:blackholes}, section ``Evaporation and the entropy carried off''}
\iamfsentry{351}{M(t)^3=M_0^3-\frac{\hbar c^4t}{5120\pi G^2},\qquad \tau_{\rm evap}=\frac{5120\pi G^2M_0^3}{\hbar c^4}.}{Mass loss by Hawking radiation and the evaporation time (black-body approximation).}{\derived}{Eq.~\eqref{eq:bh_Mt}; Chapter~\ref{ch:blackholes}, section ``Evaporation and the entropy carried off''}
\iamfsentry{352}{\frac{dS_{\rm tr}}{dt}=\Gamma\big(M(t)\big)=\frac{c^3}{1920\,GM(t)\ln2},\qquad
S_{\rm tr}(t)=\int_0^t\frac{c^3}{1920\,GM(t')\ln2}\,dt'.}{}{\derived{}}{Eq.~\eqref{eq:bh_Str}; Chapter~\ref{ch:blackholes}, section ``Evaporation and the entropy carried off''}
\iamfsentry{353}{S_{\rm tr}(t)=S_{\rm BH,0}-S_{\rm BH}(t)=S_{\rm BH,0}\Big[1-\Big(1-\frac{t}{\tau_{\rm evap}}\Big)^{2/3}\Big].}{}{\derived{}}{Eq.~\eqref{eq:bh_Str_closed}; Chapter~\ref{ch:blackholes}, section ``Evaporation and the entropy carried off''}
\iamfsentry{354}{t=\big(1-2^{-3/2}\big)\,\tau_{\rm evap}=0.646\,\tau_{\rm evap},}{}{\derived{}}{Chapter~\ref{ch:blackholes}, section ``Evaporation and the entropy carried off''}
\iamfsentry{355}{T_{\rm GH}=\frac{\hbar H}{2\pi\kB}.}{}{\calc{}}{Eq.~\eqref{eq:bh_TGH}; Chapter~\ref{ch:blackholes}, section ``Two horizons: hot and cold''}
\iamfsentry{356}{P_{\rm net}=\sigma_{\rm SB}\,A_{\rm BH}\big(T_{\rm BH}^4-T_{\rm GH}^4\big).}{}{\derived{}}{Eq.~\eqref{eq:bh_Pnet}; Chapter~\ref{ch:blackholes}, section ``Two horizons: hot and cold''}
\iamfsentry{357}{\begin{gathered}\frac{\hbar c^3}{8\pi GM_{\rm eq}\kB}=\frac{\hbar H}{2\pi\kB}\quad\Longrightarrow\quad M_{\rm eq}=\frac{c^3}{4GH}=2.32\times10^{22}\,M_\odot\\(H_0=67.4\ \text{km\,s}^{-1}\text{Mpc}^{-1},\ \text{Planck 2018}).\end{gathered}}{}{\derived{}}{Eq.~\eqref{eq:bh_Meq}; Chapter~\ref{ch:blackholes}, section ``Two horizons: hot and cold''}
\iamfsentry{358}{M_{\rm CMB}=\frac{\hbar c^3}{8\pi Gk_BT_{\rm CMB}}=4.5\times10^{22}\ \text{kg}\approx0.6\,\text{lunar masses};}{Mass at which a black hole is in radiative balance with the CMB today.}{\derived; value \calc}{Chapter~\ref{ch:blackholes}, section ``Two horizons: hot and cold''}
\iamfsentry{359}{\frac{S_{\rm info}(R)}{\kB A(R)}\ge\frac{1}{4\lP^2},}{}{\conjecture{}}{Eq.~\eqref{eq:bh_saturation}; Chapter~\ref{ch:blackholes}, section ``When a horizon forms''}
\iamfsentry{360}{\frac{4\pi Gk_BM^2/(\hbar c)}{4\pi(2GM/c^2)^2}=\frac{k_Bc^3}{4\hbar G}=\frac{k_B}{4\ell_P^2}.}{At the Schwarzschild radius the entropy per area equals the holographic bound $k_B/4\ell_P^2$.}{\derived}{Eq.~\eqref{eq:bh_hoop}; Chapter~\ref{ch:blackholes}, section ``When a horizon forms''}
\iamfsentry{361}{M_{\rm BH}=\frac{m_{\rm Pl}}{2\sqrt\pi}\sqrt{S_{\rm collapse}},}{}{\derived{}}{Eq.~\eqref{eq:bh_seed}; Chapter~\ref{ch:blackholes}, section ``When a horizon forms''}
\iamfsentry{362}{\frac{\hbar\eta}{2\pi}=\frac{c^3}{8\pi G}\qquad\Longrightarrow\qquad\eta=\frac{c^3}{4\hbar G}=\frac{1}{4\ell_P^2}.}{Jacobson's Clausius relation fixes $\eta=1/4\ell_P^2$.}{\derived{} (Jacobson 1995)}{Chapter~\ref{ch:blackholes}, section ``Where the 1/4 comes from''}
\subsection*{Chapter~\ref{ch:bekenstein}: The horizon coefficient: why $\eta=1/(4\lP^2)$}
\iamfsentry{363}{\frac14=\frac{2\pi}{8\pi},}{}{\derived{}}{Chapter~\ref{ch:bekenstein}, section ``Overview''}
\iamfsentry{364}{T=\frac{\hbar\kappa}{2\pi\kB c},}{}{\derived{}}{Eq.~\eqref{eq:bk_unruh}; Chapter~\ref{ch:bekenstein}, section ``The problem''}
\iamfsentry{365}{S=\kB\,\eta A,}{}{\derived{}}{Eq.~\eqref{eq:bk_SetaA}; Chapter~\ref{ch:bekenstein}, section ``The problem''}
\iamfsentry{366}{G=\frac{c^3}{4\hbar\eta}.}{}{\derived{}}{Eq.~\eqref{eq:bk_Geta}; Chapter~\ref{ch:bekenstein}, section ``The problem''}
\iamfsentry{367}{|\psi\rangle_{\mathcal S}\otimes|E_0\rangle\;\longrightarrow\;\sum_i c_i\,|s_i\rangle\otimes|E_i\rangle,}{}{\derived{}}{Eq.~\eqref{eq:bk_decoherence}; Chapter~\ref{ch:bekenstein}, section ``Why entropy scales with area: the causal accessibility argument''}
\iamfsentry{368}{\rho_{\mathcal S}={\rm Tr}_{\mathcal E}|\Psi\rangle\langle\Psi|\;\longrightarrow\;\sum_i|c_i|^2\,|s_i\rangle\langle s_i|.}{}{\derived{}}{Eq.~\eqref{eq:bk_diagonal}; Chapter~\ref{ch:bekenstein}, section ``Why entropy scales with area: the causal accessibility argument''}
\iamfsentry{369}{S\propto A.}{}{\interp{}}{Eq.~\eqref{eq:bk_SpropA}; Chapter~\ref{ch:bekenstein}, section ``Why entropy scales with area: the causal accessibility argument''}
\iamfsentry{370}{4=\frac{8\pi}{2\pi}.}{}{\derived{}}{Eq.~\eqref{eq:bk_four}; Chapter~\ref{ch:bekenstein}, section ``The two geometric factors that determine $\eta$''}
\iamfsentry{371}{ds^2=-\frac{\kappa^2\rho^2}{c^2}\,dt^2+d\rho^2+d\mathbf x_\perp^2,}{}{\derived{}}{Eq.~\eqref{eq:bk_rindler}; Chapter~\ref{ch:bekenstein}, section ``The two geometric factors that determine $\eta$''}
\iamfsentry{372}{ds^2_E=\frac{\kappa^2\rho^2}{c^2}\,d\tau^2+d\rho^2+d\mathbf x_\perp^2.}{}{\derived{}}{Eq.~\eqref{eq:bk_euclid}; Chapter~\ref{ch:bekenstein}, section ``The two geometric factors that determine $\eta$''}
\iamfsentry{373}{\frac{\kappa\tau}{c}\sim\frac{\kappa\tau}{c}+2\pi\quad\Longrightarrow\quad\tau\sim\tau+\frac{2\pi c}{\kappa}.}{}{\derived{}}{Eq.~\eqref{eq:bk_period}; Chapter~\ref{ch:bekenstein}, section ``The two geometric factors that determine $\eta$''}
\iamfsentry{374}{\frac{\hbar}{\kB T}=\frac{2\pi c}{\kappa}\quad\Longrightarrow\quad T=\frac{\hbar\kappa}{2\pi\kB c},}{}{\derived{}}{Chapter~\ref{ch:bekenstein}, section ``The two geometric factors that determine $\eta$''}
\iamfsentry{375}{G_{ab}+\Lambda g_{ab}=\frac{8\pi G}{c^4}\,T_{ab}.}{}{\derived{}}{Eq.~\eqref{eq:bk_einstein}; Chapter~\ref{ch:bekenstein}, section ``The two geometric factors that determine $\eta$''}
\iamfsentry{376}{\nabla^2\Phi=4\pi G\rho_m.}{}{\derived{}}{Eq.~\eqref{eq:bk_poisson}; Chapter~\ref{ch:bekenstein}, section ``The two geometric factors that determine $\eta$''}
\iamfsentry{377}{\Omega_3=\oint d\Omega=\int_0^{2\pi}\!\!d\varphi\int_0^\pi\sin\vartheta\,d\vartheta=4\pi:}{}{\derived{}}{Eq.~\eqref{eq:bk_solid}; Chapter~\ref{ch:bekenstein}, section ``The two geometric factors that determine $\eta$''}
\iamfsentry{378}{R_{00}=\frac{8\pi G}{c^4}\Big(\rho_m c^2-\tfrac12\rho_m c^2\Big)=\frac{4\pi G\rho_m}{c^2}.}{}{\derived{}}{Chapter~\ref{ch:bekenstein}, section ``The two geometric factors that determine $\eta$''}
\iamfsentry{379}{\chi^a=-\kappa_g\lambda\,k^a .}{}{\derived{}}{Chapter~\ref{ch:bekenstein}, section ``Jacobson's derivation, step by step''}
\iamfsentry{380}{\delta Q=\int_{\mathcal H}T_{ab}\chi^a\,d\Sigma^b=-\kappa_g\int\lambda\,T_{ab}k^ak^b\,d\lambda\,dA ,\qquad d\Sigma^b=k^b\,d\lambda\,dA .}{}{\derived{}}{Chapter~\ref{ch:bekenstein}, section ``Jacobson's derivation, step by step''}
\iamfsentry{381}{\frac{d\theta}{d\lambda}=-\tfrac12\theta^2-\sigma_{ab}\sigma^{ab}-R_{ab}k^ak^b .}{}{\derived{}}{Chapter~\ref{ch:bekenstein}, section ``Jacobson's derivation, step by step''}
\iamfsentry{382}{\delta A=-\int\lambda\,R_{ab}k^ak^b\,d\lambda\,dA .}{}{\derived{}}{Chapter~\ref{ch:bekenstein}, section ``Jacobson's derivation, step by step''}
\iamfsentry{383}{-\kappa_g\int\lambda\,T_{ab}k^ak^b\,d\lambda\,dA=\frac{\hbar c\,\kappa_g}{2\pi\kB}\,\kB\eta\Big(-\int\lambda\,R_{ab}k^ak^b\,d\lambda\,dA\Big).}{}{\derived{}}{Chapter~\ref{ch:bekenstein}, section ``Jacobson's derivation, step by step''}
\iamfsentry{384}{T_{ab}k^ak^b=\frac{\hbar c\,\eta}{2\pi}\,R_{ab}k^ak^b\quad\text{for all null }k^a\quad\Longrightarrow\quad T_{ab}=\frac{\hbar c\,\eta}{2\pi}\,R_{ab}+f\,g_{ab},}{}{\derived{}}{Chapter~\ref{ch:bekenstein}, section ``Jacobson's derivation, step by step''}
\iamfsentry{385}{0=\frac{\hbar c\,\eta}{2\pi}\,\tfrac12\nabla_bR+\nabla_bf\quad\Longrightarrow\quad f=-\frac{\hbar c\,\eta}{2\pi}\Big(\tfrac12R-\Lambda\Big),}{}{\derived{}}{Chapter~\ref{ch:bekenstein}, section ``Jacobson's derivation, step by step''}
\iamfsentry{386}{R_{ab}-\tfrac12Rg_{ab}+\Lambda g_{ab}=\frac{2\pi}{\hbar c\,\eta}\,T_{ab}.}{}{\derived{}}{Chapter~\ref{ch:bekenstein}, section ``Jacobson's derivation, step by step''}
\iamfsentry{387}{\frac{\hbar c\,\eta}{2\pi}=\frac{c^4}{8\pi G},\qquad\text{equivalently}\qquad\frac{\hbar\eta}{2\pi}=\frac{c^3}{8\pi G},}{}{\derived{}}{Eq.~\eqref{eq:bk_core}; Chapter~\ref{ch:bekenstein}, section ``Jacobson's derivation, step by step''}
\iamfsentry{388}{\eta=\frac{1}{4\lP^2},\qquad\frac14=\frac{2\pi}{8\pi},}{}{\derived{}}{Eq.~\eqref{eq:bk_eta}; Chapter~\ref{ch:bekenstein}, section ``The derivation of the coefficient''}
\iamfsentry{389}{\eta=\frac{c^3}{8\pi G}\cdot\frac{2\pi}{\hbar}=\frac{2\pi}{8\pi}\cdot\frac{c^3}{\hbar G}=\frac14\cdot\frac{1}{\lP^2}=\frac{1}{4\lP^2}.}{}{\derived{}}{Eq.~\eqref{eq:bk_etasolve}; Chapter~\ref{ch:bekenstein}, section ``The derivation of the coefficient''}
\iamfsentry{390}{E_{\rm nat}=\kB T_H=\frac{\hbar\kappa}{2\pi c},}{}{\derived{}}{Eq.~\eqref{eq:bk_Enat}; Chapter~\ref{ch:bekenstein}, section ``Connection to known results''}
\iamfsentry{391}{\delta(Mc^2)=\frac{\kappa c^2}{8\pi G}\,\delta A\qquad\Longrightarrow\qquad\delta A=\frac{8\pi G}{\kappa c^2}\,\delta E ,}{}{\derived{}}{Eq.~\eqref{eq:bk_firstlaw}; Chapter~\ref{ch:bekenstein}, section ``Connection to known results''}
\iamfsentry{392}{\delta A_{\min}=\frac{8\pi G}{\kappa c^2}\cdot\frac{\hbar\kappa}{2\pi c}=\frac{8\pi}{2\pi}\cdot\frac{\hbar G}{c^3}=4\lP^2 ,}{}{\derived{}}{Eq.~\eqref{eq:bk_dAmin}; Chapter~\ref{ch:bekenstein}, section ``Connection to known results''}
\iamfsentry{393}{\eta\cdot\delta A_{\min}=\frac{1}{4\lP^2}\cdot4\lP^2=1\quad\text{(one nat, $\kB$, per minimum encoding area)}.}{}{\calc{}}{Eq.~\eqref{eq:bk_onenat}; Chapter~\ref{ch:bekenstein}, section ``Connection to known results''}
\iamfsentry{394}{\underbrace{\frac{\hbar\eta}{2\pi}}_{\text{Rindler geometry}}=\underbrace{\frac{c^3}{8\pi G}}_{\text{Einstein geometry}}.}{}{\derived{}}{Eq.~\eqref{eq:bk_structure}; Chapter~\ref{ch:bekenstein}, section ``Connection to known results''; see also (\ref{app:fs:90})}
\section{\texorpdfstring{Part~\ref{part:qp}}{Part 4}: Particles and Quantum Records}
\subsection*{Chapter~\ref{ch:quantumrecords}: Records at the quantum scale}
\iamfsentry{395}{|\Psi\rangle_{S\mathcal E}=\sum_i c_i\,|s_i\rangle_S\otimes|e_i\rangle_{\mathcal E}.}{}{\derived{} (unitary evolution of a system and its environment)}{Eq.~\eqref{eq:qr_SE}; Chapter~\ref{ch:quantumrecords}, section ``What decoherence explains, and what it leaves open''}
\iamfsentry{396}{\tau_D\sim\tau_R\left(\frac{\lambda_{\rm th}}{\Delta x}\right)^2,\qquad\lambda_{\rm th}=\frac{\hbar}{\sqrt{2Mk_BT}},}{}{\derived{}}{Eq.~\eqref{eq:qr_tauD}; Chapter~\ref{ch:quantumrecords}, section ``What decoherence explains, and what it leaves open''}
\iamfsentry{397}{2\langle K\rangle=n\,\langle|V|\rangle .}{}{\derived{}}{Eq.~\eqref{eq:qr_virial_n}; Chapter~\ref{ch:quantumrecords}, section ``The virial half: how much bound energy writes''}
\iamfsentry{398}{2K=|V|,\qquad K=\tfrac12|V| .}{}{\derived{} (Clausius 1870~\cite{Clausius1870}; for a circular orbit in $V=-kr^{-n}$, $mv^2/r=nkr^{-n-1}$ gives $2K/|V|=n$ directly)}{Eq.~\eqref{eq:qr_virial_1}; Chapter~\ref{ch:quantumrecords}, section ``The virial half: how much bound energy writes''}
\iamfsentry{399}{\beta_m=\frac{\Omega_m}{2}=0.15765 .}{}{\derived{} (given IAM's Law and the virial partition; $\Omega_m=0.3153$)}{Eq.~\eqref{eq:qr_betam}; Chapter~\ref{ch:quantumrecords}, section ``The virial half: how much bound energy writes''; see also (\ref{app:fs:178}), (\ref{app:fs:319}), (\ref{app:fs:331}), (\ref{app:fs:522})}
\iamfsentry{400}{\tau_D\sim\frac{\hbar R_{\rm vir}}{GM^2}.}{}{\calc{}}{Eq.~\eqref{eq:qr_tauhalo}; Chapter~\ref{ch:quantumrecords}, section ``Collapse sets the rate''}
\iamfsentry{401}{\tau_D\simeq2.5\times10^{-87}\ {\rm s},}{}{\calc{}}{Eq.~\eqref{eq:qr_tauMW}; Chapter~\ref{ch:quantumrecords}, section ``Collapse sets the rate''}
\iamfsentry{402}{\dot I=\int\frac{dn}{dM}\,\frac{M}{m_p}\,H(a)\,\frac{d\ln F_{\rm coll}}{d\ln a}\,dM,}{}{\derived{}}{Eq.~\eqref{eq:qr_Idot_particle}; Chapter~\ref{ch:quantumrecords}, section ``Collapse sets the rate''}
\iamfsentry{403}{\frac{dF_{\rm coll}}{d\ln D}=\sqrt{\frac2\pi}\,\nu_{\min}\,e^{-\nu_{\min}^2/2},\qquad
n_{\rm eff}\equiv\frac{d\ln(dF_{\rm coll}/d\ln D)}{d\ln D}=\nu_{\min}^2-1 .}{}{\derived{} (sympy, \texttt{verify\_records\_measurement\_time.py} A5)}{Eq.~\eqref{eq:qr_neff_nu}; Chapter~\ref{ch:quantumrecords}, section ``Collapse sets the rate''}
\iamfsentry{404}{S_{\rm total}=S_{\rm geo}+S_{\rm info},}{}{\interp{}}{Eq.~\eqref{eq:qr_Stot}; Chapter~\ref{ch:quantumrecords}, section ``The exponent from both directions''; see also (\ref{app:fs:13}), (\ref{app:fs:103})}
\iamfsentry{405}{\begin{aligned}\frac{dS_{\rm info}}{da}&\propto\frac{\dot I}{T_HA_H}\,\frac{1}{aH}\\&\propto\frac{a^{-3}\,a^n\,a^{-3/2}}{a^{-3/2}\,a^{3}}\cdot\frac{1}{a\,a^{-3/2}}=a^{\,n-11/2},
\\S_{\rm info}(a)&\propto\frac{a^{\,n-9/2}}{n-9/2}.\end{aligned}}{}{\derived{}}{Eq.~\eqref{eq:qr_Sinfo}; Chapter~\ref{ch:quantumrecords}, section ``The exponent from both directions''}
\iamfsentry{406}{n=\tfrac72\qquad\text{(top-down).}}{The exponent required by $E(a)$, read from the horizon accounting (top-down).}{\derived{} (given the form; sympy, \texttt{verify\_records\_measurement\_time.py} A6: with $n=7/2$ the integral is $-1/a$)}{Eq.~\eqref{eq:qr_n72}; Chapter~\ref{ch:quantumrecords}, section ``The exponent from both directions''}
\iamfsentry{407}{\dot I\propto\rho_m(a)\,D(a)^{7/2}\,f(a)\,H(a).}{}{\calc{}}{Eq.~\eqref{eq:qr_Idot}; Chapter~\ref{ch:quantumrecords}, section ``The exponent from both directions''}
\iamfsentry{408}{E(a)=\exp\!\left(1-\frac1a\right)}{}{\derived{}}{Eq.~\eqref{eq:qr_Ea}; Chapter~\ref{ch:quantumrecords}, section ``The cosmic horizon as the last environment''; see also (\ref{app:fs:185}), (\ref{app:fs:206}), (\ref{app:fs:318}), (\ref{app:fs:332}), (\ref{app:fs:520})}
\iamfsentry{409}{\mu(a)=\frac{H^2_{\Lambda{\rm CDM}}(a)}{H^2_{\Lambda{\rm CDM}}(a)+\beta_mE(a)H_0^2}<1,\qquad\Sigma(a)=1}{Growth coupling $\mu(a)<1$ for matter and $\Sigma=1$ for light.}{\derived{} (given Eq.~\eqref{eq:Hm} read as an effective Newton constant; that reading is a \conjecture)}{Eq.~\eqref{eq:qr_mu}; Chapter~\ref{ch:quantumrecords}, section ``The cosmic horizon as the last environment''}
\iamfsentry{410}{S=\frac{4\pi\bar\lambda_C^2}{4\lP^2}=\pi\left(\frac{\mP}{m}\right)^2 ,}{}{\conjecture{} (the area law applied to the Compton sphere; it is not saturation of Bekenstein's bound, which gives $S\le2\pi$ here, Chapter~\ref{ch:electronmass})}{Eq.~\eqref{eq:qr_Scompton}; Chapter~\ref{ch:quantumrecords}, section ``Fundamental particles as stable records''}
\iamfsentry{411}{E_{\rm L}=S\,k_BT_{GH}\ln2=\pi\left(\frac{\mP}{m}\right)^2\frac{\hbar H_0\ln2}{2\pi}.}{}{\calc{}}{Eq.~\eqref{eq:qr_ELcompton}; Chapter~\ref{ch:quantumrecords}, section ``Fundamental particles as stable records''}
\iamfsentry{412}{\Gamma_{\rm BH}=\frac{P}{k_BT_{\rm BH}\ln2}=\frac{c^3}{1920\,GM\ln2}\ {\rm bits\,s^{-1}},}{}{\derived{}}{Eq.~\eqref{eq:qr_GammaBH}; Chapter~\ref{ch:quantumrecords}, section ``Black holes as local encoding surfaces''}
\iamfsentry{413}{E=T_{\rm BH}S_{\rm BH}=\tfrac12Mc^2 ,}{}{\derived{}}{Eq.~\eqref{eq:qr_smarr}; Chapter~\ref{ch:quantumrecords}, section ``Black holes as local encoding surfaces''}
\iamfsentry{414}{H_0^{(\rm matter)}=H_0^{(\rm photon)}\sqrt{1+\beta_m}=67.16\sqrt{1.15765}=72.26\ {\rm km\,s^{-1}\,Mpc^{-1}},}{}{\calc{}}{Eq.~\eqref{eq:qr_H0}; Chapter~\ref{ch:quantumrecords}, section ``Observable consequences''}
\iamfsentry{415}{\tau_{\rm IAM}=\frac{\hbar\,(k_BT)^2\ln2}{E_G^3}\;\propto\;T^2m^{-5}\quad(\text{fixed density}),\qquad\tau_{\rm PD}\propto m^{-5/3}.}{IAM gravitational decoherence time of a massive superposition, with the boundary capacity taken as $k_BT/E_G$; compared with Penrose--Di\'osi.}{\derived{} (given the postulated rate and the boundary capacity $\kB T/E_G$, both \conjecture); its $T^2$ and mass scaling are a \prediction}{Eq.~\eqref{eq:qr_tauIAM}; Chapter~\ref{ch:quantumrecords}, section ``Massive superpositions''}
\subsection*{Chapter~\ref{ch:entanglement}: Entanglement and the cost of records}
\iamfsentry{416}{S_{\max}(t)=2\sqrt{1+c(t)^2}.}{Largest CHSH value of a pair dephasing in its pointer basis, as a function of the coherence $c(t)$.}{\derived{} (standard quantum mechanics for pointer-basis dephasing; Horodecki criterion)}{Eq.~\eqref{eq:ent:smax}; Chapter~\ref{ch:entanglement}, section ``What dephasing does to a Bell test''; see also (\ref{app:fs:428})}
\subsection*{Chapter~\ref{ch:measurement}: Measurement as the writing of a record}
\iamfsentry{417}{Q_L\equiv\kB T_D\ln2}{IAM measurement criterion: an interaction is irreversible when the heat dissipated in the detector reaches $k_BT_D\ln2$.}{\conjecture{} (as labeled); see Bennett's erasure result (\openprob)}{Eq.~\eqref{eq:mp_QL}; Chapter~\ref{ch:measurement}, section ``The criterion''}
\iamfsentry{418}{\tau_{\rm IAM}=\frac{\hbar\,(\kB T)^2\ln2}{E_G^3}}{}{\conjecture{}}{Eq.~\eqref{eq:mp_tauIAM}; Chapter~\ref{ch:measurement}, section ``The gravitational channel''}
\iamfsentry{419}{F_{\mathrm{IAM}}(Q,T)=1-\frac{1}{e}\exp\!\left(1-\frac{\kB T\ln2}{Q}\right)=1-\exp\!\left(-\frac{Q_L}{Q}\right),}{Proposed erasure fidelity of a quantum eraser as a function of dissipated energy; the form carries over the activation ramp.}{\conjecture{} (form not derived); the universal curve is a \prediction}{Eq.~\eqref{eq:F}; Chapter~\ref{ch:measurement}, section ``The quantum eraser''}
\iamfsentry{420}{S_{\max}=2\sqrt{1+c^2}}{}{\derived{}}{Eq.~\eqref{eq:mp_smax}; Chapter~\ref{ch:measurement}, section ``Entanglement and Bell tests''}
\subsection*{Chapter~\ref{ch:gravdec}: Gravitational decoherence at the quantum level}
\iamfsentry{421}{\tau_{\mathrm{DP}}\simeq\frac{\hbar}{E_G},\qquad E_G\simeq\frac{Gm^2}{R},}{Di\'osi--Penrose collapse time from the gravitational self-energy of the mass-distribution difference.}{definition (the Di\'osi--Penrose proposal)}{Eq.~\eqref{eq:gd_dp}; Chapter~\ref{ch:gravdec}, section ``Di\'osi--Penrose and the IAM rate''}
\iamfsentry{422}{\frac{dN_{\mathrm{bits}}}{dt}=\frac{E_G^2}{\hbar\,\kB T\ln2},}{Postulated rate of information production of a massive superposition.}{\conjecture{} (postulated)}{Eq.~\eqref{eq:gd_bitrate}; Chapter~\ref{ch:gravdec}, section ``Di\'osi--Penrose and the IAM rate''}
\iamfsentry{423}{\ln E_q(t)=\int_0^t\frac{\Gamma_{\mathrm{info}}(t')}{S_{\mathrm{boundary}}}\,dt'=\frac{E_G^3}{\hbar\,\kB^2T^2\ln2}\,t,\qquad
 \boxed{\;\tau_{\mathrm{IAM}}=\frac{\hbar\,\kB^2T^2\ln2}{E_G^3}\;}}{Integrated decoherence of a massive superposition with boundary capacity $k_BT/E_G$: an exponential with time constant $\tau_{\rm IAM}$.}{\derived{} (given the postulated rate and the boundary capacity, both \conjecture); $T^2$ and mass scaling \prediction}{Eq.~\eqref{eq:tauIAM}; Chapter~\ref{ch:gravdec}, section ``Di\'osi--Penrose and the IAM rate''}
\iamfsentry{424}{E_q(\eta)=\exp\!\left(1-\frac1\eta\right),\qquad C_{\mathrm{IAM}}(t)=1-\frac{E_q(t/\tau_{\mathrm{IAM}})}{e},}{}{\conjecture{}}{Eq.~\eqref{eq:gd_ramp}; Chapter~\ref{ch:gravdec}, section ``Di\'osi--Penrose and the IAM rate''}
\iamfsentry{425}{P_{\rm IAM}=\frac{k_BT\ln2}{\tau_{\rm IAM}}=\frac{E_G^3}{\hbar\,k_BT},\qquad\frac{dn}{dt}=\frac{E_G^3}{\hbar^2k_BT\,\omega_0}.}{}{\calc{}}{Eq.~\eqref{eq:gd_heating}; Chapter~\ref{ch:gravdec}, section ``Quantitative predictions''}
\iamfsentry{426}{\Gamma_{\rm IAM}(\eta)=\frac{dE_q}{d\eta}=\frac{1}{\eta^2}\exp\!\left(1-\frac1\eta\right),\qquad\eta=t/\tau_{\rm IAM},}{}{\conjecture{}}{Eq.~\eqref{eq:gd_gamma}; Chapter~\ref{ch:gravdec}, section ``Lindblad dynamics with a time-dependent rate''}
\iamfsentry{427}{\frac{d\rho}{d\eta}=\Gamma(\eta)\left(L\rho L^\dagger-\tfrac12\{L^\dagger L,\rho\}\right),}{}{\conjecture{}}{Eq.~\eqref{eq:gd_lindblad}; Chapter~\ref{ch:gravdec}, section ``Lindblad dynamics with a time-dependent rate''}
\subsection*{Chapter~\ref{ch:nonlocal}: Nonlocality and the boundary of reality}
\iamfsentry{428}{S_{\max}(t)=2\sqrt{1+c(t)^2}}{Largest CHSH value under pointer-basis dephasing (repeated from Part~\ref{part:qp}).}{\derived{} (standard quantum mechanics for pointer-basis dephasing)}{Chapter~\ref{ch:nonlocal}, section ``The IAM account''; see also (\ref{app:fs:416})}
\subsection*{Chapter~\ref{ch:electroweak}: Electroweak symmetry breaking and the matter sector}
\iamfsentry{429}{2\langle K\rangle=k\,\langle V\rangle.}{Time-averaged virial theorem for a potential homogeneous of degree $k$.}{\derived{} (Clausius; Euler's theorem)}{Eq.~\eqref{eq:ew:virial}; Chapter~\ref{ch:electroweak}, section ``Which interactions bind''; see also (\ref{app:fs:48})}
\iamfsentry{430}{V(r)=-\frac43\frac{\alpha_s}{r}+\sigma r,\qquad\sigma\approx0.18\ {\rm GeV}^2=0.91\ {\rm GeV\,fm}^{-1}.}{Cornell static quark potential, with string tension $\sigma\approx0.18$~GeV$^2$.}{\calc{} (as labeled; the Cornell form is a published phenomenological potential)}{Chapter~\ref{ch:electroweak}, section ``Which interactions bind''}
\iamfsentry{431}{\beta_m=\frac{\Omega_b}{2}+\frac{\Omega_{\rm dm}}{2}=0.0247+0.1330=0.1577.}{}{\calc{}}{Eq.~\eqref{eq:ew:betam}; Chapter~\ref{ch:electroweak}, section ``The composition of $\beta_m$''}
\subsection*{Chapter~\ref{ch:koide}: Three charged leptons and the Koide relation}
\iamfsentry{432}{Q\equiv\frac{m_e+m_\mu+m_\tau}{\big(\sqrt{m_e}+\sqrt{m_\mu}+\sqrt{m_\tau}\big)^2}.}{Koide's relation for the charged-lepton pole masses (PDG 2024; $0.66666051$ with the 2022 $m_\tau$).}{\observed}{Eq.~\eqref{eq:ko:Q}; Chapter~\ref{ch:koide}, section ``Koide's relation''}
\iamfsentry{433}{Q(\mu=m_\tau)=0.667824,\qquad Q(\mu=M_Z)=0.667840,}{}{\calc{}}{Chapter~\ref{ch:koide}, section ``Koide's relation''}
\iamfsentry{434}{\vec s=\big(\sqrt{m_e},\sqrt{m_\mu},\sqrt{m_\tau}\big)=(0.7148,\;10.2790,\;42.1536)\ {\rm MeV}^{1/2}.}{}{\calc{}}{Chapter~\ref{ch:koide}, section ``The square-root vector and its angle''}
\iamfsentry{435}{\cos^2\theta=\frac{(\vec s\cdot(1,1,1))^2}{3\,|\vec s|^2}=\frac{(\sum_k\sqrt{m_k})^2}{3\sum_km_k}=\frac{1}{3Q}.}{}{\derived{}}{Eq.~\eqref{eq:ko:angle}; Chapter~\ref{ch:koide}, section ``The square-root vector and its angle''}
\iamfsentry{436}{\sqrt{m_k}=x\big[1+r\cos(\delta+2\pi k/3)\big],\qquad k=0,1,2.}{}{\derived{}}{Eq.~\eqref{eq:ko:param}; Chapter~\ref{ch:koide}, section ``The circle and the three phases''}
\iamfsentry{437}{\sum_k\cos\phi_k=0,\qquad\sum_k\cos^2\phi_k=\tfrac32,}{}{\derived{}}{Eq.~\eqref{eq:ko:Z3}; Chapter~\ref{ch:koide}, section ``The circle and the three phases''}
\iamfsentry{438}{Q=\frac13\Big(1+\frac{r^2}{2}\Big).}{}{\derived{}}{Eq.~\eqref{eq:ko:Qr}; Chapter~\ref{ch:koide}, section ``The circle and the three phases''}
\iamfsentry{439}{T=\frac{\hbar\kappa}{2\pi\kB}.}{}{\derived{}}{Eq.~\eqref{eq:ko:TU}; Chapter~\ref{ch:koide}, section ``The encoding on a charge orbit''; see also (\ref{app:fs:452})}
\iamfsentry{440}{\frac{\hbar\eta}{2\pi}=\frac{c^3}{8\pi G}\quad\Longrightarrow\quad\eta=\frac{c^3}{4\hbar G}=\frac1{4\lP^2}.}{}{\derived{}}{Eq.~\eqref{eq:ko:eta}; Chapter~\ref{ch:koide}, section ``The encoding on a charge orbit''}
\iamfsentry{441}{\eta\,\delta A_{\min}=1.}{}{\derived{}}{Eq.~\eqref{eq:ko:onebit}; Chapter~\ref{ch:koide}, section ``The encoding on a charge orbit''}
\iamfsentry{442}{\frac{\partial}{\partial p_k}\Big(H-\lambda\sum_jp_j\Big)=0\quad\Longrightarrow\quad p_k=\frac1K,\qquad w_k=\frac WK.}{}{\derived{} (given the set-up)}{Eq.~\eqref{eq:ko:maxent}; Chapter~\ref{ch:koide}, section ``The encoding on a charge orbit''}
\iamfsentry{443}{\sqrt{m(\phi)}=a_0+\sum_{n\ge1}\big(a_n\cos n\phi+b_n\sin n\phi\big).}{}{\derived{}}{Eq.~\eqref{eq:ko:fourier}; Chapter~\ref{ch:koide}, section ``The encoding on a charge orbit''}
\iamfsentry{444}{E^{\rm grad}_n=\frac{n^2}{2}\,\omega_0^2,}{}{\derived{}}{Eq.~\eqref{eq:ko:grad}; Chapter~\ref{ch:koide}, section ``The encoding on a charge orbit''}
\iamfsentry{445}{\frac{w_n}{w_1}\sim\exp\!\Big[-\frac{(n^2-1)\,\omega_0^2}{2\kB T_{\rm enc}}\Big].}{}{\conjecture{} (the encoding temperature and $\omega_0$ are left open)}{Eq.~\eqref{eq:ko:boltz}; Chapter~\ref{ch:koide}, section ``The encoding on a charge orbit''}
\iamfsentry{446}{\sqrt{m(\phi)}=x+y\cos\phi,\qquad x>0.}{}{\conjecture{}}{Eq.~\eqref{eq:ko:two}; Chapter~\ref{ch:koide}, section ``The encoding on a charge orbit''}
\iamfsentry{447}{w_{\rm DC}=\frac1{2\pi}\int_0^{2\pi}x^2\,d\phi=x^2,\qquad w_{\rm harm}=\frac1{2\pi}\int_0^{2\pi}(y\cos\phi)^2\,d\phi=\frac{y^2}{2}.}{}{\conjecture{}}{Eq.~\eqref{eq:ko:weights}; Chapter~\ref{ch:koide}, section ``The encoding on a charge orbit''}
\iamfsentry{448}{x^2=\frac{y^2}{2}\quad\Longrightarrow\quad\frac yx=\sqrt2.}{}{\derived{} (given Steps 1--7)}{Eq.~\eqref{eq:ko:ratio}; Chapter~\ref{ch:koide}, section ``The encoding on a charge orbit''}
\iamfsentry{449}{\sqrt{m_k}=x\big(1+\sqrt2\cos\phi_k\big)>0\quad\Longleftrightarrow\quad|\phi_k-\pi|>\pi/4\ ({\rm mod}\ 2\pi).}{}{\derived{}}{Eq.~\eqref{eq:ko:pos}; Chapter~\ref{ch:koide}, section ``The number of generations''}
\iamfsentry{450}{Q=\frac{3x^2+\tfrac32y^2}{(3x)^2}=\frac13\Big(1+\frac{y^2}{2x^2}\Big)=\frac13(1+1)=\frac23.}{}{\derived{}}{Eq.~\eqref{eq:ko:thm2}; Chapter~\ref{ch:koide}, section ``The Koide value''}
\iamfsentry{451}{\delta=0:\qquad m_e=m_\mu=26.92\ {\rm MeV},\qquad m_\tau=1829.26\ {\rm MeV}.}{}{\calc{}}{Eq.~\eqref{eq:ko:delta0}; Chapter~\ref{ch:koide}, section ``The offset''}
\subsection*{Chapter~\ref{ch:electronmass}: The electron rest mass as a fixed point}
\iamfsentry{452}{T=\frac{\hbar\kappa}{2\pi\kB}}{}{\derived{}}{Eq.~\eqref{eq:em:T}; Chapter~\ref{ch:electronmass}, section ``One temperature formula, two horizons''; see also (\ref{app:fs:439})}
\iamfsentry{453}{S=\frac{4\pi\bar\lambda_C^2}{4\lP^2}=\pi\Big(\frac{\mP}{m}\Big)^2=1.79\times10^{45}\quad(m=m_e).}{}{\conjecture{}}{Eq.~\eqref{eq:em:S}; Chapter~\ref{ch:electronmass}, section ``The ingredients''}
\iamfsentry{454}{E_{\rm bit}=\kB T_{GH}\ln2=\frac{\hbar H_0\ln2}{2\pi}=2.54\times10^{-53}\ {\rm J}.}{}{\derived{} (given IAM's Law and the choice of the cosmic horizon as the pricing surface)}{Eq.~\eqref{eq:em:Ebit}; Chapter~\ref{ch:electronmass}, section ``The ingredients''}
\iamfsentry{455}{N(m)=\Big(\frac{\mP}{m}\Big)^{3/2}=\Big(\frac S\pi\Big)^{3/4}=3.69\times10^{33}\quad(m=m_e).}{}{\conjecture{}}{Eq.~\eqref{eq:em:N}; Chapter~\ref{ch:electronmass}, section ``The ingredients''}
\iamfsentry{456}{f(\alpha)\big|_{\rm spatial}=\alpha^{3/2}\cdot\alpha=\alpha^{5/2}=4.55\times10^{-6}.}{}{\conjecture{}}{Eq.~\eqref{eq:em:fs}; Chapter~\ref{ch:electronmass}, section ``The ingredients''}
\iamfsentry{457}{f(\alpha)\big|_{\rm temporal}=\alpha\cdot\alpha^{3/2}=\alpha^{5/2}.}{}{\conjecture{}}{Eq.~\eqref{eq:em:ft}; Chapter~\ref{ch:electronmass}, section ``The ingredients''}
\iamfsentry{458}{mc^2=E_{\rm bit}\,\frac{N(m)}{f(\alpha)}.}{}{\conjecture{}}{Eq.~\eqref{eq:em:fp}; Chapter~\ref{ch:electronmass}, section ``The fixed point and its solution''}
\iamfsentry{459}{mc^2=\frac{\hbar H_0\ln2}{2\pi}\cdot\frac{(\mP/m)^{3/2}}{\alpha^{5/2}}\quad\Longrightarrow\quad
m^{5/2}=\frac{\hbar H_0\ln2\;\mP^{3/2}}{2\pi\,\alpha^{5/2}c^2},}{}{\derived{}}{Eq.~\eqref{eq:em:m52}; Chapter~\ref{ch:electronmass}, section ``The fixed point and its solution''}
\iamfsentry{460}{m_*=(2\pi)^{-2/5}\,\mathcal B,\qquad\mathcal B\equiv\left[\frac{\hbar H_0\ln2\;\mP^{3/2}}{\alpha^{5/2}c^2}\right]^{2/5}.}{}{\derived{} (given the ingredients)}{Eq.~\eqref{eq:em:mstar}; Chapter~\ref{ch:electronmass}, section ``The fixed point and its solution''}
\iamfsentry{461}{m_*=0.5762\,m_e.}{}{\calc{}}{Eq.~\eqref{eq:em:short}; Chapter~\ref{ch:electronmass}, section ``The fixed point and its solution''}
\iamfsentry{462}{m_e=(2\pi)^{-1/10}\left[\frac{\hbar H_0\ln2\;\mP^{3/2}}{\alpha^{5/2}c^2}\right]^{2/5}=1.0000066\,m_e\quad(H_0=67.4).}{The electron-mass fixed point; the factor $(2\pi)^{3/10}$ inside the prefactor $(2\pi)^{-1/10}$ was found by numerical search.}{\fitted{} (prefactor $(2\pi)^{3/10}$ selected numerically by search); precision limited by $H_0$ to $\pm0.3\,\%$}{Eq.~\eqref{eq:em:result}; Chapter~\ref{ch:electronmass}, section ``The fixed point and its solution''}
\section{\texorpdfstring{Part~\ref{part:3}}{Part 5}: The Quantum Order of Informational Actualization: Qubits and Semiconductors}
\subsection*{Chapter~\ref{ch:scprimer}: The superconducting qubit as an encoding surface}
\iamfsentry{463}{\xqp^{\mathrm{th}}(T)=\sqrt{\frac{2\pi\kB T}{\Delta}}\,e^{-\Delta/\kB T}}{Thermal-equilibrium quasiparticle density in a BCS superconductor at $T\ll T_c$, normalized per Cooper pair near the gap.}{\derived{} (standard BCS physics)}{Chapter~\ref{ch:scprimer}, section ``Three loss channels''}
\iamfsentry{464}{\Gamma_1^{\mathrm{qp}}=\xqp\,\frac{\omega_q}{\pi}\sqrt{\frac{2\Delta}{\hbar\omega_q}} .}{Qubit relaxation rate from quasiparticle tunneling across the junction (Catelani et al.\ 2011).}{\derived{} (published theory)}{Eq.~\eqref{eq:catelani}; Chapter~\ref{ch:scprimer}, section ``Three loss channels''}
\subsection*{Chapter~\ref{ch:xqp}: The cost of holding a qubit}
\iamfsentry{465}{x_{\rm qp}=n_{\rm qp}/n_{\rm cp}\sim10^{-8}\text{--}10^{-6},\qquad n_{\rm cp}=2\nu_0\Delta\approx4\times10^6\,\mu{\rm m}^{-3},}{Quasiparticle density in the field's normalization; $n_{\rm cp}=2\nu_0\Delta$ counts pairs within $\Delta$ of the Fermi surface. Defines $x_{\rm qp}$.}{definition; the floor $\sim10^{-7}$ is \observed}{Chapter~\ref{ch:xqp}, section ``The anomaly''}
\iamfsentry{466}{x_{\rm qp}=\frac{2N\,\tau_{\rm qp}}{\tau_{\rm TLS}\,n_{\rm cp}\,V}.}{Steady-state quasiparticle density from $N$ erasure sites, two quasiparticles per event, diffusing over the island volume $V$.}{\derived{} (from the proposal and the three facts above)}{Eq.~\eqref{eq:xqp}; Chapter~\ref{ch:xqp}, section ``From events to quasiparticles''}
\iamfsentry{467}{V_{\rm eff}=4\pi\int_0^\infty e^{-2r/\lamL}\,r^2\,dr=\pi\lamL^3}{}{\derived{} (for that weight)}{Eq.~\eqref{eq:xqp_veff}; Chapter~\ref{ch:xqp}, section ``From events to quasiparticles''}
\iamfsentry{468}{x_{\rm qp}=\frac{x_0}{1-g},\qquad g=\frac{2\phi\,\tau_{\rm qp}}{\tau_{\rm TLS}},}{}{\derived{}}{Eq.~\eqref{eq:xqp_feedback}; Chapter~\ref{ch:xqp}, section ``From events to quasiparticles''}
\iamfsentry{469}{\frac{x_{\rm qp}\,n_{\rm cp}\,V\,\tau_{\rm TLS}}{N\,\tau_{\rm qp}}=2.}{Same-device test: the dimensionless combination the proposal predicts equals the pair-breaking yield, 2.}{\prediction}{Chapter~\ref{ch:xqp}, section ``The test''}
\subsection*{Chapter~\ref{ch:ascoreqc}: The gate error and its floor: quantum processors}
\iamfsentry{470}{\varepsilon=-\ln(1-p)=p+\tfrac12p^2+\dots}{Gate error in nats from the two-qubit error probability; equals $p$ to within about $0.5\,\%$ for $p<10^{-2}$ ($\varepsilon/p-1\simeq p/2$).}{definition}{Eq.~\eqref{eq:Agate}; Chapter~\ref{ch:ascoreqc}, section ``From gate error to information cost''}
\subsection*{Chapter~\ref{ch:thermaln}: How the floor moves with temperature}
\iamfsentry{471}{\frac{d\ln p_{\rm eq}}{d\ln T}=M\,(1-p_{\rm eq})}{Local logarithmic slope of the held-record thermal floor $p_{\rm eq}=1/(1+e^M)$; not a constant (6.85 at 35\,mK, 16.0 at 15\,mK for 5\,GHz).}{\derived{}}{Eq.~\eqref{eq:thermal_slope}; Chapter~\ref{ch:thermaln}, section ``The temperature dependence of a floor''}
\subsection*{Chapter~\ref{ch:walls}: Coherence walls and the coherence optimum}
\iamfsentry{472}{p_{2Q}=p_{\mathrm{coh}}+p_{\mathrm{mat}}+p_{\mathrm{ctrl}} .}{Three-part decomposition of a superconducting two-qubit error: coherence, material, control.}{definition; the material/control split is a model, \conjecture}{Eq.~\eqref{eq:pdecomp}; Chapter~\ref{ch:walls}, section ``Three components of a gate error''}
\iamfsentry{473}{\frac{T_1^*}{T_{1,\mathrm{free}}}=\frac{r}{1+r},\qquad r=\sqrt{a/b}.}{The coherence optimum of the model $p=a/T_1+b/(T_{1,\mathrm{free}}-T_1)$; its location is set by each device's constants $a$ and $b$.}{\derived{} (within the model)}{Eq.~\eqref{eq:t1star}; Chapter~\ref{ch:walls}, section ``A coherence optimum''}
\iamfsentry{474}{p_{2Q}=\frac{A}{P}+BP+p_{\mathrm{ctrl}},\qquad B=\eta^2\Bigl(\frac{\pi}{2}\Bigr)^2\frac{\Gamma_{\mathrm{heat}}\,t_g}{P_{\mathrm{ref}}},}{}{\derived{}}{Eq.~\eqref{eq:walls_mhi}; Chapter~\ref{ch:walls}, section ``The same optimum in other classes''}
\subsection*{Chapter~\ref{ch:qplatforms}: Each platform on its own gauge: qubits}
\iamfsentry{475}{p_{\rm eq}=\frac{1}{1+e^{M}},\qquad M=\frac{E}{\kB T},}{Equilibrium occupation of the wrong level of a two-level record with gap $E=hf$ against a bath at $T$: the held-record thermal floor, the left-hand mark of the qubit gauge.}{\derived{} (detailed balance)}{Eq.~\eqref{eq:qp_peq}; Chapter~\ref{ch:qplatforms}, section ``Three marks on every qubit gauge''}
\iamfsentry{476}{p_{\rm th}(t)=\Gamma_\uparrow t=p_{\rm eq}\,\frac{t}{T_1}}{Thermal floor of one gate of duration $t\ll T_1$: $6.2\times10^{-7}$ for a 5\,GHz transmon at 35\,mK, $T_1=68\,\mu$s, $t=40$\,ns.}{\derived{} (detailed balance)}{Eq.~\eqref{eq:qp_pth}; Chapter~\ref{ch:qplatforms}, section ``Three marks on every qubit gauge''}
\iamfsentry{477}{p(\Omega)=\frac{a}{\Omega}+b\,\Omega^2+p_{\rm ctrl},\qquad \Omega^*=\Bigl(\frac{a}{2b}\Bigr)^{1/3},\qquad p(\Omega^*)-p_{\rm ctrl}=\frac{3a}{2\Omega^*},}{Drive optimum of a Rydberg gate: decay falls with the drive, blockade leakage and scattering rise with it.}{\derived{} within a \conjecture{} model; $a$, $b$ unmeasured}{Eq.~\eqref{eq:qp_rydberg}; Chapter~\ref{ch:qplatforms}, section ``Neutral atoms''}
\subsection*{Chapter~\ref{ch:cmos}: Landauer at the transistor gate: CMOS and the Dennard wall}
\iamfsentry{478}{\kB T_j\ln2=3.33\times10^{-21}\ \mathrm J=0.021\ \mathrm{eV}.}{Landauer floor for logic at a junction temperature $T_j=348$\,K.}{\calc}{Chapter~\ref{ch:cmos}, section ``The floor for logic''}
\iamfsentry{479}{E_{\min}=\kB T\ln(1/p)}{}{\calc{}}{Eq.~\eqref{eq:reliableswitch}; Chapter~\ref{ch:cmos}, section ``Dennard scaling and its end''}
\subsection*{Chapter~\ref{ch:chipgen}: Chips, generations and nodes on their own gauge}
\iamfsentry{480}{n_{\rm floor}=\frac{\ln R}{-\ln(1-s)}}{Generations after which a chip reading $R$ times its floor and improving by a constant fraction $s$ per generation would reach the floor: 6.2 at $s=0.646$, 18.5 at $s=1-1/\sqrt2\approx0.293$ for $R=600$.}{\derived; values \calc}{Eq.~\eqref{eq:cg_nfloor}; Chapter~\ref{ch:chipgen}, section ``Generations and nodes''}
\section{\texorpdfstring{Part~\ref{part:4}}{Part 6}: Cellular Physics: Thermodynamics of the Methylome}
\subsection*{Chapter~\ref{ch:bridge}: From the horizon to the nucleus}
\iamfsentry{481}{E_{\mathrm{bit}}=\kB T\ln 2 .}{Landauer's principle: the least heat dissipated in erasing one bit at reservoir temperature $T$.}{\derived{} (Landauer 1961); \observed{} (B\'erut 2012; Hong 2016)}{Eq.~\eqref{eq:landauer}; Chapter~\ref{ch:bridge}, section ``One cost, many surfaces''}
\subsection*{Chapter~\ref{ch:astrogenetics}: Astro-Genetics: the cell read as a thermometer}
\iamfsentry{482}{H(\beta)=-\beta\log_2\beta-(1-\beta)\log_2(1-\beta),}{Shannon (binary) entropy of a locus with methylated fraction $\beta$, in bits.}{definition; its properties are \derived}{Chapter~\ref{ch:astrogenetics}, section ``One instrument, one gauge''; see also (\ref{app:fs:494})}
\subsection*{Chapter~\ref{ch:landauer}: Landauer's cost at body temperature}
\iamfsentry{483}{e_{\mathrm{bit}}=\kB\Tb\ln2=(1.380649\times10^{-23})(310.15)(0.693147)
  =2.968\times10^{-21}\ \mathrm{J},}{Cost of one bit at body temperature, $310.15$\,K.}{\calc}{Eq.~\eqref{eq:ebit}; Chapter~\ref{ch:landauer}, section ``The cost of one bit''}
\iamfsentry{484}{M=\frac{\dGATP}{R\Tb}=\frac{54\,000}{8.314462618\times310.15}=20.94 .}{The Mahaffey number of the cell: cytosolic ATP drive over $RT$ at body temperature.}{\calc{} (from a definition)}{Eq.~\eqref{eq:M}; Chapter~\ref{ch:landauer}, section ``The Mahaffey number''}
\iamfsentry{485}{M=\frac{\Edrive}{\kB T}.}{The Mahaffey number for any system: drive energy per event over $k_BT$ (no $\ln2$).}{definition}{Chapter~\ref{ch:landauer}, section ``The Mahaffey number''}
\iamfsentry{486}{\frac{M}{\ln2}=\frac{\dGATP}{N_A\,\kB\Tb\ln2}=30.21 .}{Landauer bits one ATP can pay for: $M$ in Landauer units.}{\calc{} (a unit conversion)}{Eq.~\eqref{eq:bitsperATP}; Chapter~\ref{ch:landauer}, section ``The Mahaffey number''}
\iamfsentry{487}{M_{\rm transmon}=\frac{\Delta_{\rm Al}\ln2}{\kB\,(\Delta_{\rm Al}/\kB)}=\ln2,\qquad \frac{M_{\rm transmon}}{\ln2}=1 ,}{}{\derived{}}{Eq.~\eqref{eq:Mtransmon}; Chapter~\ref{ch:landauer}, section ``The cell beside a transistor and a qubit''}
\iamfsentry{488}{\frac{\text{chemical cost of one mark}}{\text{Landauer floor of one bit}}\gtrsim\frac{M}{\ln2}\approx30 .}{Cellular margin: chemical cost of one methyl mark over the Landauer floor of one bit.}{\calc{}}{Eq.~\eqref{eq:markmargin}; Chapter~\ref{ch:landauer}, section ``What a methyl mark costs the cell''}
\iamfsentry{489}{E_{\mathrm{floor}}=N\,\kB\Tb\ln2=2.822\times10^{7}\times2.968\times10^{-21}
  =8.38\times10^{-14}\ \mathrm{J},}{Landauer floor for one rewrite of all 28,217,448 CpGs.}{\calc}{Eq.~\eqref{eq:efloor}; Chapter~\ref{ch:landauer}, section ``The floor for one copy of the methylome''}
\iamfsentry{490}{\frac{\Delta\Delta G}{\kB T}=\ln\frac{1}{0.02\text{--}0.10}=2.3\text{--}3.9 .}{Discrimination energy implied by maintenance error rates of 2--10\,\% (equilibrium-discrimination limit).}{\calc{} (from Genereux et al.\ 2005)}{Chapter~\ref{ch:landauer}, section ``Margin per write, and fidelity per write''}
\iamfsentry{491}{\phi=\frac{E_{\rm hold}}{M\,\kB T}=0.163\pm0.006 .}{Fraction of one ATP spent per held bit: holding energy over $M$.}{\measured}{Eq.~\eqref{eq:phi}; Chapter~\ref{ch:landauer}, section ``Margin per write, and fidelity per write''}
\iamfsentry{492}{H(C_i)=-p_i\log_2p_i-(1-p_i)\log_2(1-p_i)}{}{\derived{}}{Eq.~\eqref{eq:sanchezH}; Chapter~\ref{ch:landauer}, section ``Where this stands: the information thermodynamics of methylation''}
\iamfsentry{493}{E_R=I_R\,\kB T\ln2 .}{}{\derived{} (from Landauer's principle, given Eq.~\eqref{eq:sanchezH})}{Eq.~\eqref{eq:sanchezER}; Chapter~\ref{ch:landauer}, section ``Where this stands: the information thermodynamics of methylation''}
\subsection*{Chapter~\ref{ch:surface}: The methylome as an encoding surface}
\iamfsentry{494}{H(\beta)=-\beta\log_2\beta-(1-\beta)\log_2(1-\beta),}{Shannon (binary) entropy of a locus with methylated fraction $\beta$, in bits.}{definition; its properties are \derived}{Eq.~\eqref{eq:H}; Chapter~\ref{ch:surface}, section ``What $\beta$ measures''; see also (\ref{app:fs:482})}
\iamfsentry{495}{H(\bbar)\ \ge\ \overline{H(\beta)},\qquad \bbar=\frac1n\sum_i\beta_i ,}{Jensen's inequality: the entropy of the mean exceeds the mean of the entropies.}{\derived}{Eq.~\eqref{eq:jensen}; Chapter~\ref{ch:surface}, section ``The mean of the entropies''}
\iamfsentry{496}{\overline{H}=\frac1n\sum_i H(\beta_i),}{Mean of the per-site entropies, the statistic of Met-A.}{definition}{Eq.~\eqref{eq:meanH}; Chapter~\ref{ch:surface}, section ``The mean of the entropies''}
\iamfsentry{497}{\beta_{\mathrm{ss}}=\frac{d}{u+d}.}{Steady state of a two-state maintenance chain with loss $u$ and gain $d$ per division.}{\derived{} (toy model)}{Eq.~\eqref{eq:steady}; Chapter~\ref{ch:surface}, section ``A pattern is a steady state''}
\subsection*{Chapter~\ref{ch:ledgers}: Two ledgers and the virial balance}
\iamfsentry{498}{2\langle K\rangle=k\,\langle U\rangle .}{Virial theorem for a potential homogeneous of degree $k$ (as in Chapter~\ref{ch:virial_law}).}{\derived{}}{Eq.~\eqref{eq:virial}; Chapter~\ref{ch:ledgers}, section ``What the virial theorem says''}
\iamfsentry{499}{Mc^2=2\,T_HS ,}{Smarr relation for a Schwarzschild black hole: rest energy is twice $T_HS$.}{\derived{} (Smarr 1973); one-solar-mass check \calc}{Eq.~\eqref{eq:smarr}; Chapter~\ref{ch:ledgers}, section ``Where the halves really are''}
\subsection*{Chapter~\ref{ch:floorbreach}: A full surface: the horizon and the cell}
\iamfsentry{500}{A_{\max}^{\text{Met-A}}=\frac{1}{0.330263}=3.03,}{Far end of the Met-A gauge for neutrophils: every identity site at a coin flip, one bit over the healthy reference.}{\calc{} (from the healthy reference alone)}{Eq.~\eqref{eq:Amax}; Chapter~\ref{ch:floorbreach}, section ``The full surface of a cell''}
\subsection*{Chapter~\ref{ch:gauge}: One gauge for a cell}
\iamfsentry{501}{A=\frac{\text{reading}}{\text{healthy reference}} .}{The informational fidelity ratio: a reading over the same cell type's healthy reading on the same instrument.}{definition}{Eq.~\eqref{eq:A}; Chapter~\ref{ch:gauge}, section ``The reading''}
\iamfsentry{502}{A=\frac{H(\text{error now})}{H(\text{error in the same cell when healthy})},}{The reading as an error rate: entropy of the error now over the entropy of the error in the same cell when healthy.}{\derived{} (from $H(\beta)=H(1-\beta)$)}{Chapter~\ref{ch:gauge}, section ``The reading is an error rate''}
\subsection*{Chapter~\ref{ch:meta}: Met-A: the cell's own reference on arrays}
\iamfsentry{503}{\text{Met-A}=\frac{\frac{1}{|\mathcal I|}\sum_{i\in\mathcal I}H(\beta_i)}{H_{\rm ref}},
  \qquad H_{\rm ref}^{\text{EPIC, neutrophil}}=0.330263\ \text{bits}.}{Met-A on arrays, with the frozen EPIC neutrophil reference $0.330263$ bits.}{\calibrated{}}{Eq.~\eqref{eq:meta}; Chapter~\ref{ch:meta}, section ``The reading''}
\iamfsentry{504}{\text{Met-A}_{\rm WB}=\frac{\overline{H(\beta_i)}}{\overline{H(e_i)}},\qquad e_i=\sum_g f_g\,\mu_{g,i},}{Met-A in whole blood: the denominator is the specimen's own composition-matched healthy expectation.}{definition; validated on known mixtures (\measured)}{Eq.~\eqref{eq:metawb}; Chapter~\ref{ch:meta}, section ``Whole blood: the person's own expectation''}
\subsection*{Chapter~\ref{ch:iama}: IAM-A: copy error on single molecules}
\iamfsentry{505}{\varepsilon=\frac{\sum\text{isolated errors}}{\sum\text{opportunities}},
  \qquad E_{\rm hold}=\kB T\ln\frac{1-\varepsilon}{\varepsilon},}{Copy error on single molecules and the holding energy it implies.}{definition; \derived{} for the form, selection rules \calibrated}{Eq.~\eqref{eq:eps}; Chapter~\ref{ch:iama}, section ``The statistic''}
\iamfsentry{506}{\varepsilon_{\min}=\frac{1}{1+e^{M}},\qquad \varepsilon_0=\frac{1}{1+e^{\phi M}},\qquad \Hmin=H(\varepsilon_{\min}),\qquad \Href=H(\varepsilon_0).}{The floor (one ATP per site) and the healthy reference (the measured holding energy), each from a two-state Boltzmann factor.}{\derived{} for the form; floor \calc{} ($2.5\times10^{-8}$ bits); reference \measured{} ($\varepsilon_0=0.032$, $\Href=0.2043$ bits)}{Eq.~\eqref{eq:eps0}; Chapter~\ref{ch:iama}, section ``The floor, the healthy reference and the ceiling''}
\iamfsentry{507}{\text{IAM-A}=\frac{H(\varepsilon)}{P_{\rm cell}\,H(\varepsilon_0)},\qquad P_{\rm neutrophil}=1.149 .}{IAM-A with the neutrophil position $P=1.149$, whole files (pipeline \texttt{loyfer\_pat\_v1}).}{\measured{}}{Eq.~\eqref{eq:iama}; Chapter~\ref{ch:iama}, section ``The healthy position, and the reading''}
\subsection*{Chapter~\ref{ch:cscore}: The C-score: where the departures sit}
\iamfsentry{508}{z_i=\frac{H(\beta_i)-H(\mathrm{ref}_i)}{s_i},}{Residual map: each identity site's entropy departure from the healthy reference, in units of the healthy spread $s_i$.}{definition (\derived{} as a construction)}{Eq.~\eqref{eq:z}; Chapter~\ref{ch:cscore}, section ``The residual map''; see also (\ref{app:fs:514})}
\iamfsentry{509}{c=\frac{\mathrm{var}\left(\sqrt{50}\;\overline{z}_{\rm block}\right)}{\mathrm{var}(z)},\qquad
  C=\frac{c}{c_{\rm healthy}},\qquad c_{\rm healthy}=1.1104,}{C-score: variance of $\sqrt{50}\times$ the 50-site block means over the site variance, divided by the healthy median $1.1104$.}{definition; $c_{\rm healthy}$ \calibrated}{Eq.~\eqref{eq:C}; Chapter~\ref{ch:cscore}, section ``The score''}
\subsection*{Chapter~\ref{ch:temperature}: Temperature, and other species}
\iamfsentry{510}{\varepsilon_0(T)=\frac{1}{1+\exp\!\left(E_{\rm hold}/\kB T\right)} ,}{The copy-error floor as a function of body temperature at a fixed holding energy.}{\derived{} (given a fixed holding energy, a \conjecture)}{Eq.~\eqref{eq:eps0T}; Chapter~\ref{ch:temperature}, section ``The floor is a function of temperature''}
\subsection*{Chapter~\ref{ch:separation}: Component separation: finding the cells}
\iamfsentry{511}{\beta_i=\sum_g f_g\,\mu_{g,i}+\epsilon_i,\qquad f_g\ge0,\ \ \sum_g f_g=1,}{Composition model of Stage~A: measured $\beta$ as a non-negative, sum-to-one mixture of blood-group profiles.}{definition (the mixture model); profiles and markers \calibrated}{Eq.~\eqref{eq:mix}; Chapter~\ref{ch:separation}, section ``The problem''}
\subsection*{Chapter~\ref{ch:atlas}: The reference: purified healthy cells on the specimen's platform}
\iamfsentry{512}{\begin{aligned}\beta^{\mathrm{obs}}_{ics}&\sim\mathcal N\!\left(\mu_{ic}+\delta_{is},\ \sigma^2_{\mathrm{obs},s}+\sigma^2_{\mathrm{loc},i}\right)\quad\text{where measured},\\
  \mu_{ic}&=m_i+e_{ic},\qquad e_{ic}\sim\mathcal N(0,\tau^2),\end{aligned}}{Atlas v2 hierarchical model: observed $\beta$ per locus, cell and source with a source term $\delta_{is}$; cell means drawn around a shared locus level.}{definition (the statistical model); the fit is \measured}{Chapter~\ref{ch:atlas}, section ``The model''}
\subsection*{Chapter~\ref{ch:skytools}: Tools from the sky}
\iamfsentry{513}{z_i=\frac{H(\beta_i)-\overline{H}_i}{s_i},}{Residual of one array's site entropy from the healthy reference mean, in units of the shrunk healthy spread (as Eq.~\eqref{eq:z}).}{definition}{Chapter~\ref{ch:skytools}, section ``Reading a cell's sky''}
\subsection*{Chapter~\ref{ch:sky}: The sky: what the reference cannot explain}
\iamfsentry{514}{z_i=\frac{H(\beta_i)-H(\mathrm{ref}_i)}{s_i},}{Sky value at identity site $i$: entropy offset from the healthy reference in units of the reference spread.}{definition (\derived{} as a construction); $s_i$ \calibrated}{Eq.~\eqref{eq:sky}; Chapter~\ref{ch:sky}, section ``The sky chain v3 draws''; see also (\ref{app:fs:508})}
\iamfsentry{515}{\sigma_{\mathrm{copies}}=\sqrt{\frac{\beta(1-\beta)}{2N}}}{Cell-count floor: binomial spread of a measured fraction from $N$ diploid genome copies (the counterpart of cosmic variance).}{\derived}{Eq.~\eqref{eq:cellcount}; Chapter~\ref{ch:sky}, section ``The cell-count floor''}
\iamfsentry{516}{\mathbf r=\underbrace{P_{\mathrm{span}}\mathbf r}_{\text{composition error}}+\underbrace{(I-P_{\mathrm{span}})\mathbf r}_{\text{what no known mixture can make}},}{The out-of-span map: the residual split into a part any mixture of atlas cells can make and a part none can.}{\derived{} (a projection); its tests are \prediction{} (pre-registered)}{Eq.~\eqref{eq:outspan}; Chapter~\ref{ch:sky}, section ``A cleaner sky: the out-of-span map''}
\section{\texorpdfstring{Part~\ref{part:5}}{Part 7}: Thirty-Seven Orders of Magnitude and the Web that is IAM}
\subsection*{Chapter~\ref{ch:theoryinterp}: One mechanism, two horizons}
\iamfsentry{517}{\frac{S_{\rm info}(R)}{A(R)}\ge\frac{1}{4\ell_P^2}.}{Collapse criterion as holographic saturation: record entropy per area reaching $1/4\ell_P^2$ (met exactly at the Schwarzschild radius).}{\derived{} (equality at $R_s$); its reading as the collapse criterion is \interp}{Chapter~\ref{ch:theoryinterp}, section ``Two regimes: local and global horizons''; see also (\ref{app:fs:162})}
\iamfsentry{518}{M_{\rm eq}=\frac{c^3}{4GH}\approx2.3\times10^{22}M_\odot\quad(H_0=67.16),}{Mass at which a black hole is as cold as the cosmic horizon (as in Chapter~\ref{ch:blackholes}).}{\calc{} (as labeled)}{Chapter~\ref{ch:theoryinterp}, section ``Two regimes: local and global horizons''}
\subsection*{Chapter~\ref{ch:time}: Duration, the arrow of time, and the boundary between potential and actual}
\iamfsentry{519}{\tau=\frac1c\int\sqrt{-g_{\mu\nu}\,dx^\mu dx^\nu}.}{}{\derived{}}{Eq.~\eqref{eq:time_tau}; Chapter~\ref{ch:time}, section ``Three kinds of time''}
\iamfsentry{520}{E(a)=\exp\!\left(1-\frac1a\right),}{}{\derived{}}{Eq.~\eqref{eq:time_Ea}; Chapter~\ref{ch:time}, section ``The second face: the accumulated record''; see also (\ref{app:fs:185}), (\ref{app:fs:206}), (\ref{app:fs:318}), (\ref{app:fs:332}), (\ref{app:fs:408})}
\iamfsentry{521}{E(a)\big|_{\rm photon}\equiv0 .}{}{\derived{}}{Eq.~\eqref{eq:time_Ephoton}; Chapter~\ref{ch:time}, section ``The second face: the accumulated record''}
\iamfsentry{522}{\beta_m=\frac{\Omega_m}{2}=0.15765}{}{\derived{}}{Eq.~\eqref{eq:time_betam}; Chapter~\ref{ch:time}, section ``What the sector split implies''; see also (\ref{app:fs:178}), (\ref{app:fs:319}), (\ref{app:fs:331}), (\ref{app:fs:399})}
\iamfsentry{523}{\ddot\delta_m+2H\dot\delta_m-\frac32\,\Omega_mH_0^2a^{-3}\,\mu(a)\,\delta_m=0,\qquad
\mu(a)=\frac{H^2_{\Lambda\rm CDM}(a)}{H^2_{\Lambda\rm CDM}(a)+\beta_mE(a)H_0^2}<1,}{}{\derived{}}{Eq.~\eqref{eq:time_growth}; Chapter~\ref{ch:time}, section ``The relational framework and shape dynamics''}
\subsection*{Chapter~\ref{ch:virial_partners}: Dark matter and dark energy as virial partners}
\iamfsentry{524}{\begin{aligned}\beta_{\rm temporal}&=0.076618\ (48.6\,\%):&&\text{kinetic decoherence events}\to\sigma_8,\ f\sigma_8,\ \text{RSD},\\
\beta_{\rm geometric}&=0.078825\ (50.0\,\%):&&\text{structural configuration bits}\to\text{lensing/dynamics ratio},\\
\beta_{\rm radiative}&=0.002207\ (1.4\,\%):&&\text{photon channel: SZ, X-ray and lensing ratios}.\end{aligned}}{}{\calc{}}{Eq.~\eqref{eq:vp_channels}; Chapter~\ref{ch:virial_partners}, section ``The three channels of $\beta_m$''}
\subsection*{Chapter~\ref{ch:virial_decoherence}: Gravitational decoherence, the virial partition and the persistence of classical structure}
\iamfsentry{525}{\frac{S_{\rm info}(R)}{A(R)}\ge\frac{k_B}{4\ell_P^2}\qquad(\text{one nat per }4\ell_P^2),}{}{\conjecture{}}{Eq.~\eqref{eq:vd_sat}; Chapter~\ref{ch:virial_decoherence}, section ``The black hole as encoding vault''}
\iamfsentry{526}{\Gamma_{\rm BH}=\frac{P_H}{k_BT_{\rm BH}\ln2}=\frac{c^3}{1920\,GM\ln2}\ \text{bits\,s}^{-1},\qquad P_H=\frac{\hbar c^6}{15360\pi G^2M^2},}{}{\derived{} (sympy, \texttt{verify\_virial\_papers.py} section E)}{Eq.~\eqref{eq:vd_gamma}; Chapter~\ref{ch:virial_decoherence}, section ``The black hole as encoding vault''}
\subsection*{Chapter~\ref{ch:onegauge}: One gauge: the reading over its floor}
\iamfsentry{527}{H(\beta)=H(\varepsilon),\qquad H(x)=-x\log_2x-(1-x)\log_2(1-x).}{At an identity site the entropy of $\beta$ is the entropy of the error; the binary entropy $H(x)$ in bits is that of Chapter~\ref{ch:surface}.}{\derived{} ($H(\beta)=H(1-\beta)$); $H(x)$ is a definition}{Chapter~\ref{ch:onegauge}, section ``A is an error measurement''}
\iamfsentry{528}{\text{Met-A}=\frac{\langle H(\beta)\rangle_{\rm identity\ sites}}{\langle H(\beta)\rangle_{\rm healthy\ reference}}.}{Met-A: mean binary entropy over the cell's identity sites over the same quantity on purified healthy cells (same type, same platform).}{definition; the reference is \calibrated{} (Stage~1)}{Chapter~\ref{ch:onegauge}, section ``A is an error measurement''}
\iamfsentry{529}{\text{IAM-A}=\frac{H(\varepsilon)}{P_{\rm cell}\,H(\varepsilon_0)},\qquad \varepsilon_0=\frac{1}{1+e^{\varphi M}}=0.032,}{IAM-A: copy-error entropy over the physics floor at the cell's position $P_{\rm cell}$; $\varepsilon_0=1/(1+e^{\varphi M})=0.032$.}{definition; the Boltzmann form of $\varepsilon_0$ is \derived, its height \measured; $P$ \measured}{Chapter~\ref{ch:onegauge}, section ``A is an error measurement''}
\subsection*{Chapter~\ref{ch:predictions}: Falsifiable predictions}
\iamfsentry{530}{x_{qp}=\frac{2N\tau_{qp}}{\tau_{\rm TLS}\,n_{cp}V}.}{}{\derived{}}{Chapter~\ref{ch:predictions}, section ``Quasiparticles''}
\subsection*{Chapter~\ref{ch:exploratory}: Exploratory: what the law permits for gravitational engineering}
\iamfsentry{531}{\Delta E_{\rm rec}=\kB T_H\ln2\ \text{per bit},\qquad T_H=\frac{\hbar H}{2\pi\kB},}{}{\calc{}}{Eq.~\eqref{eq:ge_landauer}; Chapter~\ref{ch:exploratory}, section ``Premise: what IAM already says about gravity''}
\iamfsentry{532}{E(a)=\exp\!\left(1-\frac1a\right),\qquad E(1)=1,\qquad E(a\to\infty)=e,}{}{\derived{}}{Eq.~\eqref{eq:ge_Ea}; Chapter~\ref{ch:exploratory}, section ``Premise: what IAM already says about gravity''}
\iamfsentry{533}{H^2_{\rm eff,m}(a)=H^2_{\Lambda{\rm CDM}}(a)+\beta_m\,E(a)\,H_0^2 .}{}{\derived{}}{Eq.~\eqref{eq:ge_Hm}; Chapter~\ref{ch:exploratory}, section ``Premise: what IAM already says about gravity''}
\iamfsentry{534}{\mathbf g_{\rm IAM}(\mathbf x)=-\nabla\Phi_{\rm IAM}(\mathbf x).}{}{\derived{}}{Eq.~\eqref{eq:ge_g}; Chapter~\ref{ch:exploratory}, section ``Zone I: established physics and IAM consequences''}
\iamfsentry{535}{\ddot{\mathbf x}=-\nabla\Phi_{\rm IAM}(\mathbf x).}{}{\derived{}}{Eq.~\eqref{eq:ge_eom}; Chapter~\ref{ch:exploratory}, section ``Zone I: established physics and IAM consequences''}
\iamfsentry{536}{\mathbf f_{\rm felt}=m\left(\ddot{\mathbf x}+\nabla\Phi_{\rm IAM}\right)=0,}{}{\derived{}}{Eq.~\eqref{eq:ge_felt}; Chapter~\ref{ch:exploratory}, section ``Zone I: established physics and IAM consequences''}
\iamfsentry{537}{f_{\rm tidal}\sim mL\,\partial^2\Phi_{\rm IAM},}{}{\derived{}}{Eq.~\eqref{eq:ge_tidal}; Chapter~\ref{ch:exploratory}, section ``Zone I: established physics and IAM consequences''}
\iamfsentry{538}{\nabla\Phi_{\rm IAM}=\mathbf g_{\rm amb}=-\nabla\Phi_{\rm amb},\qquad -\nabla\left(\Phi_{\rm amb}+\Phi_{\rm IAM}\right)=0 ,}{}{\derived{}}{Eq.~\eqref{eq:ge_hover}; Chapter~\ref{ch:exploratory}, section ``Zone I: established physics and IAM consequences''}
\iamfsentry{539}{\theta=\arctan\!\left(\frac{|\nabla\Phi_{\rm IAM}|_{\rm horizontal}}{|\nabla\Phi_{\rm IAM}|_{\rm vertical}}\right),}{}{\derived{}}{Eq.~\eqref{eq:ge_plumb}; Chapter~\ref{ch:exploratory}, section ``Zone I: established physics and IAM consequences''}
\iamfsentry{540}{\tau=-\frac{3GM}{2r^3}\left(I_\perp-I_\parallel\right)\sin2\theta ,}{}{\derived{}}{Eq.~\eqref{eq:ge_ggtorque}; Chapter~\ref{ch:exploratory}, section ``Zone I: established physics and IAM consequences''}
\iamfsentry{541}{v_{\rm rec}=H_0D ,}{}{\calc{}}{Eq.~\eqref{eq:ge_vrec}; Chapter~\ref{ch:exploratory}, section ``Zone I: established physics and IAM consequences''}
\iamfsentry{542}{D_H=\frac{c}{H_0}=1.46\times10^{10}\ \text{ly}\ (H_0=67.16),\qquad 1.35\times10^{10}\ \text{ly}\ (H_0=72.26),}{}{\calc{}}{Eq.~\eqref{eq:ge_DH}; Chapter~\ref{ch:exploratory}, section ``Zone I: established physics and IAM consequences''}
\iamfsentry{543}{\mathrm ds^2=-c^2\mathrm dt^2+\left[\mathrm dx-v_s\,f(r_s)\,\mathrm dt\right]^2+\mathrm dy^2+\mathrm dz^2,\qquad
r_s=\sqrt{(x-x_s(t))^2+y^2+z^2},}{}{\derived{}}{Eq.~\eqref{eq:ge_alc}; Chapter~\ref{ch:exploratory}, section ``Zone I: established physics and IAM consequences''}
\iamfsentry{544}{\mathrm d\tau_{\rm ship}=\mathrm dt_{\rm bubble},}{}{\derived{}}{Eq.~\eqref{eq:ge_tau}; Chapter~\ref{ch:exploratory}, section ``Zone I: established physics and IAM consequences''}
\iamfsentry{545}{\frac{\Delta\tau_{\rm Earth}-\Delta\tau_{\rm ship}}{\Delta\tau}\sim\frac{\Delta\Phi_{\rm IAM}}{c^2}\ll1}{}{\derived{}}{Eq.~\eqref{eq:ge_dtau}; Chapter~\ref{ch:exploratory}, section ``Zone I: established physics and IAM consequences''}
\iamfsentry{546}{\rho=-\frac1{8\pi}\,\frac{v_s^2\,(y^2+z^2)}{4r_s^2}\left(\frac{\mathrm df}{\mathrm dr_s}\right)^2\le0 ,}{}{\derived{}}{Eq.~\eqref{eq:ge_rho}; Chapter~\ref{ch:exploratory}, section ``Zone I: established physics and IAM consequences''}
\iamfsentry{547}{t_{\rm one\text-way}=\frac{D}{v_{\rm eff}}=\frac{444\ \text{ly}}{\xi c}=\frac{444}{\xi}\ \text{years}.}{}{\calc{}}{Eq.~\eqref{eq:ge_tone}; Chapter~\ref{ch:exploratory}, section ``Zone I: established physics and IAM consequences''}
\iamfsentry{548}{\xi=\frac{444\ \text{yr}}{0.0192\ \text{yr}}\approx2.3\times10^4 .}{}{\calc{}}{Eq.~\eqref{eq:ge_xi}; Chapter~\ref{ch:exploratory}, section ``Zone I: established physics and IAM consequences''}
\iamfsentry{549}{\begin{aligned}\Phi_{\rm IAM}(\mathbf x,t)&=\Phi_0(\mathbf x)+\int G(\mathbf x-\mathbf x',t-t')\,J(\mathbf x',t')\,\mathrm d^3x'\,\mathrm dt'
\\&\quad\text{(conjectured form; derivation open)},\end{aligned}}{}{\conjecture{}}{Eq.~\eqref{eq:ge_kernel}; Chapter~\ref{ch:exploratory}, section ``Zone II: the conjectural engineering coupling''}
\iamfsentry{550}{\Phi_{\rm drive}(\mathbf x)=\sum_i A_i\,G(\mathbf x-\mathbf r_i)\,e^{i\phi_i}\qquad\text{(conditional on the unknown }G).}{}{\derived{}}{Eq.~\eqref{eq:ge_drive}; Chapter~\ref{ch:exploratory}, section ``Zone II: the conjectural engineering coupling''}
\iamfsentry{551}{\mathbf F_A+\mathbf F_B=-\iint\rho_A(\mathbf x)\,\rho_B(\mathbf x')\left[\nabla G(\mathbf x-\mathbf x')+\nabla G(\mathbf x'-\mathbf x)\right]\mathrm d^3x\,\mathrm d^3x' ,}{}{\derived{}}{Eq.~\eqref{eq:pr_noforce}; Chapter~\ref{ch:exploratory}, section ``Propulsion and the conservation laws''}
\iamfsentry{552}{F\le\frac{P}{c},\qquad P\ge m\,a\,c .}{}{\derived{}}{Eq.~\eqref{eq:pr_thrust}; Chapter~\ref{ch:exploratory}, section ``Propulsion and the conservation laws''}
\iamfsentry{553}{f_{\rm felt}=m\,\eta\,a ,\qquad |\eta|\le\frac{\varepsilon\,g}{a}\ \text{to keep the felt load below }\varepsilon\,g .}{}{\derived{}}{Eq.~\eqref{eq:pr_eta}; Chapter~\ref{ch:exploratory}, section ``Propulsion and the conservation laws''}
\iamfsentry{554}{\Delta a=\frac{2GML}{r^3}=\frac{2aL}{r} .}{}{\derived{}}{Eq.~\eqref{eq:pr_focus}; Chapter~\ref{ch:exploratory}, section ``Propulsion and the conservation laws''}
\iamfsentry{555}{E=-\frac{v_s^2}{12}\int_0^\infty r^2\left(\frac{\mathrm df}{\mathrm dr}\right)^2\mathrm dr ,}{}{\derived{}}{Eq.~\eqref{eq:pr_Etot}; Chapter~\ref{ch:exploratory}, section ``Propulsion and the conservation laws''}
\iamfsentry{556}{E=-\frac{v_s^2}{12}\left(\frac{R^2}{\Delta}+\frac{\Delta}{12}\right) ,}{}{\derived{}}{Eq.~\eqref{eq:pr_Ewall}; Chapter~\ref{ch:exploratory}, section ``Propulsion and the conservation laws''}
\iamfsentry{557}{F_{\rm net}=\frac{G\,m_{p1}m_{p2}}{r^2}\left(\frac{m_{a2}}{m_{p2}}-\frac{m_{a1}}{m_{p1}}\right) .}{}{\derived{}}{Eq.~\eqref{eq:pr_selfforce}; Chapter~\ref{ch:exploratory}, section ``Propulsion and the conservation laws''}
\iamfsentry{558}{(1-\delta)\,m_d+m_{\rm pay}=0,\qquad \delta=1+\frac{m_{\rm pay}}{m_d} .}{}{\derived{}}{Eq.~\eqref{eq:pr_hover}; Chapter~\ref{ch:exploratory}, section ``Propulsion and the conservation laws''}
